\pdfoutput=1
\documentclass[11pt]{amsart}

\usepackage[margin=1.1in]{geometry}
\usepackage{amsmath,amssymb,amsthm,mathtools}
\usepackage{enumitem}
\usepackage{array}
\usepackage{aliascnt}
\usepackage{microtype}
\usepackage{hyperref}
\usepackage[nameinlink,capitalize]{cleveref}

\hypersetup{colorlinks=true,linkcolor=blue,citecolor=blue,urlcolor=blue}
\raggedbottom

\theoremstyle{plain}
\newtheorem{theorem}{Theorem}[section]

\newaliascnt{lemma}{theorem}
\newtheorem{lemma}[lemma]{Lemma}
\aliascntresetthe{lemma}

\newaliascnt{proposition}{theorem}
\newtheorem{proposition}[proposition]{Proposition}
\aliascntresetthe{proposition}

\newaliascnt{corollary}{theorem}
\newtheorem{corollary}[corollary]{Corollary}
\aliascntresetthe{corollary}

\newaliascnt{hypothesis}{theorem}
\newtheorem{hypothesis}[hypothesis]{Hypothesis}
\aliascntresetthe{hypothesis}

\theoremstyle{definition}
\newaliascnt{definition}{theorem}
\newtheorem{definition}[definition]{Definition}
\aliascntresetthe{definition}

\theoremstyle{remark}
\newaliascnt{remark}{theorem}
\newtheorem{remark}[remark]{Remark}
\aliascntresetthe{remark}

\crefname{theorem}{Theorem}{Theorems}
\Crefname{theorem}{Theorem}{Theorems}
\crefname{lemma}{Lemma}{Lemmas}
\Crefname{lemma}{Lemma}{Lemmas}
\crefname{proposition}{Proposition}{Propositions}
\Crefname{proposition}{Proposition}{Propositions}
\crefname{corollary}{Corollary}{Corollaries}
\Crefname{corollary}{Corollary}{Corollaries}
\crefname{hypothesis}{Hypothesis}{Hypotheses}
\Crefname{hypothesis}{Hypothesis}{Hypotheses}
\crefname{definition}{Definition}{Definitions}
\Crefname{definition}{Definition}{Definitions}
\crefname{remark}{Remark}{Remarks}
\Crefname{remark}{Remark}{Remarks}
\crefalias{theorem}{theorem}
\crefalias{lemma}{lemma}
\crefalias{proposition}{proposition}
\crefalias{corollary}{corollary}
\crefalias{hypothesis}{hypothesis}
\crefalias{definition}{definition}
\crefalias{remark}{remark}

\numberwithin{equation}{section}

\newcommand{\R}{\mathbb R}
\newcommand{\N}{\mathbb N}
\newcommand{\eps}{\varepsilon}
\newcommand{\loc}{\mathrm{loc}}
\newcommand{\supp}{\operatorname{supp}}
\newcommand{\divergence}{\operatorname{div}}
\newcommand{\dist}{\operatorname{dist}}
\newcommand{\esssup}{\operatorname*{ess\,sup}}
\newcommand{\calS}{\mathcal S}
\newcommand{\calR}{\mathcal R}
\newcommand{\calC}{\mathcal C}
\newcommand{\calU}{\mathcal U}
\newcommand{\calH}{\mathcal H}
\DeclareMathOperator{\diver}{div}

\title[Conditional Local Type I Navier--Stokes Reduction]
{Conditional Local Type I Navier--Stokes Reduction: Concentration, Seregin Limits, and Liouville-Class Closures}

\author{Guillem Duran Ballester}
\email{guillem@fragile.tech}
\date{}

\subjclass[2020]{35Q30, 35B44, 35B40, 76D05}
\keywords{Navier--Stokes equations, Type I singularities, ancient solutions, Seregin limits, local concentration, Liouville theorems}

\begin{document}

\begin{abstract}
We prove the concentration, compactness, and Liouville reductions needed for a
conditional exclusion of local pointwise Type I Navier--Stokes singularities.
Starting from a finite-energy local pointwise Type I singular
point, we obtain a raw extracted bounded centered Seregin ancient profile
with positive compact local \(L^3\) velocity mass, then pass to the raw
generated Seregin state space with invariant local nonvanishing.  We define the
small-amplitude, stationary \(L^3\), uniformly \(L^3\)-tight, and
structure and decay classes, and prove their exclusion by
perturbative, stationary, endpoint \(L^3\), and axisymmetric or
weighted-vorticity Liouville theorems.  The exclusion of the uniformly
\(L^3\)-tight class
rests on the Albritton--Barker endpoint ancient Liouville theorem after a
physical pullback on truncated ancient intervals.  The only conditional
ingredient is the residual closure hypothesis \Cref{p1:hyp:no-remainder}.
The final Type I exclusion is conditional on residual closure for the
raw generated Seregin state space
\(\calS_{\mathrm{raw\text{-}gen}}(u,p,z_*;R_0,\eta_*)\) defined below.
\end{abstract}

\maketitle
\tableofcontents

\section*{Introduction and Structure of the Argument}

The conditional local Type I reduction reduces a local
pointwise Type I singularity to a raw extracted bounded centered ancient
solution, then to a raw generated Seregin state with invariant local
nonvanishing, and excludes the non-residual Liouville classes obtained from
standard local compactness and explicit Liouville theorems.  The residual
class is handled only by the residual-class hypothesis
\Cref{p1:hyp:no-remainder}; no unconditional residual exclusion is claimed in
this paper.

The generated state space used in the final assembly is a raw local Seregin
state space.  The projected mild formulation is not imposed on every raw
generated state.  Instead, mildness is treated as a branch condition in those
Liouville classes whose exclusion theorems require it.  Generated states not
known to satisfy the full hypotheses of one of the explicit Liouville branches
belong to the remaining class \(\calR(\calS)\), and are removed only by the
residual closure hypothesis.

The reduction has four components.  First, the local concentration theorem
shows that a singular point carries positive Caffarelli--Kohn--Nirenberg
critical mass at arbitrarily small scales and gives the local Type I/local Type
II dichotomy.  It also identifies the endpoint weak-Serrin estimate used to
obtain terminal velocity no-escape for local pointwise Type I rescalings.
Second, the Type I blow-up analysis gives the admissible Seregin extraction:
an admissible Type I extraction produces a raw smooth bounded ancient solution
of the centered self-similar equation with positive compact local \(L^3\)
velocity mass.  The compact mass is converted into invariant local
nonvanishing before forming the raw generated descendant hull.  The paper then
defines the Liouville classes and the residual class on that raw hull.

Third, the uniformly tight analysis excludes nonzero uniformly \(L^3\)-tight
generated raw Seregin states that lie in the mild stratum.  Boundedness plus uniform
\(L^3\)-tightness gives a backward sequence of finite critical norm after
physical pullback, so the Albritton--Barker endpoint ancient Liouville theorem
applies.  Fourth, the structured part gives the axisymmetric and
perturbative weighted-vorticity exclusions used in the Liouville-class analysis.
These results turn the corresponding structured hypotheses into the zero
ancient solution, contradicting the invariant local nonvanishing retained by
the raw generated Seregin state space.

Consequently, a generated raw Seregin state in
\(\calS_{\mathrm{raw\text{-}gen}}(u,p,z_*;R_0,\eta_*)\) must either belong to one of the excluded
Liouville classes or to the residual class.  The Type I contradiction is
therefore conditional on the residual-class hypothesis
\Cref{p1:hyp:no-remainder}.
No unconditional Type I regularity theorem is claimed in this paper; the final
regularity conclusion is conditional on the residual closure hypothesis for the
raw generated state space.

\subsection*{Dependency table}

\begin{center}
\small
\setlength{\tabcolsep}{4pt}
\begin{tabular}{@{}>{\raggedright\arraybackslash}p{0.23\textwidth}
                >{\raggedright\arraybackslash}p{0.38\textwidth}
                >{\raggedright\arraybackslash}p{0.27\textwidth}@{}}
\hline
Step & Status/source & Role \\
\hline
Velocity epsilon criterion & imported from Gustafson--Kang--Tsai & singularity gives \(L^3\) concentration \\
Weak-Serrin Type I bridge & imported from Albritton--Barker, Lemma~2.5; used here & local Type I gives \(A,C,D,E\) bounds \\
Terminal no-escape & proved here using Albritton--Barker & prevents mass from escaping into \(t=T\) \\
Seregin extraction & proved here & produces a raw extracted centered ancient profile \\
Raw generated closure & defined here & gives the ambient state space for residual closure \\
Mild stratum & defined here & supplies the branch condition for mild Liouville theorems \\
Tight \(L^3\) branch & proved here using Albritton--Barker, Theorem~1.2 & excludes a non-residual class \\
Structured Liouville branches & imported theorem by theorem plus Type I constant removal & exclude the structured classes \\
Residual closure & conditional input from the companion residual paper & removes \(\calR(\calS_{\mathrm{raw\text{-}gen}}(u,p,z_*;R_0,\eta_*))\) \\
Contradiction & conditional on residual closure & rules out local pointwise Type I singularities \\
\hline
\end{tabular}
\end{center}

\subsection*{How to read this paper}

Readers interested only in the construction of the ancient limit should read
Parts I and II together with the final assembly
\Cref{p1:thm:final-assembly}.  Readers interested in the non-residual
Liouville exclusions should read Parts III and IV.  The residual closure
hypothesis \Cref{p1:hyp:no-remainder} is stated as a conditional analytic
assumption.

\subsection*{Pressure notation}

Lowercase \(p\) denotes the source pressure of a suitable weak solution.
\(Q\) denotes a physical ancient pressure obtained as a local distributional
limit after blow-up.  \(\Pi\) denotes a centered pressure obtained from
\(Q\) by the self-similar change of variables.  All pressures are
considered modulo functions of time only; this gauge ambiguity is invisible
to the equations and to the local energy inequality
(\Cref{p1:lem:gauges,p1:rem:centered-gauge,p1:lem:pressure-atlas,p1:rem:pressure-atlas}).

\clearpage
\section*{Part I. Local Concentration and the Type I/Type II Dichotomy}

\subsection*{Overview}

This part proves a local analytic reduction for possible finite-time
singularities of the three-dimensional Navier--Stokes equations.  For a
suitable weak solution,
every singular point at time \(T\) has arbitrarily small backward parabolic
cylinders on which both the standard mean-subtracted
Caffarelli--Kohn--Nirenberg scale-invariant quantity \(C+D\) and the
scale-invariant velocity quantity \(C\) are bounded below by positive
constants.  If these local scale-invariant quantities tend to zero at each
point of the singular set, the solution is regular at time \(T\) by the
\(\varepsilon\)-regularity theorem.  Along any such positive concentration
sequence, the solution either satisfies a local Type I scale-invariant bound or
belongs to the local Type II alternative.  For local pointwise Type I terminal
singularities of suitable weak solutions, the endpoint weak-Serrin argument gives
the velocity no-escape condition
used in the Type I blow-up analysis.

\section{Introduction}

Let \((u,p)\) be a suitable weak solution of the three-dimensional
incompressible Navier--Stokes equations
\[
  \partial_t u+u\cdot\nabla u+\nabla p=\Delta u,
  \qquad \nabla\cdot u=0.
\]
The viscosity is normalized to one.  The weak
solution theory goes back to Leray \cite{Leray1934}.  Suitable weak solutions
and their partial regularity theory were introduced and developed by Scheffer
\cite{Scheffer1976} and by Caffarelli, Kohn, and Nirenberg \cite{CKN1982}; see
also Lin \cite{Lin1998}.  The local reduction below connects a possible
singularity at time \(T\) with the Type I/Type II blow-up rate dichotomy.

Scale-invariant regularity and blow-up criteria provide the standard background
for this formulation; see Serrin \cite{Serrin1962}, Giga \cite{Giga1986},
Escauriaza--Seregin--{\v S}verak \cite{ESS2003}, and Seregin
\cite{Seregin2012}.  The principal regularity theorem is the
Caffarelli--Kohn--Nirenberg \(\varepsilon\)-regularity theorem.

The proof has four components.  First, the local scale-invariant
quantities \(C\) and \(D\) are finite for suitable weak solutions and invariant
under the usual parabolic scaling.  Second, the Caffarelli--Kohn--Nirenberg
\(\varepsilon\)-regularity theorem implies that vanishing of \(C+D\) at a
putative singular point gives local regularity, while the velocity
\(\varepsilon\)-regularity criterion gives positive concentration in the
scale-invariant \(L^3_{x,t}\) quantity.  Third, a finite-time singularity
therefore produces sequences of shrinking cylinders on which \(C+D\), and
separately \(C\), are bounded below by positive constants.  Finally, once such
a sequence is fixed, it either satisfies a Type I scale-invariant bound or
belongs to the local Type II alternative.

The concentration and Type I/Type II dichotomy proofs are local in
spacetime.  Apart from the standard Caffarelli--Kohn--Nirenberg
\(\varepsilon\)-regularity theorem, all quantities and alternatives used in
that local analysis are defined in this part.  The local Type I reduction section
shows how the endpoint weak-Serrin estimate upgrades the local pointwise Type I
envelope to terminal velocity no-escape through the local endpoint weak-Serrin
Type I estimate.

\section{Suitable Weak Solutions and Singular Points}

\begin{definition}[Suitable weak solution]
Let \(I\subset\mathbb R\) be an open interval.  A pair \((u,p)\) is a suitable
weak solution on \(\mathbb R^3\times I\) if
\[
  u\in L^\infty_tL^2_{x,\mathrm{loc}}
  \cap L^2_tH^1_{x,\mathrm{loc}},
  \qquad
  p\in L^{3/2}_{\mathrm{loc}},
  \qquad
  \nabla\cdot u=0,
\]
the Navier--Stokes equations hold in the sense of distributions, and the local
energy inequality holds: for almost every \(t_1<t_2\) in \(I\) and every
nonnegative \(\phi\in C_c^\infty(\mathbb R^3\times I)\),
\[
\begin{aligned}
&\int_{\mathbb R^3}\frac{|u(x,t_2)|^2}{2}\phi(x,t_2)\,dx
 +\int_{t_1}^{t_2}\int_{\mathbb R^3}|\nabla u|^2\phi\,dx\,dt \\
&\le
 \int_{\mathbb R^3}\frac{|u(x,t_1)|^2}{2}\phi(x,t_1)\,dx
 +\int_{t_1}^{t_2}\int_{\mathbb R^3}
   \frac{|u|^2}{2}(\partial_t\phi+\Delta\phi)\,dx\,dt \\
&\quad
 +\int_{t_1}^{t_2}\int_{\mathbb R^3}
   \left(\frac{|u|^2}{2}+p\right)u\cdot\nabla\phi\,dx\,dt .
\end{aligned}
\]
The pressure is understood modulo functions of time.
\end{definition}

For \(z_0=(x_0,t_0)\) and \(r>0\), write
\[
  Q_r(z_0)=B_r(x_0)\times(t_0-r^2,t_0).
\]
Also write \(Q_r(x_0,t_0)\) for \(Q_r((x_0,t_0))\).

\begin{definition}[Singular set at a fixed time]
For a time \(T\), define
\[
  \Sigma(T)=
  \{x\in\mathbb R^3:\ u\text{ is not locally bounded in any }
  Q_r(x,T)\}.
\]
Equivalently, \(x\notin\Sigma(T)\) if there are \(r>0\) and \(M<\infty\) such
that \(\|u\|_{L^\infty(Q_r(x,T))}\le M\).
\end{definition}

If a finite-time singularity occurs at time \(T\), then
\(\Sigma(T)\ne\emptyset\) in this local sense.

\subsection*{Temporal cutoff convention}

\begin{lemma}[Temporal cutoff convention]\label{p0:lem:temporal-cutoff}
For a suitable weak solution, every local energy inequality used below with
endpoint times or compact time slabs is obtained by applying the standard
distributional local energy inequality to a test function of the form
\(\phi(x,t)\chi_\varepsilon(t)\), where \(\chi_\varepsilon\) is a smooth
nonincreasing temporal cutoff with \(\chi_\varepsilon=1\) on a closed
sub-interval and supported in a slightly larger open interval, and then
letting \(\varepsilon\downarrow0\) along an admissible set of times where the
\(L^2\) trace exists.
\end{lemma}

\noindent We invoke this convention silently in subsequent proofs.

\section{Parabolic Rescaling and Scale-Invariant Quantities}

Fix \(z_0=(x_0,T)\).  For \(r>0\), define the parabolic rescaling
\[
  u^{(r)}(y,s)=r\,u(x_0+ry,T+r^2s),
  \qquad
  p^{(r)}(y,s)=r^2p(x_0+ry,T+r^2s).
\]
The rescaled variables are defined on the set of \((y,s)\) for which
\((x_0+ry,T+r^2s)\) belongs to the original spacetime domain.

\begin{proposition}[Invariance under parabolic rescaling]
If \((u,p)\) is a suitable weak solution in a cylinder containing
\((x_0,T)\), then
\((u^{(r)},p^{(r)})\) is a suitable weak solution on the rescaled domain.  The
local energy inequality is preserved, and the pressure remains defined modulo
functions of the rescaled time.
\end{proposition}

\begin{proof}
The distributional equations and the divergence condition follow from the
change of variables \(x=x_0+ry\), \(t=T+r^2s\).  The local energy inequality
is verified to fix the normalization.

Let \(\psi\in C_c^\infty\) be nonnegative and compactly supported in the
rescaled domain.  Define
\[
  \phi(x,t)=
  \psi\left(\frac{x-x_0}{r},\frac{t-T}{r^2}\right).
\]
If \(y=(x-x_0)/r\) and \(s=(t-T)/r^2\), then
\[
  \partial_t\phi=r^{-2}(\partial_s\psi)(y,s),\qquad
  \nabla_x\phi=r^{-1}(\nabla_y\psi)(y,s),\qquad
  \Delta_x\phi=r^{-2}(\Delta_y\psi)(y,s).
\]
For almost every \(s_1<s_2\), the corresponding times
\(T+r^2s_1<T+r^2s_2\) are admissible in the original local energy inequality.
Apply that inequality for \((u,p)\) to \(\phi\) on
\([T+r^2s_1,T+r^2s_2]\).  After the change of variables
\(x=x_0+ry\), \(t=T+r^2s\), using
\[
  u(x_0+ry,T+r^2s)=r^{-1}u^{(r)}(y,s),\qquad
  p(x_0+ry,T+r^2s)=r^{-2}p^{(r)}(y,s),
\]
each term contains the same common factor \(r\).  Dividing by this factor gives
the rescaled inequality
\[
\begin{aligned}
&\int_{\mathbb R^3}\frac{|u^{(r)}(y,s_2)|^2}{2}\psi(y,s_2)\,dy
 +\int_{s_1}^{s_2}\int_{\mathbb R^3}|\nabla_y u^{(r)}|^2\psi\,dy\,ds \\
&\le
 \int_{\mathbb R^3}\frac{|u^{(r)}(y,s_1)|^2}{2}\psi(y,s_1)\,dy
 +\int_{s_1}^{s_2}\int_{\mathbb R^3}
   \frac{|u^{(r)}|^2}{2}(\partial_s\psi+\Delta_y\psi)\,dy\,ds \\
&\quad
 +\int_{s_1}^{s_2}\int_{\mathbb R^3}
   \left(\frac{|u^{(r)}|^2}{2}+p^{(r)}\right)
   u^{(r)}\cdot\nabla_y\psi\,dy\,ds .
\end{aligned}
\]
Thus \((u^{(r)},p^{(r)})\) is suitable on the rescaled domain.  The pressure
transforms according to the Navier--Stokes scaling, and addition of a function
of \(t\) becomes addition of a function of \(s\).
\end{proof}

\begin{definition}[Scale-invariant local quantities]
For \(z_0=(x_0,t_0)\) and \(r>0\), set
\[
  C(z_0,r)=r^{-2}\int_{Q_r(z_0)} |u|^3\,dx\,dt
\]
and
\[
  D(z_0,r)=r^{-2}
  \int_{t_0-r^2}^{t_0}\int_{B_r(x_0)}
  |p(x,t)-(p)_{B_r(x_0)}(t)|^{3/2}\,dx\,dt.
\]
Here \((p)_{B_r(x_0)}(t)\) denotes the spatial mean of \(p(\cdot,t)\) on
\(B_r(x_0)\).
\end{definition}

The subtraction of the spatial mean removes the pressure ambiguity in the form
used by the Caffarelli--Kohn--Nirenberg criterion.

\begin{proposition}[Local finiteness and scaling]\label{p0:prop:critical-scaling}
Let \((u,p)\) be a suitable weak solution in a cylinder containing
\(z_0=(x_0,T)\).  Then
\(C(z_0,r)\) and \(D(z_0,r)\) are finite for every sufficiently small \(r\).
Moreover, for every fixed \(R<\infty\),
\[
\begin{aligned}
&\int_{-R^2}^{0}\int_{B_R}
\left(
|u^{(r)}|^3+
|p^{(r)}-(p^{(r)})_{B_R}(s)|^{3/2}
\right)\,dy\,ds \\
&\qquad
=R^2\bigl(C(z_0,rR)+D(z_0,rR)\bigr).
\end{aligned}
\]
\end{proposition}

\begin{proof}
Local finiteness of the velocity term follows from
\[
  u\in L^\infty_tL^2_{x,\mathrm{loc}}
  \cap L^2_tH^1_{x,\mathrm{loc}}
  \subset L^3_{\mathrm{loc}}
\]
on finite cylinders.  Indeed, if \(B\Subset\mathbb R^3\) and
\((a,b)\Subset I\), Sobolev and interpolation give
\[
  \int_a^b \|u(t)\|_{L^3(B)}^3\,dt
  \le
  C
  \left(\operatorname*{ess\,sup}_{a<t<b}\|u(t)\|_{L^2(B)}\right)^{3/2}
  \left(\int_a^b \|u(t)\|_{H^1(B)}^2\,dt\right)^{3/4}
  (b-a)^{1/4}.
\]
Thus \(u\in L^3_{\mathrm{loc}}\).  The pressure term is finite because
\(p\in L^{3/2}_{\mathrm{loc}}\), and subtracting a spatial mean over the same
ball preserves local \(L^{3/2}\) integrability.  More explicitly, for a ball
\(B\),
\[
  |p-(p)_B(t)|^{3/2}\le C\left(|p|^{3/2}+|(p)_B(t)|^{3/2}\right),
\]
while Jensen's inequality gives
\[
  |(p)_B(t)|^{3/2}
  \le \frac1{|B|}\int_B |p(x,t)|^{3/2}\,dx.
\]
After integration over \(B\) and over time, the mean-subtracted pressure term is
controlled by the local \(L^{3/2}\) norm of \(p\).

For the scaling identity, use \(x=x_0+ry\), \(t=T+r^2s\).  The velocity term
transforms as
\[
  \int_{-R^2}^{0}\int_{B_R}|u^{(r)}|^3\,dy\,ds
  =r^{-2}\int_{T-(rR)^2}^{T}\int_{B_{rR}(x_0)}|u|^3\,dx\,dt.
\]
The pressure means transform according to
\[
  (p^{(r)})_{B_R}(s)=r^2(p)_{B_{rR}(x_0)}(T+r^2s),
\]
because
\[
  \frac1{|B_R|}\int_{B_R} r^2p(x_0+ry,T+r^2s)\,dy
  =
  \frac{r^2}{|B_{rR}|}\int_{B_{rR}(x_0)}p(x,T+r^2s)\,dx .
\]
Therefore
\[
\begin{aligned}
&\int_{-R^2}^{0}\int_{B_R}
|p^{(r)}-(p^{(r)})_{B_R}(s)|^{3/2}\,dy\,ds \\
&\qquad =
r^{-2}\int_{T-(rR)^2}^{T}\int_{B_{rR}(x_0)}
|p-(p)_{B_{rR}(x_0)}(t)|^{3/2}\,dx\,dt .
\end{aligned}
\]
Adding the velocity and pressure identities gives the displayed formula.
\end{proof}

\section{\texorpdfstring{\(\varepsilon\)-Regularity}{Epsilon-Regularity} and the Vanishing Case}

\begin{theorem}[Caffarelli--Kohn--Nirenberg \(\varepsilon\)-regularity]
\label{p0:thm:ckn}
There exists a universal constant \(\varepsilon_{\rm CKN}>0\) with the
following property.  Let \((u,p)\) be a suitable weak solution in \(Q_r(z_0)\).
If
\[
  C(z_0,r)+D(z_0,r)<\varepsilon_{\rm CKN},
\]
then \(u\) is locally bounded in a smaller cylinder, for instance in
\(Q_{r/2}(z_0)\).  In particular, if \(z_0=(x_0,T)\), then
\(x_0\notin\Sigma(T)\).
\end{theorem}

\begin{proof}
We use the local \(\varepsilon\)-regularity theorem of
Caffarelli--Kohn--Nirenberg \cite{CKN1982}; see also Lin \cite{Lin1998}.  The
statement above is written with the spatial mean subtraction used in \(D\).
Equivalent representatives of the pressure differ by functions of time, and
the spatial mean subtraction on \(B_r(x_0)\) removes this ambiguity.
\end{proof}

\begin{theorem}[Velocity \(\varepsilon\)-regularity]\label{p0:thm:velocity-ckn}
There exists a universal constant \(\varepsilon_v>0\) with the following
property.  Let \((u,p)\) be a suitable weak solution in \(Q_r(z_0)\).  If
\[
  C(z_0,r)<\varepsilon_v,
\]
then \(u\) is locally bounded in \(Q_{r/2}(z_0)\).  In particular, if
\(z_0=(x_0,T)\), then \(x_0\notin\Sigma(T)\).
\end{theorem}

\begin{proof}
We use the one-scale velocity criterion of
Gustafson--Kang--Tsai \cite{GustafsonKangTsai2007} in the case
\(p=q=3\).  In their notation the scaled \(L^{p,q}_{x,t}\) velocity norm is
small under the range \(3/p+2/q\le2\); taking \(p=q=3\) gives the
scale-invariant quantity \(C(z_0,r)\) used here, up to the harmless
normalization convention.  Thus a singular point must
satisfy the concentration consequence
\[
   \limsup_{r\downarrow0} C(z_0,r)\ge\varepsilon_v .
\]
The statement is scale invariant because \(C\) is invariant under the
parabolic Navier--Stokes scaling.
\end{proof}

\begin{definition}[Vanishing local scale-invariant quantities at time \(T\)]
We say that the local scale-invariant quantities vanish at time \(T\) if, for
every \(x_0\in\Sigma(T)\),
\[
  \limsup_{r\downarrow0}
  \bigl(C((x_0,T),r)+D((x_0,T),r)\bigr)=0.
\]
If \(\Sigma(T)=\emptyset\), there is no condition to impose.
\end{definition}

\begin{proposition}[Equivalent rescaled formulation]
\label{p0:prop:vanishing-rescaled}
The preceding condition is equivalent to the following statement: for every
\(x_0\in\Sigma(T)\) and every fixed \(R<\infty\),
\[
  \lim_{r\downarrow0}
  \int_{-R^2}^{0}\int_{B_R}
  \left(
  |u^{(r)}|^3+
  |p^{(r)}-(p^{(r)})_{B_R}(s)|^{3/2}
  \right)\,dy\,ds=0.
\]
\end{proposition}

\begin{proof}
By \cref{p0:prop:critical-scaling}, the rescaled integral over
\(B_R\times(-R^2,0)\) equals
\[
  R^2\bigl(C((x_0,T),rR)+D((x_0,T),rR)\bigr).
\]
Thus vanishing of \(C+D\) implies vanishing on every fixed rescaled cylinder.
Conversely, taking \(R=1\) gives the original pointwise condition.
\end{proof}

\begin{theorem}[Regularity under vanishing local quantities]\label{p0:thm:no-concentration}
Let \((u,p)\) be a suitable weak solution on
\(\mathbb R^3\times(T-\delta,T)\) for some \(\delta>0\).  If the local
scale-invariant quantities vanish at time \(T\), then
\(\Sigma(T)=\emptyset\).  Thus no finite-time singularity occurs at time \(T\)
in the no-concentration case.
\end{theorem}

\begin{proof}
Suppose, toward a contradiction, that \(x_0\in\Sigma(T)\).  By the vanishing
hypothesis, choose \(r>0\) such that
\[
  C((x_0,T),r)+D((x_0,T),r)<\varepsilon_{\rm CKN}.
\]
\Cref{p0:thm:ckn} implies that \(u\) is bounded in \(Q_{r/2}(x_0,T)\), hence
\(x_0\notin\Sigma(T)\).  This contradicts the choice of \(x_0\).  Hence
\(\Sigma(T)=\emptyset\).
\end{proof}

\begin{corollary}[Vanishing local quantities exclude a terminal singularity]
\label{p0:cor:escape}
Let \((u,p)\) be suitable near \((x_0,T)\).  If
\[
  \limsup_{r\downarrow0}
  \bigl(C((x_0,T),r)+D((x_0,T),r)\bigr)=0,
\]
then \(u\) is bounded in some backward cylinder \(Q_\rho(x_0,T)\), and hence
\(x_0\notin\Sigma(T)\).  Equivalently, membership in \(\Sigma(T)\) requires
that \(C+D\) have a positive lower bound along a sequence of shrinking backward
parabolic cylinders.
\end{corollary}

\begin{proof}
The hypothesis gives a scale at which \(C+D\) is below the universal threshold
in \cref{p0:thm:ckn}.  The \(\varepsilon\)-regularity theorem gives local boundedness in
a backward cylinder ending at \(T\), so \(x_0\notin\Sigma(T)\).
\end{proof}

\section{Positive Local Concentration}

\begin{lemma}[Failure of vanishing gives a positive concentration sequence]
\label{p0:lem:positive-concentration}
Suppose the local scale-invariant quantities do not vanish at time \(T\).  Then
there exist \(x_0\in\Sigma(T)\), a constant \(\eta>0\), and radii
\(r_n\downarrow0\) such that
\[
  C((x_0,T),r_n)+D((x_0,T),r_n)\ge\eta
  \qquad\text{for all }n.
\]
Equivalently, the rescaled sequence around \((x_0,T)\) has a positive lower
bound for the corresponding local scale-invariant quantities on a fixed
bounded cylinder.
\end{lemma}

\begin{proof}
Negating the definition of vanishing gives a point \(x_0\in\Sigma(T)\) with
positive limsup of
\[
  C((x_0,T),r)+D((x_0,T),r)
\]
as \(r\downarrow0\).  Choose \(\eta>0\) below that limsup and choose a sequence
\(r_n\downarrow0\) along which the displayed lower bound holds.  The rescaled
formulation follows from \cref{p0:prop:critical-scaling}.
\end{proof}

\begin{definition}[Positive local concentration sequence]
\label{p0:def:positive-concentration-sequence}
Let \(x_0\in\Sigma(T)\).  A sequence \(r_n\downarrow0\) is a positive local
concentration sequence at \((x_0,T)\) if there is \(\eta>0\) such that
\[
  C((x_0,T),r_n)+D((x_0,T),r_n)\ge\eta
  \qquad\text{for all }n.
\]
The corresponding parabolic rescalings are those defined by
\[
  u_n(y,s)=r_nu(x_0+r_ny,T+r_n^2s),
  \qquad
  p_n(y,s)=r_n^2p(x_0+r_ny,T+r_n^2s).
\]
\end{definition}

\begin{proposition}[Persistence of positive concentration]
\label{p0:prop:compactness-persistence}
Let \((u_n,p_n)\) be suitable weak solutions on a fixed cylinder
\(B_R\times(-R^2,0)\).  Suppose that
\[
  \int_{-R^2}^{0}\int_{B_R}
  \left(|u_n|^3+|p_n-(p_n)_{B_R}(s)|^{3/2}\right)\,dy\,ds
  \ge\eta>0
\]
for all \(n\).  Suppose further that, after passing to a subsequence,
\[
  u_n\to U\quad\text{in }L^3(B_R\times(-R^2,0))
\]
and
\[
  p_n-(p_n)_{B_R}(s)\to \Pi
  \quad\text{in }L^{3/2}(B_R\times(-R^2,0)).
\]
Then
\[
  \int_{-R^2}^{0}\int_{B_R}
  \left(|U|^3+|\Pi|^{3/2}\right)\,dy\,ds\ge\eta.
\]
In particular, the limiting pair retains the same positive lower bound for the
mean-subtracted local scale-invariant quantity.  If
\(\Pi=P-(P)_{B_R}(s)\), this is the corresponding
Caffarelli--Kohn--Nirenberg quantity of \((U,P)\) on
\(B_R\times(-R^2,0)\).
\end{proposition}

\begin{proof}
Strong convergence in \(L^3\) gives
\[
  \int_{-R^2}^{0}\int_{B_R}|u_n|^3\,dy\,ds
  \longrightarrow
  \int_{-R^2}^{0}\int_{B_R}|U|^3\,dy\,ds.
\]
Strong convergence in \(L^{3/2}\) gives the analogous convergence of the
mean-subtracted pressure integrals.  Passing to the limit in the lower bound
yields the asserted inequality.
\end{proof}

\begin{remark}
\Cref{p0:prop:compactness-persistence} is an auxiliary compactness observation,
independent of the proof of the dichotomy below.  The Seregin extraction in
Part II uses only the velocity-only version below; strong
\(L^{3/2}\)-pressure convergence is not needed for nontriviality.
\end{remark}

\begin{proposition}[Persistence of velocity concentration]
\label{p0:prop:velocity-persistence}
If \(u_n\to U\) strongly in \(L^3(B_R\times(-R^2,0))\) and
\[
  \int_{-R^2}^{0}\int_{B_R}|u_n|^3\,dy\,ds\ge\eta,
\]
then
\[
  \int_{-R^2}^{0}\int_{B_R}|U|^3\,dy\,ds\ge\eta.
\]
\end{proposition}

\begin{proof}
Immediate from continuity of \(\|\cdot\|_{L^3}^3\) under strong
\(L^3\)-convergence.
\end{proof}

\begin{proposition}[Local alternative]\label{p0:prop:local-alternative}
Let \((u,p)\) be a suitable weak solution on
\(\mathbb R^3\times(T-\delta,T)\) for some \(\delta>0\).  Exactly one of the
following alternatives holds.
\begin{enumerate}[label=\textup{(\roman*)}]
\item For every \(x_0\in\Sigma(T)\),
\[
  \limsup_{r\downarrow0}
  \bigl(C((x_0,T),r)+D((x_0,T),r)\bigr)=0.
\]
\item There exist \(x_0\in\Sigma(T)\), \(\eta>0\), and \(r_n\downarrow0\) such
that
\[
  C((x_0,T),r_n)+D((x_0,T),r_n)\ge\eta
  \qquad\text{for all }n.
\]
\end{enumerate}
\end{proposition}

\begin{proof}
The first alternative is a universal assertion over the singular set, and the
second is its failure expressed by a sequence.  If the first alternative fails,
then for some \(x_0\in\Sigma(T)\) the corresponding limsup is positive;
choosing a positive number below this limsup and a sequence along which the
lower bound holds gives the second alternative.  Conversely, the second
alternative contradicts the first.  Hence precisely one alternative holds.
\end{proof}

\begin{theorem}[Finite-time singularity yields positive local scale-invariant concentration]
\label{p0:thm:singular-positive}
Let \((u,p)\) be a suitable weak solution on
\(\mathbb R^3\times(T-\delta,T)\).  If \(\Sigma(T)\ne\emptyset\), then there
are \(x_0\in\Sigma(T)\), \(\eta>0\), and \(r_n\downarrow0\) such that
\[
  C((x_0,T),r_n)+D((x_0,T),r_n)\ge\eta
  \qquad\text{for all }n.
\]
\end{theorem}

\begin{proof}
If the first alternative in \cref{p0:prop:local-alternative} held, then
\cref{p0:thm:no-concentration} would imply \(\Sigma(T)=\emptyset\), contrary to
the hypothesis.  Therefore the second alternative in
\cref{p0:prop:local-alternative} holds.
\end{proof}

\begin{corollary}[Singular points carry positive velocity concentration]
\label{p0:cor:velocity-concentration}
Let \((u,p)\) be a suitable weak solution on
\(\mathbb R^3\times(T-\delta,T)\).  If \(x_0\in\Sigma(T)\), then there are
\(\eta_v>0\) and \(r_n\downarrow0\) such that
\[
  C((x_0,T),r_n)\ge\eta_v
  \qquad\text{for all }n.
\]
\end{corollary}

\begin{proof}
If the conclusion failed, then
\[
  \limsup_{r\downarrow0} C((x_0,T),r)=0.
\]
Choose \(r>0\) so small that
\[
  C((x_0,T),r)<\varepsilon_v,
\]
where \(\varepsilon_v\) is the constant in
\cref{p0:thm:velocity-ckn}.  The velocity \(\varepsilon\)-regularity theorem gives
boundedness of \(u\) in \(Q_{r/2}(x_0,T)\), contradicting
\(x_0\in\Sigma(T)\).  Hence the limsup is positive; taking any
\(\eta_v\) below it and passing to a sequence with
\(C((x_0,T),r_n)\ge\eta_v\) gives the result.
\end{proof}

\begin{proposition}[Concentration consequences at a terminal singular point]
\label{p0:prop:concentration-package}
Let \((u,p)\) be a suitable weak solution on \(\mathbb R^3\times(T-\delta,T)\)
and let \(x_*\in\Sigma(T)\).  Then:
\begin{enumerate}[label=\textup{(\arabic*)}]
\item \(\limsup_{r\downarrow0}\bigl(C+D\bigr)((x_*,T),r)>0\) by
\Cref{p0:thm:no-concentration};
\item \(\limsup_{r\downarrow0}C((x_*,T),r)>0\) by
\Cref{p0:cor:velocity-concentration};
\item if the local pointwise Type I bound
\eqref{p0:eq:paper0-local-typeI-bound} holds at \(z_*=(x_*,T)\) (or
the global pointwise Type I bound holds), then
\Cref{p0:thm:local-weak-serrin-typeI} (respectively
\Cref{p0:thm:auto-velocity-no-escape}) supplies an admissible local Type I
concentration sequence in the sense of
\Cref{p0:def:paper0-admissible-local-typeI-sequence} suitable for the
Seregin extraction.
\end{enumerate}
\end{proposition}

\begin{proof}
Items (1) and (2) are restatements of the cited results at the singular
point \(x_*\).  Item (3) combines (2) with the no-escape conclusion in
\Cref{p0:thm:local-weak-serrin-typeI} or
\Cref{p0:thm:auto-velocity-no-escape}, together with
\Cref{p0:cor:paper0-local-typeI-admissible} for the local entry and
\Cref{p0:cor:auto-paperI-admissibility} for the global entry.
\end{proof}

\section{No-Escape for Finite-Energy Type I Singularities}
\label{p0:sec:auto-no-escape}

The preceding concentration result gives the positive velocity concentration
sequence used in the Type I blow-up analysis.  In the Seregin extraction
theorem \Cref{p1:thm:extraction}, the
additional compactness hypothesis is a no-escape condition excluding the possibility
that all of the selected scale-invariant velocity mass collapses into the
final rescaled time slice.  The next theorem shows that this no-escape condition is
forced by the finite-energy Type I hypotheses.

For this section we use, in addition to \(C\) and \(D\), the standard local
kinetic-energy and dissipation quantities
\[
  A(z_0,r)
  =r^{-1}\operatorname*{ess\,sup}_{t_0-r^2<t<t_0}
    \int_{B_r(x_0)} |u(x,t)|^2\,dx,
\]
and
\[
  E(z_0,r)
  =r^{-1}\int_{t_0-r^2}^{t_0}\int_{B_r(x_0)}
    |\nabla u(x,t)|^2\,dx\,dt .
\]
These quantities are invariant under the Navier--Stokes scaling.

\begin{definition}[Finite-energy Type I terminal singularity]
\label{p0:def:finite-energy-typeI-paper0}
Let \((u,p)\) be a suitable weak solution on \(\mathbb R^3\times(0,T)\).
We say that \(z_*=(x_*,T)\) is a finite-energy Type I terminal singularity if
\(x_*\in\Sigma(T)\),
\[
  u\in L^\infty(0,T;L^2(\mathbb R^3))
  \cap L^2(0,T;\dot H^1(\mathbb R^3)),
\]
and there are \(T_0<T\) and \(M<\infty\) such that
\begin{equation}\label{p0:eq:paper0-global-typeI}
  \|u(t)\|_{L^\infty(\mathbb R^3)}
  \le \frac{M}{\sqrt{T-t}},
  \qquad T_0<t<T .
\end{equation}
\end{definition}

The main estimate is a uniformly local energy propagation bound.  It is stated
at unit scale because the application is obtained by parabolic rescaling from
\((x_*,T)\).  The estimate is a local differential inequality: the unknown is the
uniformly-local kinetic energy, and the Type I envelope turns the nonlinear
pressure contribution into an integrable linear coefficient.  No bound below
depends on the size of the global \(L^2\) energy or the global dissipation;
finite energy is used to place the pressure in the standard whole-space
Leray gauge and to provide fixed-scale terminal \(L^3\) absolute continuity
for each fixed rescaling.

The standard Leray pressure gauge is used throughout this section.  At almost
every terminal time on which the Type I bound holds,
\(u(t)\in L^2(\mathbb R^3)\cap L^\infty(\mathbb R^3)\), hence
\(u_i(t)u_j(t)\in L^{3/2}(\mathbb R^3)\).  The whole-space pressure
\[
  p^L(t)=\mathcal R_i\mathcal R_j(u_i(t)u_j(t))
\]
is then defined by the Calderon--Zygmund theory.  The distributional pressure
equation gives
\(-\Delta p=\partial_i\partial_j(u_i u_j)=-\Delta p^L\), and the whole-space
de Rham argument applied to the momentum equation shows that \(p-p^L\) is a
function of time.  Thus replacing \(p\) by \(p^L\) is only a pressure-gauge
change and does not alter suitability.  This gauge choice uses the
finite-energy setting to define the representative; the uniform constants
below, in particular \(B_{\mathrm{uloc}}\) and \(B_A\), do not contain
\(\|u\|_{L^\infty(0,T;L^2(\mathbb R^3))}\) or
\(\int_0^T\|\nabla u(t)\|_{L^2(\mathbb R^3)}^2\,dt\).

\begin{lemma}[Uniformly-local energy propagation under a Type I envelope]
\label{p0:lem:uloc-typeI-propagation}
Let \((v,\pi)\) be a suitable weak solution on
\(\mathbb R^3\times(-4,0)\) with \(v(s)\in L^2(\mathbb R^3)\) for a.e.
\(s\), and take \(\pi\) to be the whole-space Leray pressure representative,
modulo functions of time.  Assume that
\begin{equation}\label{p0:eq:unit-typeI-envelope}
  |v(x,s)|\le M(-s)^{-1/2},
  \qquad \text{for a.e. }(x,s)\in\mathbb R^3\times(-4,0),
\end{equation}
and that for some admissible initial time \(s_0\in(-4,-3)\),
\begin{equation}\label{p0:eq:unit-uloc-initial}
  \alpha_0
  :=\sup_{a\in\mathbb R^3}\int_{B_2(a)} |v(x,s_0)|^2\,dx
  \le K .
\end{equation}
Then there is a constant \(B_{\mathrm{uloc}}=B_{\mathrm{uloc}}(M,K)<\infty\)
such that
\begin{equation}\label{p0:eq:unit-uloc-bound}
  \operatorname*{ess\,sup}_{-1<s<0}\sup_{a\in\mathbb R^3}
  \int_{B_1(a)} |v(x,s)|^2\,dx
  +
  \sup_{a\in\mathbb R^3}\int_{-1}^{0}\int_{B_1(a)}
  |\nabla v|^2\,dx\,ds
  \le B_{\mathrm{uloc}} .
\end{equation}
\end{lemma}

\begin{proof}
The estimate supplies the terminal no-escape bound for
\Cref{p0:thm:auto-velocity-no-escape}.  Let
\[
  \alpha(s)=\sup_{a\in\mathbb R^3}\int_{B_1(a)} |v(x,s)|^2\,dx .
\]
By \eqref{p0:eq:unit-typeI-envelope},
\begin{equation}\label{p0:eq:alpha-typeI-trivial}
  \alpha(s)^{1/2}\le C M(-s)^{-1/2},
  \qquad -4<s<0 .
\end{equation}
Fix \(a\in\mathbb R^3\) and choose
\(\phi_a\in C_c^\infty(B_2(a))\) with
\(0\le\phi_a\le1\), \(\phi_a\equiv1\) on \(B_1(a)\), and
\(|\nabla\phi_a|+|\Delta\phi_a|\le C\), uniformly in \(a\).  Applying the
local energy inequality on \((s_0,t)\), using the standard monotone
approximation of temporal cutoffs and with the pressure gauge chosen below,
gives for a.e. \(t\in(s_0,0)\)
\begin{align}
 &\int |v(t)|^2\phi_a\,dx
   +2\int_{s_0}^{t}\int |\nabla v|^2\phi_a\,dx\,ds       \notag \\
 &\quad\le
   \int |v(s_0)|^2\phi_a\,dx
   +C\int_{s_0}^{t}\int_{B_2(a)} |v|^2\,dx\,ds
   +C\int_{s_0}^{t}\int_{B_2(a)} |v|^3\,dx\,ds          \notag \\
 &\qquad
   +C\int_{s_0}^{t}\int_{B_2(a)}
       |\pi-c_a(s)|\,|v|\,dx\,ds .
\label{p0:eq:lei-unit-estimate}
\end{align}
Here \(c_a(s)\) is an arbitrary function of time; adding or subtracting such a
function does not change the local energy inequality, since
\(\nabla\cdot v=0\) and \(\phi_a\) is compactly supported.

The velocity terms are controlled by the uniformly local kinetic energy.  Since \(B_2(a)\) is covered by
finitely many unit balls,
\[
  \int_{B_2(a)} |v|^2\,dx\le C\alpha(s),
\]
and by \eqref{p0:eq:unit-typeI-envelope},
\begin{equation}\label{p0:eq:cubic-unit-bound}
  \int_{B_2(a)} |v|^3\,dx
  \le M(-s)^{-1/2}\int_{B_2(a)} |v|^2\,dx
  \le C M(-s)^{-1/2}\alpha(s).
\end{equation}
It remains to estimate the pressure term uniformly in \(a\).

Use the whole-space pressure representative
\(\pi=\mathcal R_i\mathcal R_j(v_i v_j)\), modulo functions of time.  Let
\(\chi_a\in C_c^\infty(B_8(a))\), \(\chi_a\equiv1\) on \(B_4(a)\), and write
inside \(B_4(a)\)
\[
  \pi=q_a+h_a,
  \qquad
  q_a=\mathcal R_i\mathcal R_j(\chi_a v_i v_j).
\]
Then \(h_a\) is harmonic in \(B_4(a)\).  Calderon--Zygmund boundedness
\cite{Stein1970} and \eqref{p0:eq:unit-typeI-envelope} give, for a.e. \(s\),
\begin{align*}
  \|q_a(s)\|_{L^2(B_4(a))}
  &\le C\|\chi_a v(s)\otimes v(s)\|_{L^2(\mathbb R^3)}  \\
  &\le C M(-s)^{-1/2}\|v(s)\|_{L^2(B_8(a))}             \\
  &\le C M(-s)^{-1/2}\alpha(s)^{1/2},
\end{align*}
where the last step again uses a finite covering by unit balls.  Hence
\begin{equation}\label{p0:eq:local-pressure-unit}
  \int_{B_2(a)} |q_a|\,|v|\,dx
  \le C M(-s)^{-1/2}\alpha(s).
\end{equation}

For the harmonic part, set \(c_a(s)=h_a(a,s)\).  The kernel estimate
\[
  |K_{ij}(x-z)-K_{ij}(a-z)|
  \le C |x-a|\,|z-a|^{-4},
  \qquad x\in B_2(a),\quad |z-a|>4,
\]
where \(K_{ij}\) is the kernel of \(\mathcal R_i\mathcal R_j\), gives
\[
  |h_a(x,s)-c_a(s)|
  \le
  C\sum_{m=2}^{\infty}2^{-4m}
    \int_{2^m<|z-a|<2^{m+1}} |v(z,s)|^2\,dz .
\]
The annulus \(\{2^m<|z-a|<2^{m+1}\}\) is covered by at most \(C2^{3m}\)
unit balls.  Therefore
\begin{equation}\label{p0:eq:harmonic-pressure-unit}
  \|h_a(s)-c_a(s)\|_{L^\infty(B_2(a))}
  \le C\sum_{m=2}^{\infty}2^{-m}\alpha(s)
  \le C\alpha(s).
\end{equation}
Consequently,
\begin{align}
  \int_{B_2(a)} |h_a-c_a|\,|v|\,dx
  &\le C\alpha(s)\int_{B_2(a)} |v|\,dx                    \notag\\
  &\le C\alpha(s)^{3/2}                                    \notag\\
  &\le C M(-s)^{-1/2}\alpha(s),
\label{p0:eq:harmonic-pressure-flux-unit}
\end{align}
where the last inequality is precisely the Type I absorbing estimate
\eqref{p0:eq:alpha-typeI-trivial}.
Combining \eqref{p0:eq:local-pressure-unit} and
\eqref{p0:eq:harmonic-pressure-flux-unit},
\begin{equation}\label{p0:eq:pressure-flux-unit}
  \int_{B_2(a)} |\pi-c_a(s)|\,|v|\,dx
  \le C M(-s)^{-1/2}\alpha(s).
\end{equation}

Insert \eqref{p0:eq:cubic-unit-bound} and \eqref{p0:eq:pressure-flux-unit} into
\eqref{p0:eq:lei-unit-estimate}, then take the supremum over \(a\).  Since
\((-s)^{-1/2}\in L^1(-4,0)\), Gronwall's inequality yields
\[
  \operatorname*{ess\,sup}_{s_0<s<0}\alpha(s)
  \le C K
     \exp\left(C\int_{s_0}^{0}\bigl(1+M(-s)^{-1/2}\bigr)\,ds\right)
  \le B_1(M,K).
\]
Returning to \eqref{p0:eq:lei-unit-estimate}, using this bound, and then letting
\(t\uparrow0\) through admissible times gives the corresponding
uniformly-local dissipation estimate on \((-1,0)\).  This proves
\eqref{p0:eq:unit-uloc-bound}.
\end{proof}

\begin{remark}[Local energy estimate as an absorbing bound]
\label{p0:rem:state-space-stratification}
The preceding proof can be read as a three-term estimate for the local quantity
\(\alpha(s)\).  The diffusion and transport terms are linear in
\(\alpha\), with coefficient \(1+M(-s)^{-1/2}\).  The local pressure term is
also linear after Calderon--Zygmund and the Type I envelope.  The remaining
nonlinear term is the harmonic pressure flux, which is
\(O(\alpha^{3/2})\); the pointwise Type I envelope implies
\(\alpha^{1/2}\le C M(-s)^{-1/2}\), reducing it to the same integrable linear
coefficient.  Thus large values of \(\alpha\) cannot persist: the resulting
differential inequality bounds \(\alpha\) in terms of constants depending
only on \(M\) and the initial uniformly local energy level.
\end{remark}

\begin{proposition}[Terminal kinetic tightness]
\label{p0:prop:auto-terminal-A}
Let \((u,p,z_*)\), \(z_*=(x_*,T)\), be a finite-energy Type I terminal
singularity in the sense of \cref{p0:def:finite-energy-typeI-paper0}.  Then there
exist \(r_0>0\) and \(B_A<\infty\), with \(B_A\) depending only on the Type I
constant \(M\), such that
\begin{equation}\label{p0:eq:auto-A-bound}
  A(z_*,r)\le B_A,
  \qquad 0<r<r_0 .
\end{equation}
Moreover,
\begin{equation}\label{p0:eq:auto-E-bound}
  r^{-1}\int_{T-r^2}^{T}\int_{B_r(x_*)}|\nabla u|^2\,dx\,dt
  \le B_A,
  \qquad 0<r<r_0 .
\end{equation}
\end{proposition}

\begin{proof}
Choose \(r_0>0\) so small that \(T-4r^2>T_0\) for \(0<r<r_0\).  For such
\(r\), use the whole-space Leray pressure gauge fixed above and define the
rescaled solution
\[
  v^{(r)}(y,s)=r u(x_*+ry,T+r^2s),
  \qquad
  \pi^{(r)}(y,s)=r^2 p(x_*+ry,T+r^2s),
\]
on \(\mathbb R^3\times(-4,0)\).  The Type I bound
\eqref{p0:eq:paper0-global-typeI} gives
\[
  |v^{(r)}(y,s)|\le M(-s)^{-1/2},
  \qquad -4<s<0 .
\]
For any admissible \(s_0\in(-4,-3)\) at which the Type I bound holds, and any
\(a\in\mathbb R^3\),
\begin{align*}
  \int_{B_2(a)} |v^{(r)}(y,s_0)|^2\,dy
  &= r^{-1}\int_{B_{2r}(x_*+ra)}
       |u(x,T+r^2s_0)|^2\,dx                                    \\
  &\le r^{-1}|B_{2r}|\frac{M^2}{r^2|s_0|}
   \le C M^2 .
\end{align*}
Such times form a full-measure subset of \((-4,-3)\).  The preceding estimate
is uniform in this choice, so \cref{p0:lem:uloc-typeI-propagation} applies with
\(K=CM^2\), uniformly in \(r\).  In particular \(K\) depends only on \(M\),
so \(B_A=B_{\mathrm{uloc}}(M,CM^2)\) depends only on \(M\).  Hence
\[
  \operatorname*{ess\,sup}_{-1<s<0}\int_{B_1}|v^{(r)}(y,s)|^2\,dy
  +\int_{-1}^{0}\int_{B_1}|\nabla v^{(r)}|^2\,dy\,ds
  \le B_A .
\]
Undoing the parabolic change of variables gives precisely
\eqref{p0:eq:auto-A-bound} and \eqref{p0:eq:auto-E-bound}.  The constant \(B_A\)
depends only on \(M\); no global energy norm of \(u\) has been used in the
estimate.
\end{proof}

\begin{theorem}[Global pointwise Type I implies finite-energy velocity no-escape]
\label{p0:thm:auto-velocity-no-escape}
Let \((u,p,z_*)\), \(z_*=(x_*,T)\), be a finite-energy Type I terminal
singularity.  Let \(\lambda_k\downarrow0\) be any sequence of scales, and set
\[
  u_k(x,t)=\lambda_k u(x_*+\lambda_k x,T+\lambda_k^2 t).
\]
Then
\begin{equation}\label{p0:eq:auto-no-escape}
  \lim_{\sigma\downarrow0}
  \sup_k\int_{-\sigma}^{0}\int_{B_1}|u_k(x,t)|^3\,dx\,dt=0 .
\end{equation}
\end{theorem}

\begin{proof}
By \cref{p0:prop:auto-terminal-A}, after discarding at most finitely many indices,
\[
  \operatorname*{ess\,sup}_{-1<t<0}\int_{B_1}|u_k(x,t)|^2\,dx\le B_A .
\]
The Type I bound gives, for the same tail of the sequence,
\[
  \|u_k(t)\|_{L^\infty(\mathbb R^3)}
  \le M(-t)^{-1/2},
  \qquad -1<t<0 .
\]
Therefore, for \(0<\sigma<1\),
\[
\begin{aligned}
  \int_{-\sigma}^{0}\int_{B_1}|u_k|^3\,dx\,dt
  &\le
  \int_{-\sigma}^{0}\|u_k(t)\|_{L^\infty(B_1)}
       \int_{B_1}|u_k(x,t)|^2\,dx\,dt  \\
  &\le M B_A\int_{-\sigma}^{0}(-t)^{-1/2}\,dt
   =2M B_A\sigma^{1/2} .
\end{aligned}
\]
This bound is uniform over the tail.  For each of the finitely many discarded
indices \(k\), choose \(\sigma_k>0\) so that
\(T+\lambda_k^2t>T_0\) for \(-\sigma_k<t<0\).  On this terminal subinterval
the Type I envelope and the finite-energy assumption imply
\[
  \int_{-\sigma_k}^{0}\int_{B_1}|u_k|^3\,dx\,dt
  \le
  \int_{-\sigma_k}^{0}\|u_k(t)\|_{L^\infty(B_1)}
       \int_{B_1}|u_k(x,t)|^2\,dx\,dt
  <\infty .
\]
Indeed, on \((-\sigma_k,0)\),
\[
  \|u_k(t)\|_{L^\infty(B_1)}\le M(-t)^{-1/2},
  \qquad
  \int_{B_1}|u_k(x,t)|^2\,dx
  \le \lambda_k^{-1}\|u(T+\lambda_k^2t)\|_{L^2(\mathbb R^3)}^2,
\]
and the right-hand side is finite for this fixed \(k\) by finite energy.
Thus each discarded index has absolute continuity of its \(L^3\)-integral at
the terminal slice.  Taking the maximum over this finite set gives the same
limit for the full supremum over \(k\).  This proves
\eqref{p0:eq:auto-no-escape}.
\end{proof}

\begin{corollary}[Admissibility of finite-energy Type I concentration sequences]
\label{p0:cor:auto-paperI-admissibility}
Let \((u,p,z_*)\), \(z_*=(x_*,T)\), be a finite-energy Type I terminal
singularity.  Then there exist \(\eta_v>0\) and \(\lambda_k\downarrow0\) such
that
\[
  C(z_*,\lambda_k)\ge\eta_v
  \qquad\text{for every }k,
\]
and the corresponding rescaled velocities satisfy the no-escape condition
\eqref{p0:eq:auto-no-escape}.  Hence the concentration and no-escape
admissibility condition used in the Type I blow-up limit theorem
\Cref{p1:thm:extraction} follows for finite-energy Type I singularities.
\end{corollary}

\begin{proof}
Since \(x_*\in\Sigma(T)\), \cref{p0:cor:velocity-concentration} gives
\(\eta_v>0\) and \(\lambda_k\downarrow0\) with
\(C(z_*,\lambda_k)\ge\eta_v\) for all \(k\).  The no-escape condition for this
same sequence follows from \cref{p0:thm:auto-velocity-no-escape}.
\end{proof}

\section{The Type I/Type II Blow-Up Rate Dichotomy}

Let \(x_0\in\Sigma(T)\), \(\eta>0\), and \(r_n\downarrow0\) be given by
\cref{p0:thm:singular-positive}.  Define
\[
  u_n(y,s)=r_nu(x_0+r_ny,T+r_n^2s),
  \qquad
  p_n(y,s)=r_n^2p(x_0+r_ny,T+r_n^2s).
\]
Then \((u_n,p_n)\) are suitable weak solutions on expanding backward
cylinders, and by \cref{p0:prop:critical-scaling} they retain a positive lower
bound for the mean-subtracted local scale-invariant quantity on a fixed
bounded cylinder.

\begin{definition}[Type I concentration sequence]
\label{p0:def:typeI-concentration}
Let \((u,p)\) be a suitable weak solution on
\(\mathbb R^3\times(T-\delta,T)\), and let \(r_n\downarrow0\) be a positive
local concentration sequence at \((x_0,T)\).  The sequence is called Type I
if the concentration point \((x_0,T)\) satisfies a local Type I
scale-invariant bound:
there are \(\rho>0\) and \(M<\infty\), with \(\rho^2<\delta\), such that
\[
  \operatorname*{ess\,sup}_{T-\rho^2<t<T}
  \sqrt{T-t}\,\|u(t)\|_{L^\infty(B_\rho(x_0))}\le M.
\]
This is the Type I alternative.
\end{definition}

\begin{definition}[Local Type II alternative]
\label{p0:def:typeII-alternative}
Let \((u,p)\) be a suitable weak solution on
\(\mathbb R^3\times(T-\delta,T)\), and let \(r_n\downarrow0\) be a positive
local concentration sequence at \((x_0,T)\).  The sequence belongs to the local
Type II alternative if the concentration point \((x_0,T)\) satisfies no
local Type I scale-invariant bound.  Equivalently, for every \(\rho>0\) with
\(\rho^2<\delta\),
\[
  \operatorname*{ess\,sup}_{T-\rho^2<t<T}
  \sqrt{T-t}\,\|u(t)\|_{L^\infty(B_\rho(x_0))}=\infty.
\]
\end{definition}

\begin{theorem}[Local blow-up rate dichotomy]\label{p0:thm:typeI-typeII-dichotomy}
Let \((u,p)\) be a suitable weak solution on
\(\mathbb R^3\times(T-\delta,T)\), and suppose \(\Sigma(T)\ne\emptyset\).
Then there exist \(x_0\in\Sigma(T)\),
\(\eta>0\), and \(r_n\downarrow0\) such that
\[
  C((x_0,T),r_n)+D((x_0,T),r_n)\ge\eta
  \qquad\text{for all }n.
\]
The sequence \(r_n\) is a positive local concentration sequence.  The
corresponding concentration regime satisfies either the local Type I bound or
the local Type II alternative.  Consequently, once the vanishing case has been
excluded by \(\varepsilon\)-regularity, every finite-time singular case belongs
to precisely one of the two blow-up rate alternatives.
\end{theorem}

\begin{proof}
If the local scale-invariant quantities vanished at every point of
\(\Sigma(T)\), then \cref{p0:thm:no-concentration} would give
\(\Sigma(T)=\emptyset\), contrary to the hypothesis.  Therefore
\cref{p0:thm:singular-positive} gives
\(x_0\in\Sigma(T)\), \(\eta>0\), and \(r_n\downarrow0\) with the stated lower
bound.  The rescaled sequence has a positive local scale-invariant lower bound
by \cref{p0:prop:critical-scaling}.  By definition, either the Type I bound in
\cref{p0:def:typeI-concentration} holds, or the local Type II alternative in
\cref{p0:def:typeII-alternative} holds.
\end{proof}

\begin{remark}
This theorem is local.  Its conclusion is the existence of a positive
mean-subtracted scale-invariant concentration sequence and the dichotomy
between a local Type I bound and failure of every such bound at the
concentration point.
\end{remark}

\section{Local Pointwise Type I Reduction to Seregin Extraction}
\label{p0:sec:local-typeI-terminal-reduction-paper0}

The following consequence of the weak-Serrin endpoint argument is used.
The local pointwise Type I alternative in
\cref{p0:def:typeI-concentration} is strong enough, after rescaling at the singular
point, to give compactness on every compact negative-time cylinder.  The
endpoint weak-Serrin implication reads:
for suitable weak solutions the local pointwise Type I envelope belongs to
\(L^{2,\infty}_tL^\infty_x\), hence the Albritton--Barker weak Serrin local
Type I estimate gives the scale-invariant \(A\)-bound on smaller cylinders and
therefore terminal velocity no-escape.  Thus positive local Type I
concentration sequences are admissible for the local Seregin extraction in
\Cref{p0:lem:local-seregin-extraction}.

Throughout this section \(z_*=(x_*,T)\), and
\[
  Q_1(0,0)=B_1(0)\times(-1,0).
\]

\begin{definition}[Admissible local Type I concentration sequence]
\label{p0:def:paper0-admissible-local-typeI-sequence}
For a local pointwise Type I point \(z_*=(x_*,T)\), a sequence
\(r_n\downarrow0\) is an admissible local Type I concentration sequence if the rescaled
velocities
\[
  u_n(x,t)=r_nu(x_*+r_nx,T+r_n^2t)
\]
satisfy, for some \(\eta_v>0\),
\[
   \int_{Q_1(0,0)} |u_n|^3\,dx\,dt\ge\eta_v
   \qquad\text{for every }n,
\]
and
\[
   \lim_{\sigma\downarrow0}
   \sup_n
   \int_{-\sigma}^{0}\int_{B_1}|u_n|^3\,dx\,dt=0 .
\]
\end{definition}

\begin{proposition}[Local pointwise Type I rescalings have terminal compactness]
\label{p0:prop:paper0-local-typeI-terminal-compactness}
Let \((u,p)\) be a suitable weak solution on
\(\mathbb R^3\times(0,T)\) satisfying the finite-energy bounds
\[
  u\in L^\infty(0,T;L^2(\mathbb R^3))
  \cap L^2(0,T;\dot H^1(\mathbb R^3)).
\]
Let \(z_*=(x_*,T)\) be a singular point, and assume that \(z_*\) satisfies the
local pointwise Type I bound: there exist \(\rho>0\) and \(M<\infty\) such that
\begin{equation}\label{p0:eq:paper0-local-typeI-bound}
  \operatorname*{ess\,sup}_{T-\rho^2<t<T}
  \sqrt{T-t}\,\|u(t)\|_{L^\infty(B_\rho(x_*))}\le M .
\end{equation}
For any sequence \(r_n\downarrow0\), define
\begin{equation}\label{p0:eq:paper0-local-typeI-rescaling}
  u_n(x,t)=r_nu(x_*+r_nx,T+r_n^2t),
  \qquad
  p_n(x,t)=r_n^2p(x_*+r_nx,T+r_n^2t).
\end{equation}
Then the following hold.
\begin{enumerate}[label=\textup{(\roman*)}]
\item If
\[
  K=B_R\times[-S,-\sigma]\Subset\mathbb R^3\times(-\infty,0),
  \qquad R,S<\infty,\quad \sigma>0,
\]
then, for all sufficiently large \(n\),
\begin{equation}\label{p0:eq:paper0-local-typeI-Linfty}
  \|u_n\|_{L^\infty(K)}\le M\sigma^{-1/2}.
\end{equation}
\item On each such compact cylinder \(K\), after subtracting time-dependent
local pressure gauges \(a_{n,K}(t)\), the sequence satisfies uniform local
Seregin bounds
\begin{equation}\label{p0:eq:paper0-local-typeI-Seregin-bounds}
  \|u_n\|_{L^\infty_tL^2_x(K)}
  +\|\nabla u_n\|_{L^2(K)}
  +\|p_n-a_{n,K}(t)\|_{L^{3/2}(K)}
  \le C_K .
\end{equation}
Here \(C_K\) may depend on \(K,M,\rho\), and on the finite Leray energy level
of the original solution, but it does not depend on \(n\) and does not use a
terminal no-escape assumption.
\item After passing to a subsequence, \((u_n,p_n-a_{n,K})\) converges on every
compact \(K\Subset\mathbb R^3\times(-\infty,0)\) to an ancient suitable weak
solution \((U,Q)\), with
\[
  u_n\to U\quad\text{strongly in }L^q(K),\qquad 1\le q<\infty,
\]
\[
  p_n-a_{n,K}(t)\rightharpoonup Q
  \quad\text{weakly in }L^{3/2}(K),
\]
and
\[
  |U(x,t)|\le M(-t)^{-1/2},\qquad t<0.
\]
\end{enumerate}
\end{proposition}

\begin{proof}
Fix \(K=B_R\times[-S,-\sigma]\Subset\mathbb R^3\times(-\infty,0)\).  For
large \(n\),
\[
  r_nR<\rho/2,\qquad r_n^2S<\rho^2/2.
\]
Thus \(x_*+r_nx\in B_\rho(x_*)\) and
\(T+r_n^2t\in(T-\rho^2,T)\) for a.e. \((x,t)\in K\).  The local pointwise Type I bound
\eqref{p0:eq:paper0-local-typeI-bound} gives
\[
  |u_n(x,t)|
  =r_n|u(x_*+r_nx,T+r_n^2t)|
  \le r_nM(T-(T+r_n^2t))^{-1/2}
  =M(-t)^{-1/2}
  \le M\sigma^{-1/2}.
\]
This proves \eqref{p0:eq:paper0-local-typeI-Linfty}.  It also gives
uniform local \(L^\infty_tL^2_x\) and \(L^3_{x,t}\) bounds on \(K\).

It remains to explain the pressure and dissipation bounds.  At this point the
ambient solution is a suitable Leray--Hopf solution on
\(\R^3\times(0,T)\), not merely a local suitable weak solution.  Hence the
whole-space Leray pressure representative may be chosen, modulo functions of
time, as \(p=\mathcal R_i\mathcal R_j(u_i u_j)\).  All local pressure
estimates below are made for this representative, followed by local
time-dependent mean subtraction.  Choose
\(\chi_R\in C_c^\infty(B_{8R})\), \(\chi_R\equiv1\) on \(B_{4R}\), and write
inside \(B_{4R}\)
\[
  p_n=q_n+h_n,\qquad
  q_n=\mathcal R_i\mathcal R_j(\chi_Ru_{n,i}u_{n,j}).
\]
Calderon--Zygmund boundedness and the local \(L^3\)-bound for \(u_n\) give a
uniform \(L^{3/2}\)-bound for \(q_n\).  For \(h_n\), take the gauge
\(a_{n,K}(t)=h_n(0,t)\).  Let
\[
   A_m=\{2^mR<|z|<2^{m+1}R\},\qquad m\ge2 .
\]
For \(x\in B_{2R}\), the harmonic pressure kernel estimate gives
\[
  |h_n(x,t)-a_{n,K}(t)|
  \le
  C R\sum_{m\ge2}(2^mR)^{-4}
  \int_{A_m}|u_n(z,t)|^2\,dz .
\]
Choose \(m_n\) by
\[
   2^{m_n}Rr_n\le \rho <2^{m_n+1}Rr_n .
\]
After discarding finitely many indices, assume \(m_n\ge2\).
If \(m\le m_n-1\), then the physical image of \(A_m\) is contained in
\(B_\rho(x_*)\).  On the time slab \([-S,-\sigma]\), the local Type I bound
therefore gives
\[
  \int_{A_m}|u_n(z,t)|^2\,dz
  \le C M^2\sigma^{-1}(2^mR)^3 .
\]
Consequently the inner part of the harmonic tail is bounded by
\[
\begin{aligned}
  C R\sum_{m\le m_n-1}(2^mR)^{-4}
  \int_{A_m}|u_n(z,t)|^2\,dz
  &\le C R M^2\sigma^{-1}\sum_{m\le m_n-1}(2^mR)^{-1} \\
  &\le C M^2\sigma^{-1}.
\end{aligned}
\]
For \(m\ge m_n\), the finite Leray energy gives
\[
  \int_{A_m}|u_n(z,t)|^2\,dz
  \le
  \int_{\mathbb R^3}|u_n(z,t)|^2\,dz
  =
  r_n^{-1}\int_{\mathbb R^3}|u(y,T+r_n^2t)|^2\,dy .
\]
Since \(2^{m_n}R\simeq \rho/r_n\), the outer part satisfies
\[
\begin{aligned}
  C R\sum_{m\ge m_n}(2^mR)^{-4}
  \int_{A_m}|u_n(z,t)|^2\,dz
  &\le
  C R r_n^{-1}\|u\|_{L^\infty(0,T;L^2)}^2
  \sum_{m\ge m_n}(2^mR)^{-4}  \\
  &\le
  C_{\rho,R}\,r_n^3\|u\|_{L^\infty(0,T;L^2)}^2 .
\end{aligned}
\]
Thus \(h_n-a_{n,K}\) is bounded in
\(L^\infty([-S,-\sigma];L^\infty(B_{2R}))\), uniformly in \(n\).  Together
with the \(L^{3/2}\)-bound for \(q_n\), this gives the
\(L^{3/2}(K)\) pressure estimate in
\eqref{p0:eq:paper0-local-typeI-Seregin-bounds}.

Apply the local energy inequality for \((u_n,p_n)\) with a cutoff equal to one
on \(K\) and supported in a slightly larger compact cylinder.  The local
kinetic, cubic velocity, and pressure-flux terms are bounded by the preceding
estimates, so \(\|\nabla u_n\|_{L^2(K)}\le C_K\).  The Navier--Stokes
equation then bounds \(\partial_tu_n\) locally in a negative Sobolev space.
Aubin--Lions compactness and the pressure bound give the asserted Seregin
subsequence convergence.  Passing to the limit in the equations and local
energy inequality gives an ancient suitable weak solution, and the pointwise
Type I bound passes to the limit.
\end{proof}

\begin{lemma}[Terminal \(A\)-bound from the local weak-Serrin estimate]
\label{p0:lem:terminal-A-from-weak-serrin}
Let \((u,p)\) be a suitable weak solution in a backward cylinder
\(Q_\rho(z_*)\), \(z_*=(x_*,T)\), and assume the local pointwise Type I bound
\eqref{p0:eq:paper0-local-typeI-bound}.  Then, for every \(0<\theta<1\),
there is \(K_\theta<\infty\) such that, for all sufficiently small \(r\),
\[
   r^{-1}
   \esssup_{T-r^2<t<T}
   \int_{B_r(x_*)}|u(x,t)|^2\,dx
   \le K_\theta .
\]
\end{lemma}

\begin{proof}
After translating and scaling, assume \(z_*=(0,0)\) and work in
\[
   Q_\rho=B_\rho\times(-\rho^2,0).
\]
Put
\[
   f(t)=\|u(t)\|_{L^\infty(B_\rho)} .
\]
The pointwise Type I bound gives \(f(t)\le M(-t)^{-1/2}\).  Hence, for every
\(\alpha>0\),
\[
   |\{t\in(-\rho^2,0): f(t)>\alpha\}|
   \le \min\{\rho^2,M^2\alpha^{-2}\},
\]
and therefore
\[
   \|f\|_{L^{2,\infty}(-\rho^2,0)}\le C M .
\]
Thus
\[
   u\in L_t^{2,\infty}L_x^\infty(Q_\rho).
\]

Fix \(0<\theta<1\).  Apply the Albritton--Barker weak-Serrin Type I lemma
\cite[Lemma 2.5]{AlbrittonBarker2019} with \(p=\infty\) and \(q=2\).  In
their notation,
\[
   I(\omega)
   =
   \sup_{Q'\Subset\omega}\bigl(A(Q')+C(Q')+D(Q')+E(Q')\bigr),
\]
and the endpoint weak-Serrin hypothesis above gives a finite bound
\[
   I(Q_{\theta\rho})\le K_\theta .
\]
Let \(r<(\theta\rho)/\sqrt{2}\) and \(0<\eps<r^2\).  Then
\[
   Q_r((0,-\eps))=B_r\times(-r^2-\eps,-\eps)\Subset Q_{\theta\rho},
\]
so the Albritton--Barker estimate gives
\[
   r^{-1}\esssup_{-r^2-\eps<t<-\eps}
   \int_{B_r}|u(x,t)|^2\,dx
   \le K_\theta .
\]
Since \((-r^2,-\eps)\subset(-r^2-\eps,-\eps)\), it follows that
\[
   r^{-1}\esssup_{-r^2<t<-\eps}
   \int_{B_r}|u(x,t)|^2\,dx
   \le K_\theta .
\]
As \(\eps\downarrow0\), the intervals \((-r^2,-\eps)\) increase to
\((-r^2,0)\).  By the definition of essential supremum over nested measurable
sets,
\[
   r^{-1}\esssup_{-r^2<t<0}
   \int_{B_r}|u(x,t)|^2\,dx
   \le K_\theta .
\]
Undoing the translation and scaling gives the stated estimate at \(z_*\).
\end{proof}

\begin{theorem}[Local endpoint weak-Serrin Type I estimate]
\label{p0:thm:local-weak-serrin-typeI}
Let \((u,p)\) be a suitable weak solution in a backward cylinder
\(Q_\rho(z_*)\), \(z_*=(x_*,T)\), and assume the local pointwise Type I bound
\eqref{p0:eq:paper0-local-typeI-bound}
\[
  \esssup_{T-\rho^2<t<T}
  \sqrt{T-t}\,\|u(t)\|_{L^\infty(B_\rho(x_*))}\le M.
\]
Then, after possibly shrinking \(\rho\), the rescaled velocities
\(u_n(y,s)=r_nu(x_*+r_ny,T+r_n^2s)\) along every scale sequence
\(r_n\downarrow0\) satisfy the terminal velocity no-escape condition
\[
   \lim_{\sigma\downarrow0}\sup_n
   \int_{-\sigma}^{0}\int_{B_1}|u_n(y,s)|^3\,dy\,ds=0 .
\]
\end{theorem}

\begin{proof}
After translating and scaling, assume \(z_*=(0,0)\) and work in
\[
   Q_\rho=B_\rho\times(-\rho^2,0).
\]
Let \(r_n\downarrow0\) and
\[
   u_n(y,s)=r_n u(r_ny,r_n^2s).
\]
Fix \(0<\theta<1\).  By \Cref{p0:lem:terminal-A-from-weak-serrin}, the terminal
\(A\)-quantity is bounded uniformly along the chosen scales.
For all sufficiently large \(n\),
\[
   \esssup_{-1<s<0}\int_{B_1}|u_n(y,s)|^2\,dy
   =
   r_n^{-1}\esssup_{-r_n^2<t<0}
   \int_{B_{r_n}}|u(x,t)|^2\,dx
   \le K_\theta .
\]
Moreover,
\[
   \|u_n(s)\|_{L^\infty(B_1)}
   \le M(-s)^{-1/2}.
\]
Therefore, for \(0<\sigma<1\),
\[
\begin{aligned}
   \int_{-\sigma}^{0}\int_{B_1}|u_n|^3\,dy\,ds
   &\le
   \int_{-\sigma}^{0}
   \|u_n(s)\|_{L^\infty(B_1)}
   \int_{B_1}|u_n(y,s)|^2\,dy\,ds  \\
   &\le
   M K_\theta \int_{-\sigma}^{0}(-s)^{-1/2}\,ds
   =
   2M K_\theta\sigma^{1/2}.
\end{aligned}
\]
This proves the required terminal velocity no-escape for all sufficiently
large \(n\).  The finitely many remaining indices are harmless by absolute
continuity of the local \(L^3\) integral, since suitable weak solutions are
locally in \(L^3\) on terminal subcylinders.  Taking the
supremum in \(n\) and then letting \(\sigma\downarrow0\) gives the claim.
\end{proof}

\begin{corollary}[Local pointwise Type I concentration sequences are admissible]
\label{p0:cor:paper0-local-typeI-admissible}
Under the hypotheses of
\cref{p0:prop:paper0-local-typeI-terminal-compactness}, let \(r_n\downarrow0\) be
any sequence such that the rescaled velocities satisfy
\[
  \int_{Q_1(0,0)} |u_n|^3\,dx\,dt\ge\eta_v>0
  \qquad\text{for all }n .
\]
Then \(r_n\) is an admissible local Type I concentration sequence in the sense of
\cref{p0:def:paper0-admissible-local-typeI-sequence}.
\end{corollary}

\begin{proof}
The displayed positive concentration lower bound is the first requirement in
\cref{p0:def:paper0-admissible-local-typeI-sequence}.  The
Albritton--Barker local weak-Serrin Type I estimate
\cite{AlbrittonBarker2019} gives terminal velocity no-escape from the local
pointwise Type I envelope for every rescaled sequence; the precise statement used here is
\Cref{p0:thm:local-weak-serrin-typeI}.  This gives the second requirement in
\cref{p0:def:paper0-admissible-local-typeI-sequence}.
\end{proof}

\begin{lemma}[Local Type I concentration enters the Seregin extraction]
\label{p0:lem:paper0-local-typeI-into-terminal-dichotomy}
Under the hypotheses of
\cref{p0:prop:paper0-local-typeI-terminal-compactness}, let \(r_n\downarrow0\) be
a velocity concentration sequence satisfying
\[
  \int_{Q_1(0,0)} |u_n|^3\,dx\,dt\ge\eta_v>0
  \qquad\text{for all }n .
\]
Then \(r_n\) is admissible by
\cref{p0:cor:paper0-local-typeI-admissible}.  There are
\(\sigma_0\in(0,1)\) and \(\eta_1>0\) such that
\begin{equation}\label{p0:eq:paper0-local-typeI-away-mass}
  \int_{-1}^{-\sigma_0}\int_{B_1}|u_n|^3\,dx\,dt\ge\eta_1
  \qquad\text{for all }n.
\end{equation}
Consequently, the rescaled sequence retains positive local \(L^3\) mass on a
compact negative-time cylinder.  In particular the associated Seregin
extraction has no loss of the local \(L^3\) mass needed for
nontriviality.
\end{lemma}

\begin{proof}
The no-escape condition in
\cref{p0:def:paper0-admissible-local-typeI-sequence} gives \(\sigma_0>0\) such
that the terminal layer \(B_1\times(-\sigma_0,0)\) carries at most half of the
fixed \(Q_1\) velocity mass.  This proves
\eqref{p0:eq:paper0-local-typeI-away-mass}.  The compactness part follows
from \cref{p0:prop:paper0-local-typeI-terminal-compactness} on the compact
negative-time cylinder \(B_1\times[-1,-\sigma_0]\).  Strong local \(L^3\)
convergence on this cylinder passes the positive mass to the ancient limit.
Terminal no-escape is used only at this step.

Thus the local Type I concentration sequence is admissible for the Seregin
extraction and retains positive localized \(L^3\)-mass away from the terminal
slice.
\end{proof}

\begin{lemma}[Local Seregin extraction from an admissible local Type I sequence]
\label{p0:lem:local-seregin-extraction}
Under the hypotheses of
\cref{p0:prop:paper0-local-typeI-terminal-compactness}, let \(r_n\downarrow0\)
be an admissible local Type I concentration sequence in the sense of
\cref{p0:def:paper0-admissible-local-typeI-sequence}.  Then, after passing to
a subsequence, the rescaled velocities produce a smooth bounded centered
solution \(V\) of \eqref{p1:eq:centered} on
\(\mathbb R^3\times\mathbb R\).  The solution \(V\) satisfies the
normalization used for Seregin ancient limits: there are a compact cylinder
\(K\Subset\mathbb R^3\times\mathbb R\) and \(\eta_0>0\) such that
\[
   \iint_K |V|^3\,dy\,d\tau\ge\eta_0 .
\]
\end{lemma}

\begin{proof}
\Cref{p0:prop:paper0-local-typeI-terminal-compactness} gives local Seregin
compactness for the rescaled sequence on every compact cylinder
\(K\Subset\mathbb R^3\times(-\infty,0)\).  The local Type I bound passes to
the ancient limit as
\[
   |U(x,t)|\le M(-t)^{-1/2},\qquad t<0.
\]
Interior Serrin regularity then gives smoothness of \(U\) on
\(\mathbb R^3\times(-\infty,0)\) by the argument used in
\Cref{p1:lem:smoothness}.  The centered change of variables
\[
   y=\frac{x}{\sqrt{-t}},\qquad \tau=-\log(-t),\qquad
   V(y,\tau)=\sqrt{-t}\,U(y\sqrt{-t},t)
\]
therefore gives a smooth bounded solution of the centered equation
\eqref{p1:eq:centered}.  The no-escape part of local admissibility and
\cref{p0:lem:paper0-local-typeI-into-terminal-dichotomy} give positive
compact \(L^3\)-mass for \(U\) away from \(t=0\).  On the compact image of
that cylinder under \((x,t)=(e^{-\tau/2}y,-e^{-\tau})\), the factor
\(e^{-\tau}\) in
\[
   |U|^3\,dx\,dt=e^{-\tau}|V|^3\,dy\,d\tau
\]
is bounded above and below by positive constants.  Hence the positive lower
bound for \(U\) becomes the displayed local \(L^3\) nonvanishing for \(V\).
\end{proof}

\begin{proposition}[Local Type I entry into the final assembly]
\label{p0:prop:local-typeI-entry-final-assembly}
Assume the residual-class hypothesis \Cref{p1:hyp:no-remainder}.
Let \((u,p)\) be a finite-energy suitable weak solution on
\(\mathbb R^3\times(0,T)\).  Then every terminal point \(z_*=(x_*,T)\) satisfying
the local pointwise Type I bound \eqref{p0:eq:paper0-local-typeI-bound}
is regular.  More precisely, if such a point were singular, it would generate
a raw extracted normalized Seregin ancient limit with compact spacetime
nonvanishing; \Cref{p1:lem:compact-to-invariant-nonzero} would then give the
invariant local nonvanishing used to form the associated generated descendant
hull.  That raw hull would be exhausted by the Liouville alternatives and the
residual class removed by \Cref{p1:hyp:no-remainder}.
\end{proposition}

\begin{proof}
Assume \Cref{p1:hyp:no-remainder}.
Suppose, toward a contradiction, that \(z_*=(x_*,T)\) is a singular point
satisfying \eqref{p0:eq:paper0-local-typeI-bound}.
Choose the positive \(Q_1\) velocity concentration sequence supplied by the
local concentration theorem.  By
\cref{p0:cor:paper0-local-typeI-admissible}, this sequence satisfies
terminal no-escape and is admissible.
By \cref{p0:lem:paper0-local-typeI-into-terminal-dichotomy}, the local Type I
sequence has local compactness and positive compact \(L^3\)-mass away from the
terminal slice.
By \cref{p0:lem:local-seregin-extraction}, the sequence generates a raw
extracted normalized Seregin ancient limit \(V\).  By
\Cref{p1:lem:compact-to-invariant-nonzero}, \(V\) satisfies the invariant
local nonvanishing condition for some data \((R_0,\eta_*)\).  Form the raw
generated descendant hull
\[
   \calS=\calS_{\mathrm{raw\text{-}gen}}(u,p,z_*;R_0,\eta_*).
\]
By construction, \(V\in\calS\).
By the class exhaustion in
\cref{p1:prop:exhaustion}, \(V\) lies either in one of the Liouville
classes excluded by the stated Liouville theorems or in the remaining class
\(\calR(\calS)\).  The Liouville classes are excluded by the small-amplitude,
stationary, uniformly tight, and structure-and-decay Liouville results; the
last of these is routed through \Cref{p3:thm:typeI-zero-hypotheses}.  The
residual alternative is removed by the residual-class
hypothesis applied to
\(\calS_{\mathrm{raw\text{-}gen}}(u,p,z_*;R_0,\eta_*)\).
This contradicts the generated nonvanishing guaranteed by
\Cref{p1:cor:normalized-nonzero}.  Hence \(z_*\) is regular.
\end{proof}

\begin{corollary}[Conditional exclusion of the local Type I alternative]
\label{p0:cor:paper0-local-typeII-after-local-typeI-reduction}
Assume the residual-class hypothesis \Cref{p1:hyp:no-remainder}.  Then the
local Type I alternative is excluded for every finite-energy singular point.
Hence every finite-energy singular point belongs to the local Type II
alternative in the dichotomy above.
\end{corollary}

\begin{proof}
This is \cref{p0:prop:local-typeI-entry-final-assembly} restated in the terminology
of the local Type I/local Type II dichotomy.
\end{proof}


\subsection*{Part I dependency diagram}

\begin{center}\small
\begin{tabular}{l}
singular point \(z_*=(x_*,T)\)\\
\(\downarrow\) \Cref{p0:cor:velocity-concentration} (CKN)\\
positive velocity concentration scales\\
\(\downarrow\) \Cref{p0:thm:local-weak-serrin-typeI} (local) \emph{or} \Cref{p0:thm:auto-velocity-no-escape} (global)\\
no-escape away from terminal slice\\
\(\downarrow\) \Cref{p0:cor:paper0-local-typeI-admissible} / \Cref{p0:cor:auto-paperI-admissibility}\\
admissible concentration sequence\\
\(\downarrow\) \Cref{p0:lem:local-seregin-extraction} \emph{or} \Cref{p1:thm:extraction}\\
raw extracted Seregin ancient limit, then raw generated state space and mild strata (Part II)\\
\(\downarrow\) class decomposition \Cref{p1:def:remaining-class}\\
non-residual classes (excluded here) \(\cup\) residual class\\
\(\downarrow\) \Cref{p1:hyp:no-remainder} (conditional residual closure hypothesis)\\
contradiction\(\Rightarrow\) \Cref{p1:thm:final-assembly}.
\end{tabular}
\end{center}

This is a reduction theorem, not a self-contained residual proof; the residual
closure assumption is isolated in \Cref{p1:hyp:no-remainder}.

\clearpage
\section*{Part II. Type I Blow-up Limits and the Residual Hypothesis}

\subsection*{Overview}

An admissible finite-energy Type I singularity for the
three-dimensional Navier--Stokes equations is first reduced to a raw extracted
nonzero bounded ancient solution of the centered self-similar Navier--Stokes
equation, and then to the raw generated state space used in the residual
argument.  The
construction keeps track of pressure gauges, local compactness, stability of
the local energy inequality, and the velocity concentration supplied by local
Caffarelli--Kohn--Nirenberg theory.  It yields a smooth bounded ancient solution
for the centered equation on \(\R^3\times\R\), together with a quantitative
local \(L^3\) nonvanishing condition.

The possible generated raw Seregin states are organized by the
analytic hypotheses used in the class decomposition: small amplitude,
stationary \(L^3\), uniform \(L^3\)-tightness, and structure or decay assumptions to which
cited Liouville theorems apply.  The small-amplitude case is excluded by the bounded ancient
perturbative Liouville theorem in \Cref{p1:lem:small-liouville}, and the stationary \(L^3\) case by
the theorem of Ne\v{c}as--R\r{u}\v{z}i\v{c}ka--\v{S}ver\'ak.  The
structure and decay class contains only hypotheses for which the cited Liouville
theorems force \(V\equiv0\).  The uniformly \(L^3\)-tight class is excluded in
\Cref{p2:thm:tight-liouville}.  The residual class is handled only through the
residual-class hypothesis \Cref{p1:hyp:no-remainder}.
Consequently the Type I conclusion obtained from this extraction is conditional
on that residual hypothesis within the raw generated finite-energy and local
pointwise Type I extraction classes.

\section{Introduction and main results}\label{p1:sec:introduction}

\subsection{Type I blow-up limits}

Let \(u:\R^3\times(0,T)\to\R^3\) solve the incompressible Navier--Stokes
equations
\begin{equation}\label{p1:eq:NS}
   \partial_tu+(u\cdot\nabla)u+\nabla p=\Delta u,\qquad
   \divergence u=0.
\end{equation}
The local regularity theory of Caffarelli--Kohn--Nirenberg \cite{CKN1982}
detects singular points through scale-invariant concentration of velocity and
pressure.  A terminal point \(z_*=(x_*,T)\) is of Type I if
\begin{equation}\label{p1:eq:type-I-bound}
   \limsup_{t\uparrow T}\sqrt{T-t}\,
   \|u(t)\|_{L^\infty(\R^3)}<\infty .
\end{equation}
The blow-up extraction uses a selected velocity concentration sequence from
the local concentration theorem and assumes only that this velocity
concentration does not escape into the terminal slice after rescaling.  The
resulting admissibility condition is defined in
\cref{p1:def:admissible-type-I}.  It replaces a stronger all-scale compactness
assumption on scale-invariant energy, pressure, and dissipation quantities by
the non-escape condition needed for nontriviality of the ancient limit.

\subsection{Normalized Seregin limits}

For scales \(\lambda\downarrow0\), the parabolic rescaling about \(z_*\) is
\[
   u^{(\lambda)}(x,t)
   =
   \lambda u(x_*+\lambda x,T+\lambda^2t),\qquad
   p^{(\lambda)}(x,t)
   =
   \lambda^2p(x_*+\lambda x,T+\lambda^2t).
\]
The rescaled time intervals exhaust \((-\infty,0)\).  The Type I bound
becomes a uniform bound
\[
   \|u^{(\lambda)}(t)\|_{L^\infty(\R^3)}
   \le M(-t)^{-1/2}
\]
on compact subintervals of \((-\infty,0)\), and compactness produces an
unrenormalized ancient solution \((U,Q)\).  The associated centered variables
are
\[
   y=\frac{x}{\sqrt{-t}},\qquad
   \tau=-\log(-t),\qquad
   V(y,\tau)=\sqrt{-t}\,U(y\sqrt{-t},t).
\]
They transform the ancient Navier--Stokes solution into the centered equation
\begin{equation}\label{p1:eq:centered}
   \partial_\tau V+(V\cdot\nabla)V+\nabla\Pi
   =
   \Delta V-\frac12(V+y\cdot\nabla V),
   \qquad \divergence V=0.
\end{equation}
The selected velocity concentration at the singular point persists through
the extraction and gives a compact cylinder \(K\Subset\R^3\times\R\) and a
constant \(\eta_0>0\) such that
\begin{equation}\label{p1:eq:nonvanishing}
   \iint_K |V|^3\,dy\,d\tau\ge\eta_0 .
\end{equation}

\subsection{Liouville alternatives and the remaining case}

No unconditional Liouville theorem is known for all bounded ancient solutions
of \eqref{p1:eq:centered}.  The reduction therefore uses a finite decomposition
based on explicit Liouville theorems.  A generated raw Seregin state is compared with the
following analytic classes:
\[
   \calC_{\mathrm{small}},\qquad
   \calC_{\mathrm{stat},L^3},\qquad
   \calC_{L^3\mathrm{-tight}},\qquad
   \calC_{\mathrm{sd}},\qquad
   \calR(\calS).
\]
They denote, respectively, small amplitude, stationary \(L^3\), uniform
\(L^3\)-tightness, and structure or decay hypotheses to which the cited
Liouville theorems apply
(\(\mathrm{sd}\)), and the complement of these four classes inside a raw
generated descendant hull
\(\calS=\calS_{\mathrm{raw\text{-}gen}}(u,p,z_*;R_0,\eta_*)\).  The last set carries no additional
analytic hypothesis; it is the remaining case after the stated Liouville
alternatives have been removed.

\subsection{Main theorems}

\begin{theorem}[Normalized Type I blow-up limit]\label{p1:thm:extraction}
Let \((u,p,z_*)\) be an admissible finite-energy Type I singularity.  Then
there exist scales \(\lambda_k\downarrow0\), pressure gauges \(a_k(t)\), an
ancient suitable weak solution \((U,Q)\) on
\(\R^3\times(-\infty,0)\), and a centered solution \((V,\Pi)\) on
\(\R^3\times\R\), such that, after passing to a subsequence,
\[
   u_k\to U
   \quad\hbox{strongly in }L^q_{\loc}(\R^3\times(-\infty,0)),
   \qquad 1\le q<\infty,
\]
and
\[
   p_k-a_k(t)\rightharpoonup Q
   \quad\hbox{weakly in }L^{3/2}_{\loc}(\R^3\times(-\infty,0)).
\]
The centered velocity \(V\) is smooth, bounded, solves \eqref{p1:eq:centered},
and is a raw extracted normalized Seregin ancient limit with the compact
spacetime \(L^3\) nonvanishing condition \eqref{p1:eq:nonvanishing}.
\end{theorem}

For a raw generated descendant hull
\(\calS=\calS_{\mathrm{raw\text{-}gen}}(u,p,z_*;R_0,\eta_*)\), the class definitions give, for
each \(V\in\calS\), the exhaustive alternative
\[
   V\in
   \calC_{\mathrm{small}}
   \cup
   \calC_{\mathrm{stat},L^3}
   \cup
   \calC_{L^3\mathrm{-tight}}
   \cup
   \calC_{\mathrm{sd}}
   \cup
   \calR(\calS).
\]
This exhaustive alternative is recorded in \cref{p1:prop:exhaustion}.  It is a
direct consequence of the definition of the remaining class, not an
independent Liouville theorem.

\begin{theorem}[Conditional Type I exclusion]
\label{p1:thm:remaining-criterion}
Assume the residual-class hypothesis \Cref{p1:hyp:no-remainder}.  Then no
admissible finite-energy Type I singularity exists.
\end{theorem}

The argument uses one closure theorem and one explicit residual-class
hypothesis.  The uniformly \(L^3\)-tight class is excluded by
\Cref{p2:thm:tight-liouville}.
The remaining class is excluded only under the residual-class hypothesis
\Cref{p1:hyp:no-remainder}.  The structure and decay
class is restricted to hypotheses to which cited Liouville theorems apply.

The resulting implication is
\[
\begin{array}{c}
\text{admissible finite-energy Type I singularity}
\\
\Downarrow
\\
\text{raw extracted nonzero Seregin limit}
\\
\Downarrow
\\
\text{raw generated Seregin state with invariant nonvanishing}
\\
\Downarrow
\\
\text{small/stationary/tight/structure and decay/remaining alternative}
\\
\Downarrow
\\
\text{contradiction by the Liouville-class exclusions and the residual-class hypothesis.}
\end{array}
\]

\begin{itemize}[leftmargin=*,itemsep=2pt]
\item Small amplitude: excluded in \Cref{p1:sec:liouville}.
\item Stationary \(L^3\): excluded by \cite{NRS1996}.
\item Uniformly \(L^3\)-tight: proved in \Cref{p1:thm:tight-liouville}, using
\Cref{p2:thm:tight-liouville}.
\item Structure and decay: routed through
\Cref{p3:thm:typeI-zero-hypotheses} and excluded in
\Cref{p1:thm:known-structure-decay-liouville}.
\item Remaining class: excluded under the residual-class hypothesis
\Cref{p1:hyp:no-remainder}.
\end{itemize}

\subsection{Dependencies and conditional hypothesis}

The proof of the extraction theorem uses the local velocity concentration
result stated as \cref{p1:prop:local-velocity-concentration}; no other result
from the local concentration theory enters the blow-up construction.  The
extraction first produces a raw extracted limit carrying the local \(L^3\)
lower bound \eqref{p1:eq:nonvanishing}.  The subsequent hull construction uses
\Cref{p1:lem:compact-to-invariant-nonzero} and
\Cref{p1:lem:raw-generated-closure-stability}, which provide the invariant
nonvanishing and the raw generated structure.  Mildness is imposed only in the
branch definitions that need it.  The uniformly \(L^3\)-tight Liouville theorem and the
residual-class hypothesis are invoked only after this raw generated class is
formed and the class exhaustion is applied.  The statement
\cref{p1:thm:tight-liouville} follows from the endpoint \(L^3\) closure in
\Cref{p2:thm:tight-liouville}, and \cref{p1:hyp:no-remainder} is the
residual-class hypothesis for
\(\calS_{\mathrm{raw\text{-}gen}}(u,p,z_*;R_0,\eta_*)\).  The
structure and decay class is routed through
\Cref{p3:thm:typeI-zero-hypotheses}, whose hypotheses are recorded in the
structured Liouville section.  Thus the final Type I conclusion is conditional
on the admissibility condition in \Cref{p0:cor:auto-paperI-admissibility}, the
extraction theorem, the raw generated closure statements, the Liouville-class
exclusions, and the residual-class hypothesis.

\subsection{Two entry routes for the reduction}
\label{p1:sec:two-entry-points}

There are two routes to the same raw generated Seregin state space; they have
different hypotheses and first produce the same type of raw extracted ancient
solution.
\begin{enumerate}[label=\textup{(\arabic*)}]
\item \emph{Local pointwise Type I.}  The Part I local concentration and
local endpoint weak-Serrin estimate \Cref{p0:thm:local-weak-serrin-typeI}
produce an admissible local Type I concentration sequence at any singular point
satisfying \eqref{p0:eq:paper0-local-typeI-bound}, and
\Cref{p0:lem:paper0-local-typeI-into-terminal-dichotomy} feeds the resulting
rescalings into the Seregin extraction with positive compact \(L^3\)
velocity mass on a fixed negative-time cylinder.  This is the entry point
used by the final assembly theorem \Cref{p1:thm:final-assembly}.
\item \emph{Global pointwise Type I.}  The Part II extraction estimates,
including the pressure decomposition in \Cref{p1:lem:pressure}, are written
for an admissible Type I solution in the sense of
\Cref{p1:def:admissible-type-I}, which encodes a global pointwise Type I
bound on \(\R^3\).  This streamlines the harmonic-pressure tail estimates,
which would otherwise require the local pressure-tail decomposition recorded
in \Cref{p0:prop:paper0-local-typeI-terminal-compactness}.
\end{enumerate}
In both cases the conclusion is first a raw extracted normalized Seregin
ancient limit satisfying \eqref{p1:eq:nonvanishing}.  The subsequent hull
construction uses the
invariant normalization \eqref{p1:eq:invariant-nonvanishing}.  The local
pointwise entry uses Part I to produce
an admissible concentration sequence and pressure-tail control, while Part II
then runs on the resulting compact cylinders.  No analytic hypothesis beyond
local pointwise Type I and finite Leray energy is added in this reduction.

\begin{center}
\begin{tabular}{p{0.18\textwidth}p{0.30\textwidth}p{0.40\textwidth}}
Global route & global pointwise Type I &
global no-escape \(\Rightarrow\) finite-energy Seregin extraction \\
Local route & local pointwise Type I &
AB endpoint weak-Serrin bounds \(\Rightarrow\) terminal no-escape
\(\Rightarrow\) local Seregin extraction
\end{tabular}
\end{center}

\subsection{Organization}

\Cref{p1:sec:prelim} introduces suitable solutions, scale-invariant quantities,
the admissibility condition, scaling, pressure gauges, and local concentration.
\Cref{p1:sec:blowup-compactness} proves the compactness theorem for Type I
rescalings away from the terminal time and constructs the ancient suitable
weak limit.  \Cref{p1:sec:nontriviality} proves persistence of concentration and
nontriviality.  \Cref{p1:sec:seregin} derives the centered equation and defines
raw extracted limits, the raw generated state space, and its mild stratum.
\Cref{p1:sec:classes} defines the
Liouville classes and the remaining class.  \Cref{p1:sec:liouville} proves the
class exclusions and the uniformly tight theorem.  \Cref{p1:sec:criterion} proves
the conditional remaining-case criterion of
\Cref{p1:thm:remaining-criterion} and the residual closure hypothesis.
Appendix \ref{p1:app:centered-symmetries} gives
the elementary symmetries of the centered equation.

\section{Suitable solutions, scale-invariant quantities, and local concentration}
\label{p1:sec:prelim}

For \(z_0=(x_0,t_0)\), write
\[
   Q_r(z_0)=B_r(x_0)\times(t_0-r^2,t_0).
\]
When \(z_0=(0,0)\), write \(Q_r=Q_r(0,0)\).  The spatial average of \(f\) over
\(B_r(x_0)\) is
\[
   (f)_{B_r(x_0)}(t)=\frac1{|B_r|}\int_{B_r(x_0)}f(x,t)\,dx.
\]
Throughout this part, \(\mathbb P\) denotes the Helmholtz--Leray projection
onto divergence-free vector fields on \(\R^3\).

\begin{definition}[Terminal singular point]\label{p1:def:terminal-singular}
Let \((u,p)\) be a suitable weak solution on \(\R^3\times(0,T)\).  A terminal
point \(z_*=(x_*,T)\) is regular if there are \(r>0\) and \(K<\infty\) such
that
\[
   \esssup_{Q_r(z_*)}|u|\le K.
\]
It is singular otherwise.
\end{definition}

\begin{definition}[Suitable weak solution]\label{p1:def:suitable}
This definition repeats the Part I notion in the notation used here; the two
definitions agree.
Let \(\Omega\subset\R^3\times\R\) be open.  A pair \((u,p)\) is a suitable
weak solution of \eqref{p1:eq:NS} in \(\Omega\) if
\[
   u\in L^\infty_tL^2_x(\Omega')\cap L^2_t\dot H^1_x(\Omega'),
   \qquad
   p\in L^{3/2}(\Omega')
\]
for every compact cylinder \(\Omega'\Subset\Omega\), if \((u,p)\) solves
\eqref{p1:eq:NS} distributionally, \(\divergence u=0\), and if the local energy
inequality
\begin{align}\label{p1:eq:lei}
   &\int_{\R^3}|u(x,t_2)|^2\phi(x,t_2)\,dx
   +2\int_{t_1}^{t_2}\int_{\R^3}|\nabla u|^2\phi\,dx\,dt  \notag\\
   &\le
   \int_{\R^3}|u(x,t_1)|^2\phi(x,t_1)\,dx
   +\int_{t_1}^{t_2}\int_{\R^3}|u|^2(\partial_t\phi+\Delta\phi)\,dx\,dt
   \notag\\
   &\quad
   +\int_{t_1}^{t_2}\int_{\R^3}(|u|^2+2p)u\cdot\nabla\phi\,dx\,dt
\end{align}
holds for every non-negative \(\phi\in C_c^\infty(\Omega)\) and almost every
\(t_1<t_2\) for which the displayed terms are defined.
\end{definition}

This is the Caffarelli--Kohn--Nirenberg notion of suitability
\cite{CKN1982}; the global finite-energy background is Leray's
Leray--Hopf class \cite{Leray1934}.

\begin{definition}[Suitable Leray--Hopf solution on \(\R^3\times(0,T)\)]
\label{p1:def:suitable-leray-hopf}
A suitable Leray--Hopf solution on \(\R^3\times(0,T)\) is a suitable weak
solution on \(\R^3\times(0,T)\) satisfying
\[
   u\in L^\infty(0,T;L^2(\R^3))
   \cap L^2(0,T;\dot H^1(\R^3)),
   \qquad
   u\in C_w([0,T);L^2(\R^3)),
\]
and the global Leray energy inequality.  The pressure is taken in
\(L^{3/2}_{\loc}(\R^3\times(0,T))\), modulo functions of time.
\end{definition}

Throughout the global extraction and final assembly, the phrase
\emph{finite-energy suitable weak solution on \(\R^3\times(0,T)\)} means a
suitable Leray--Hopf solution in the sense of
\Cref{p1:def:suitable-leray-hopf}.  Purely local statements use the phrase
``suitable weak solution in a cylinder'' and do not invoke the global energy
or whole-space pressure representative unless explicitly stated.

Define
\begin{equation}\label{p1:eq:C-D-def}
   A(u;z_0,r)=r^{-1}\esssup_{t_0-r^2<t<t_0}
   \int_{B_r(x_0)}|u(x,t)|^2\,dx,
\end{equation}
\begin{equation}\label{p1:eq:C-D-def-2}
   C(u;z_0,r)=r^{-2}\iint_{Q_r(z_0)}|u|^3\,dx\,dt,
   \qquad
   D(p;z_0,r)=r^{-2}\iint_{Q_r(z_0)}
   |p-(p)_{B_r(x_0)}(t)|^{3/2}\,dx\,dt,
\end{equation}
\begin{equation}\label{p1:eq:E-I-def}
   E(u;z_0,r)=r^{-1}\iint_{Q_r(z_0)}|\nabla u|^2\,dx\,dt,
\end{equation}
\begin{definition}[Admissible finite-energy Type I singularity]
\label{p1:def:admissible-type-I}
A triple \((u,p,z_*)\), with \(z_*=(x_*,T)\), is an admissible finite-energy
Type I singularity if:
\begin{enumerate}[label=(\roman*)]
\item \((u,p)\) is a suitable Leray--Hopf solution on
\(\R^3\times(0,T)\);
\item \(z_*=(x_*,T)\) is a singular point;
\item the Type I bound \eqref{p1:eq:type-I-bound} holds;
\item there are scales \(\lambda_k\downarrow0\) and a constant
\(\eta_v>0\) such that
\[
   C(u;z_*,\lambda_k)\ge\eta_v
   \qquad\hbox{for every }k;
\]
\item for the corresponding rescaled velocities
\[
   u_k(x,t)=\lambda_k u(x_*+\lambda_kx,T+\lambda_k^2t),
\]
the velocity no-escape condition holds:
\begin{equation}\label{p1:eq:velocity-no-escape}
   \lim_{\sigma\downarrow0}
   \sup_k
   \int_{-\sigma}^{0}\int_{B_1}|u_k(x,t)|^3\,dx\,dt=0 .
\end{equation}
\end{enumerate}
The sequence \(\{\lambda_k\}\) is called an admissible concentration sequence
at \(z_*\).
\end{definition}

\begin{remark}
Item (iv) is the velocity-only concentration supplied by the local
concentration theorem; see \cref{p1:prop:local-velocity-concentration}.  Thus,
for a singular point, the additional admissibility requirement is the
no-escape condition \eqref{p1:eq:velocity-no-escape}.  It is the sole
terminal-slice compactness assumption used to prove nontriviality of the
ancient limit.  No uniform bound on \(A+C+D+E\) over all terminal scales is
assumed.
\end{remark}

\begin{lemma}[Uniform kinetic energy implies no escape]
\label{p1:lem:A-seq-no-escape}
Let \((u,p)\) satisfy the Type I bound \eqref{p1:eq:type-I-bound}, and let
\(\lambda_k\downarrow0\).  Suppose that
\[
   \sup_k A(u;z_*,\lambda_k)\le B<\infty .
\]
Then the corresponding rescaled velocities satisfy the no-escape condition
\eqref{p1:eq:velocity-no-escape}.
\end{lemma}

\begin{proof}
Choose \(M<\infty\) and \(t_0<T\) such that
\[
   \|u(t)\|_{L^\infty(\R^3)}
   \le M(T-t)^{-1/2},
   \qquad t_0<t<T.
\]
For \(k\) sufficiently large, \(T+\lambda_k^2t>t_0\) whenever
\(-1<t<0\).  Hence
\[
   \|u_k(t)\|_{L^\infty(B_1)}
   \le M(-t)^{-1/2},\qquad -1<t<0.
\]
By scale invariance of \(A\),
\[
   \esssup_{-1<t<0}\int_{B_1}|u_k(x,t)|^2\,dx
   =
   A(u;z_*,\lambda_k)\le B.
\]
Therefore, for \(0<\sigma<1\),
\[
\begin{aligned}
   \int_{-\sigma}^{0}\int_{B_1}|u_k|^3\,dx\,dt
   &\le
   \int_{-\sigma}^{0}
   \|u_k(t)\|_{L^\infty(B_1)}
   \int_{B_1}|u_k(x,t)|^2\,dx\,dt  \\
   &\le MB\int_{-\sigma}^{0}(-t)^{-1/2}\,dt
    =2MB\sigma^{1/2}.
\end{aligned}
\]
Taking the supremum over the tail \(k\ge k_0\) and then letting
\(\sigma\downarrow0\) gives \eqref{p1:eq:velocity-no-escape} for the tail of the
sequence.  For each of the finitely many remaining indices
\(1\le k<k_0\), the rescaled velocity is defined on
\(B_1\times(-\delta_k,0)\) for some \(\delta_k>0\), and belongs to
\(L^3\) there.  Therefore absolute continuity of the \(L^3\)-integral gives
\[
   \lim_{\sigma\downarrow0}
   \max_{1\le k<k_0}
   \int_{-\sigma}^{0}\int_{B_1}|u_k|^3\,dx\,dt=0.
\]
Combining the finite initial segment with the tail proves the full supremum
in \eqref{p1:eq:velocity-no-escape}.
\end{proof}

\begin{proposition}[Local velocity concentration]
\label{p1:prop:local-velocity-concentration}
Let \((u,p)\) be suitable in a parabolic neighborhood of
\(z_*=(x_*,T)\).  If \(z_*\) is singular, then there exist
\(\lambda_k\downarrow0\) and \(\eta_v>0\) such that
\[
   C(u;z_*,\lambda_k)\ge\eta_v
   \qquad\hbox{for every }k.
\]
\end{proposition}

\begin{proof}
This is the velocity-only concentration corollary of the local concentration
result in \Cref{p0:thm:typeI-typeII-dichotomy}.  That result follows from the
velocity \(\varepsilon\)-regularity criterion: if
\[
   \limsup_{r\downarrow0}C(u;z_*,r)=0,
\]
then \(u\) is bounded in a backward parabolic cylinder ending at \(z_*\),
contradicting singularity of \(z_*\).
\end{proof}

\begin{lemma}[Scaling]\label{p1:lem:scaling}
Let \((u,p)\) be suitable in a neighborhood of \((x_*,T)\).  For
\(\lambda>0\), set
\[
   u^\lambda(x,t)=\lambda u(x_*+\lambda x,T+\lambda^2t),
   \qquad
   p^\lambda(x,t)=\lambda^2p(x_*+\lambda x,T+\lambda^2t).
\]
Then \((u^\lambda,p^\lambda)\) is suitable on the rescaled domain.  Moreover
\(A,C,D,E\) are invariant under this scaling.
\end{lemma}

\begin{proof}
Suitability of the rescaled pair and preservation of pressure gauges follow
by applying \eqref{p1:eq:lei} to the test function
\[
   \phi^\lambda(X,S)
   =
   \phi\left(\frac{X-x_*}{\lambda},\frac{S-T}{\lambda^2}\right)
\]
and then changing variables
\((X,S)=(x_*+\lambda x,T+\lambda^2t)\).  The differential identities
\[
   \partial_tu^\lambda=\lambda^3(\partial_tu)(X,S),\quad
   (u^\lambda\cdot\nabla)u^\lambda
      =\lambda^3((u\cdot\nabla)u)(X,S),
\]
\[
   \nabla p^\lambda=\lambda^3(\nabla p)(X,S),\qquad
   \Delta u^\lambda=\lambda^3(\Delta u)(X,S)
\]
give the rescaled equation and the divergence condition.

For the pressure average,
\[
   (p^\lambda)_{B_r}(t)
   =
   \lambda^2(p)_{B_{\lambda r}(x_*)}(S).
\]
Consequently
\[
   \iint_{Q_r}|u^\lambda|^3\,dx\,dt
   =
   \lambda^{-2}\iint_{Q_{\lambda r}(z_*)}|u|^3\,dX\,dS
\]
and
\[
   \iint_{Q_r}|p^\lambda-(p^\lambda)_{B_r}(t)|^{3/2}\,dx\,dt
   =
   \lambda^{-2}
   \iint_{Q_{\lambda r}(z_*)}
   |p-(p)_{B_{\lambda r}(x_*)}(S)|^{3/2}\,dX\,dS .
\]
Multiplication by \(r^{-2}\) therefore gives invariance of \(C\) and \(D\).
The remaining two quantities have the same critical dimensional balance:
\begin{align*}
   \int_{B_r}|u^\lambda(x,t)|^2\,dx
   &=
   \lambda^{-1}\int_{B_{\lambda r}(x_*)}|u(X,S)|^2\,dX,\\
   \iint_{Q_r}|\nabla u^\lambda|^2\,dx\,dt
   &=
   \lambda^{-1}\iint_{Q_{\lambda r}(z_*)}|\nabla u|^2\,dX\,dS .
\end{align*}
Multiplication by \(r^{-1}\) on the rescaled cylinder is the same as
multiplication by \((\lambda r)^{-1}\) on the original cylinder.  Hence
\(A\) and \(E\) are invariant as well.
\end{proof}

\begin{lemma}[Pressure gauges do not change the equations]
\label{p1:lem:gauges}
Let \((v,q)\) be suitable in a cylinder \(Q\), and let
\(\alpha(t)\in L^{3/2}_{\loc}\) depend only on time.  Then
\((v,q-\alpha(t))\) satisfies the same distributional Navier--Stokes equation,
the same divergence condition, and the same local energy inequality.
\end{lemma}

\begin{proof}
In the distributional equation, replacing \(q\) by \(q-\alpha(t)\) adds
\(-\nabla\alpha(t)=0\).  In the local energy inequality the only new term is
\[
   -2\int \alpha(t)v(x,t)\cdot\nabla\phi(x,t)\,dx\,dt .
\]
The product is integrable on the support of \(\phi\): by suitability,
\(v\in L^\infty_tL^2_x\) locally, and
\[
   \left|\int v(\cdot,t)\cdot\nabla\phi(\cdot,t)\,dx\right|
   \le
   \|v(t)\|_{L^2(\supp\phi(t))}
   \|\nabla\phi(t)\|_{L^2}.
\]
For almost every \(t\), \(\divergence v(\cdot,t)=0\) in distributions and
\(\phi(\cdot,t)\) is compactly supported, so
\[
   \int v(x,t)\cdot\nabla\phi(x,t)\,dx=0 .
\]
Thus the additional term vanishes.
\end{proof}

\section{The Type I blow-up sequence and compactness}
\label{p1:sec:blowup-compactness}

\subsection{The rescaled sequence}

Choose \(M<\infty\) so that, after decreasing \(T_0<T\) if necessary,
\begin{equation}\label{p1:eq:type-I-pointwise}
   \|u(t)\|_{L^\infty(\R^3)}
   \le \frac{M}{\sqrt{T-t}},
   \qquad T_0<t<T.
\end{equation}
For \(\lambda_k\downarrow0\), define
\begin{equation}\label{p1:eq:rescaled-sequence}
   u_k(x,t)=\lambda_k u(x_*+\lambda_k x,T+\lambda_k^2t),
   \qquad
   p_k(x,t)=\lambda_k^2p(x_*+\lambda_k x,T+\lambda_k^2t).
\end{equation}
Then \(u_k\) is defined on \(\R^3\times(-T/\lambda_k^2,0)\), and these
domains exhaust \(\R^3\times(-\infty,0)\).

\subsection{Inherited Type I bounds}

\begin{lemma}[Inherited Type I bound]\label{p1:lem:inherited-bound}
For every compact interval \(I\Subset(-\infty,0)\),
\[
   \sup_k\|u_k\|_{L^\infty(\R^3\times I)}<\infty.
\]
More precisely, for every fixed \(t<0\) and all sufficiently large \(k\),
\[
   \|u_k(t)\|_{L^\infty(\R^3)}
   \le \frac{M}{\sqrt{-t}}.
\]
\end{lemma}

\begin{proof}
Since \(T-(T+\lambda_k^2t)=\lambda_k^2(-t)\), \eqref{p1:eq:type-I-pointwise}
gives, for all large \(k\),
\[
   \|u_k(t)\|_\infty
   =
   \lambda_k\|u(T+\lambda_k^2t)\|_\infty
   \le
   \lambda_k\frac{M}{\sqrt{\lambda_k^2(-t)}}
   =
   \frac{M}{\sqrt{-t}}.
\]
Taking the supremum over \(t\in I\) gives the interval statement.
\end{proof}

\subsection{Local energy, pressure, and time-derivative estimates}

\begin{lemma}[Pressure gauges and local pressure bounds]\label{p1:lem:pressure}
After fixing the whole-space Leray pressure representative below, set
\[
   a_k(t)=(p_k)_{B_1}(t)
\]
whenever the rescaled solution is defined.  Then for every
\[
   Q=B_R\times[-S,-\sigma]\Subset\R^3\times(-\infty,0)
\]
one has
\[
   \sup_k\|p_k-a_k(t)\|_{L^{3/2}(Q)}<\infty .
\]
The estimates below use the global pointwise Type I bound encoded in the
admissible Type I assumption \Cref{p1:def:admissible-type-I} to control
\(L^2\)-mass on every annulus uniformly in \(k\); the corresponding
local-Type-I version replaces the outer annular control by the Leray
finite-energy tail estimate already recorded in
\Cref{p0:prop:paper0-local-typeI-terminal-compactness}.
\end{lemma}

\begin{proof}
For large \(k\), the original times \(T+\lambda_k^2t\), \(t\in[-S,-\sigma]\),
lie in the interval where \eqref{p1:eq:type-I-pointwise} holds.  At the
point where the whole-space Riesz pressure representation is used, the
ambient solution is a suitable Leray--Hopf solution on
\(\R^3\times(0,T)\) in the sense of
\Cref{p1:def:suitable-leray-hopf}; the pressure representative is
therefore chosen, modulo functions of time, as the whole-space Leray
pressure
\[
   p=\mathcal R_i\mathcal R_j(u_i u_j) .
\]
On those slices
\(u(t)\in L^2(\R^3)\cap L^\infty(\R^3)\), hence
\[
   u_i(t)u_j(t)\in L^1(\R^3)\cap L^{3/2}(\R^3).
\]
Use the whole-space Leray pressure representative
\[
   p^L=\mathcal R_i\mathcal R_j(u_i u_j),
\]
modulo functions of time.  The standard convention is
\[
   \widehat{\mathcal R_i f}(\xi)=i\xi_i|\xi|^{-1}\widehat f(\xi),
\]
so that \(\mathcal R_i\mathcal R_j\) has multiplier
\(-\xi_i\xi_j|\xi|^{-2}\) and
\(-\Delta(\mathcal R_i\mathcal R_j f)=\partial_i\partial_jf\).  Thus this is
the pressure reconstruction from
\(-\Delta p=\partial_i\partial_j(u_i u_j)\) and the Calderon--Zygmund theory
of Riesz transforms \cite{Stein1970,Sohr2001}.  Since
\(u(t)\in L^2(\R^3)\cap L^\infty(\R^3)\), interpolation gives
\(u(t)\in L^3(\R^3)\), hence \(u_i(t)u_j(t)\in L^{3/2}(\R^3)\) and
\(p^L(t)\in L^{3/2}(\R^3)\) for a.e. \(t\).

The given pressure and \(p^L\) differ by a function of time.  Indeed, applying
the Helmholtz projection to the whole-space weak momentum equation gives
\[
   \partial_tu-\Delta u+\mathbb P\divergence(u\otimes u)=0
\]
in solenoidal distributions.  The complementary gradient part of
\(\divergence(u\otimes u)\) is precisely
\(-\nabla p^L\) with the convention above.  Therefore the same velocity field
solves the distributional equation with \(p^L\) in place of \(p\).  Subtracting
the two formulations gives
\[
   \nabla(p-p^L)=0
   \quad\hbox{in }\mathcal D'(\R^3)
\]
for a.e. time.  By the de Rham theorem on \(\R^3\), \(p-p^L\) is spatially
constant for a.e. time.  Replacing the pressure by \(p^L\) is therefore a
time-dependent gauge change and does not affect suitability by
\cref{p1:lem:gauges}.

Fix \(\chi_R\in C_c^\infty(B_{8R})\) with \(\chi_R=1\) on \(B_{4R}\).  For
\(t\in[-S,-\sigma]\), choose the representative so that, in \(B_{2R}\),
\[
   p_k=q_k+h_k,\qquad
   q_k=\mathcal R_i\mathcal R_j(\chi_Ru_{k,i}u_{k,j}),
\]
and \(h_k(\cdot,t)\) is harmonic in \(B_{4R}\) for a.e. \(t\).  This choice
is made before \(a_k(t)=(p_k)_{B_1}(t)\) is evaluated.  The harmonicity holds
because \((1-\chi_R)u_{k,i}u_{k,j}\) is supported outside \(B_{4R}\), so the
localized Riesz potential generated by this exterior part has zero Laplacian
in \(B_{4R}\).  Calderon--Zygmund boundedness gives, for a.e. \(t\),
\[
   \|q_k(t)\|_{L^{3/2}(B_{2R})}
   \le C\|u_k(t)\|_{L^3(B_{8R})}^2
   \le C(R,\sigma,M).
\]
Integrating over \([-S,-\sigma]\) gives a uniform
\(L^{3/2}_{x,t}\)-bound for \(q_k\).
Let \(c_{k,R}(t)=h_k(0,t)\), where \(h_k(\cdot,t)\) is represented by its
smooth harmonic representative.  For \(x\in B_R\), the kernel estimate
\[
   |K_{ij}(x-z)-K_{ij}(-z)|\le C R|z|^{-4},\qquad |z|>4R,
\]
implies
\[
   |h_k(x,t)-c_{k,R}(t)|
   \le C R\sum_{m=2}^\infty(2^mR)^{-4}
   \int_{2^mR<|z|<2^{m+1}R}|u_k(z,t)|^2\,dz .
\]
Using \cref{p1:lem:inherited-bound},
\[
   \int_{2^mR<|z|<2^{m+1}R}|u_k(z,t)|^2\,dz
   \le C M^2\sigma^{-1}(2^mR)^3,
\]
and therefore
\[
   |h_k(x,t)-c_{k,R}(t)|
   \le C M^2\sigma^{-1}\sum_{m=2}^{\infty}2^{-m},
   \qquad x\in B_R,
\]
which implies
\[
   \|h_k(t)-c_{k,R}(t)\|_{L^\infty(B_R)}\le C(M,\sigma).
\]
Thus
\[
   \sup_k\|p_k-c_{k,R}(t)\|_{L^{3/2}(B_R\times[-S,-\sigma])}<\infty .
\]
If \(R<1\), enlarge \(R\) to \(1\).  Otherwise \(B_1\subset B_R\).  Jensen's
inequality gives, for a.e. \(t\),
\[
   |a_k(t)-c_{k,R}(t)|^{3/2}
   \le
   |B_1|^{-1}\int_{B_1}|p_k-c_{k,R}(t)|^{3/2}\,dx .
\]
Integrating in time and adding constants on \(B_R\) gives the desired bound
for \(p_k-a_k(t)\).
\end{proof}

\begin{lemma}[Compatibility of pressure gauges]\label{p1:lem:gauge-compatibility}
Let \(B_1\subset B_R\), \(I\Subset(-\infty,0)\), and suppose
\[
   p_k-c_k(t)
   \quad\hbox{is bounded in }L^{3/2}(B_R\times I)
\]
for some time-dependent gauges \(c_k(t)\).  If
\(a_k(t)=(p_k)_{B_1}(t)\), then \(p_k-a_k(t)\) is bounded in
\(L^{3/2}(B_{R'}\times I)\) for every \(R'\le R\).  Changing between such
gauges changes the pressure only by a function of time and therefore does not
affect the equation, the local energy inequality, or the centered equation.
\end{lemma}

\begin{proof}
For a.e. \(t\), Jensen's inequality gives
\[
   |a_k(t)-c_k(t)|^{3/2}
   \le |B_1|^{-1}\int_{B_1}|p_k-c_k(t)|^{3/2}\,dx .
\]
Integrating in \(t\) and using \(B_{R'}\subset B_R\) yields
\[
   \|a_k-c_k\|_{L^{3/2}(I)}^{3/2}
   \le
   |B_1|^{-1}\|p_k-c_k(t)\|_{L^{3/2}(B_R\times I)}^{3/2}.
\]
Thus
\[
   \|p_k-a_k(t)\|_{L^{3/2}(B_{R'}\times I)}
   \le
   \|p_k-c_k(t)\|_{L^{3/2}(B_{R'}\times I)}
   +|B_{R'}|^{2/3}\|a_k-c_k\|_{L^{3/2}(I)}
\]
and the right-hand side is bounded by
\[
   C(R')\|p_k-c_k(t)\|_{L^{3/2}(B_R\times I)} .
\]
This proves the asserted bound.
The last assertion is \cref{p1:lem:gauges}, and centering transforms the
time-dependent gauge into another time-dependent gauge.
\end{proof}

\begin{lemma}[Pressure gauge atlas]\label{p1:lem:pressure-atlas}
Let \((u_n,p_n)\) be a rescaled sequence of suitable weak solutions satisfying
the local compactness bounds on every compact cylinder
\(K\Subset\R^3\times(-\infty,0)\).  For every ball \(B_R\) and compact time
interval \(I\Subset(-\infty,0)\) there are time-dependent functions
\(a_{n,R}(t)\) such that
\[
   p_n-a_{n,R}(t)
   \quad\hbox{is bounded in }L^{3/2}(B_R\times I).
\]
On nested balls, any two such gauges for the same \(n\) differ by functions of
time only.  After passing to a diagonal subsequence over \(R\) and \(I\), the
gauge-fixed pressures converge weakly in \(L^{3/2}_{\loc}\) to a pressure
\(Q\) defined modulo functions of time.  These gauge changes do not affect the
distributional Navier--Stokes equations or the local energy inequality, and
the same compatibility passes to pressure limits.
\end{lemma}

\begin{proof}
The local pressure bounds are those proved in \Cref{p1:lem:pressure} for the
global Type I extraction and in
\Cref{p0:prop:paper0-local-typeI-terminal-compactness} for the local Type I
rescaling construction.  On each fixed \(B_R\times I\), choose any of the gauges
provided there, for instance a spatial average on a unit ball contained in
\(B_R\) after enlarging \(R\) if necessary.  Compatibility on nested balls is
\Cref{p1:lem:gauge-compatibility}: subtracting two admissible representatives
leaves the same pressure gradient, hence the difference is spatially constant
for a.e. time.  A countable exhaustion by balls and compact time intervals,
together with Banach--Alaoglu in \(L^{3/2}\), gives a diagonal subsequence and
a compatible local pressure limit.  Invariance of the equation and local
energy inequality under time-only pressure changes is \Cref{p1:lem:gauges}.
\end{proof}

\begin{remark}[Pressure-gauge atlas]
\label{p1:rem:pressure-atlas}
In the sequel, a pressure limit \(Q\) means a compatible family of local
representatives modulo time-dependent gauges.  On overlaps, the gradients
agree; hence the representatives differ only by functions of time.  This is
the pressure atlas of \Cref{p1:lem:pressure-atlas} for local suitable limits.
When local
compactness is run on a fixed compact cylinder \(K\), we use the gauge
\(a_{k,K}(t)=(p_k)_{B_1}(t)\) supplied by \Cref{p1:lem:pressure}, and
\Cref{p1:lem:gauge-compatibility} identifies a single pressure limit on
overlaps modulo functions of time.
\end{remark}

\begin{lemma}[Local energy and dissipation]\label{p1:lem:energy-diss}
Let
\[
   Q=B_R\times[-S,-\sigma]\Subset\R^3\times(-\infty,0).
\]
Then
\[
   \sup_k\|u_k\|_{L^\infty_tL^2_x(Q)}
   +\sup_k\|\nabla u_k\|_{L^2(Q)}<\infty .
\]
\end{lemma}

\begin{proof}
The \(L^\infty_tL^2_x\) estimate follows from
\[
   \|u_k(t)\|_{L^2(B_R)}
   \le |B_R|^{1/2}M\sigma^{-1/2},
   \qquad t\in[-S,-\sigma].
\]
Choose \(\phi\in C_c^\infty(B_{2R}\times[-2S,-\sigma/2])\), \(0\le\phi\le1\),
with \(\phi=1\) on \(Q\).  Apply the local energy inequality to
\((u_k,p_k-a_k(t))\), using \cref{p1:lem:pressure} on the larger cylinder.  For
a.e. admissible times and then by monotone approximation of temporal cutoffs,
\[
   2\iint |\nabla u_k|^2\phi
   \le
   \iint |u_k|^2|\partial_t\phi+\Delta\phi|
   +\iint (|u_k|^2+2|p_k-a_k(t)|)|u_k||\nabla\phi|.
\]
The first term is bounded because \(\partial_t\phi+\Delta\phi\) is bounded
and \(u_k\) is uniformly bounded in \(L^\infty\) on the finite-measure support
of \(\phi\).  The cubic velocity term satisfies
\[
   \iint |u_k|^3|\nabla\phi|
   \le
   \|\nabla\phi\|_\infty
   |\supp\nabla\phi|\,\|u_k\|_{L^\infty(\supp\nabla\phi)}^3
   \le C(Q,\phi,M,\sigma).
\]
For the pressure term, H\"older's inequality on \(\supp\nabla\phi\) gives
\[
   \iint |p_k-a_k(t)|\,|u_k|\,|\nabla\phi|
   \le
   \|p_k-a_k(t)\|_{L^{3/2}}
   \|u_k\nabla\phi\|_{L^3},
\]
and both factors are uniformly bounded by \cref{p1:lem:pressure} and
\cref{p1:lem:inherited-bound}.  Since \(\phi=1\) on \(Q\), the dissipation bound
follows.
\end{proof}

\begin{lemma}[Time derivative bound]\label{p1:lem:time-derivative}
For every compact cylinder \(Q\Subset\R^3\times(-\infty,0)\),
\[
   \partial_tu_k
   \quad\hbox{is bounded in }L^{3/2}_tW^{-1,3/2}_x(Q).
\]
\end{lemma}

\begin{proof}
The equation is
\[
   \partial_tu_k
   =
   \Delta u_k-\divergence(u_k\otimes u_k)-\nabla(p_k-a_k(t)).
\]
Let \(Q=B_R\times[-S,-\sigma]\).  Throughout the proof we use the
embedding \(W^{1,3}_0(B_R)\hookrightarrow H^1_0(B_R)\) on the bounded
domain \(B_R\); duality then gives
\(H^{-1}(B_R)\hookrightarrow W^{-1,3/2}(B_R)\), so that any
\(L^2_tH^{-1}_x(Q)\)-bound becomes an \(L^{3/2}_tW^{-1,3/2}_x(Q)\)-bound
by H\"older in time on the finite interval \([-S,-\sigma]\).  The term
\(\Delta u_k\) is bounded in
\(L^2_tH^{-1}_x(Q)\) by \cref{p1:lem:energy-diss}; since
\(W^{1,3}_0(B_R)\subset H^1_0(B_R)\) and the time interval is finite, this is
also an \(L^{3/2}_tW^{-1,3/2}_x\)-bound.  Explicitly, for
\(\zeta\in W^{1,3}_0(B_R)\),
\[
   |\langle\Delta u_k,\zeta\rangle|
   =
   \left|\int_{B_R}\nabla u_k:\nabla\zeta\,dx\right|
   \le
   C(R)\|\nabla u_k\|_{L^2(B_R)}\|\zeta\|_{W^{1,3}(B_R)}.
\]
The \(L^2_t\)-bound on \(\nabla u_k\), together with the finite length of
\([-S,-\sigma]\), gives the claimed \(L^{3/2}_t\)-bound.

For the nonlinear term, for
\(\zeta\in W^{1,3}_0(B_R)\),
\[
   \left|\int_{B_R}(u_k\otimes u_k):\nabla\zeta\,dx\right|
   \le
   \|u_k\otimes u_k\|_{L^{3/2}(B_R)}
   \|\nabla\zeta\|_{L^3(B_R)}.
\]
Moreover
\[
   \|u_k\otimes u_k\|_{L^{3/2}(Q)}
   \le
   \|u_k\|_{L^\infty(Q)}\|u_k\|_{L^{3/2}(Q)}
   \le C(Q).
\]
Thus \(\divergence(u_k\otimes u_k)\) is bounded in
\(L^{3/2}_tW^{-1,3/2}_x(Q)\).
For the pressure term,
\[
   \left|\int_{B_R}(p_k-a_k(t))\,\divergence\zeta\,dx\right|
   \le
   \|p_k-a_k(t)\|_{L^{3/2}(B_R)}
   \|\nabla\zeta\|_{L^3(B_R)}.
\]
\Cref{p1:lem:pressure} gives
\[
   p_k-a_k(t)\quad\hbox{bounded in }L^{3/2}(Q).
\]
Hence \(\nabla(p_k-a_k(t))\) is bounded in
\(L^{3/2}_tW^{-1,3/2}_x(Q)\).  Summing the three terms gives the result.
\end{proof}

\subsection{Compactness away from the terminal time}

\begin{proposition}[Compactness]\label{p1:prop:compactness}
After passing to a subsequence,
\[
   u_k\to U
\]
strongly in \(L^q_{\loc}(\R^3\times(-\infty,0))\) for every finite \(q\), and
\[
   p_k-a_k(t)\rightharpoonup Q
   \quad\hbox{weakly in }L^{3/2}_{\loc}(\R^3\times(-\infty,0)).
\]
\end{proposition}

\begin{proof}
Fix \(Q_0\Subset\R^3\times(-\infty,0)\).  By
\cref{p1:lem:energy-diss}, \(u_k\) is bounded in \(L^2_tH^1_x(Q_0)\), and by
\cref{p1:lem:time-derivative}, \(\partial_tu_k\) is bounded in
\(L^{3/2}_tW^{-1,3/2}_x(Q_0)\).  On bounded balls,
\[
   H^1\Subset L^2\hookrightarrow W^{-1,3/2}.
\]
The Aubin--Lions--Simon theorem \cite{Aubin1963,LionsMagenes,Simon1987} gives
relative compactness in \(L^2(Q_0)\).  Hence, after passing to a subsequence,
\(u_k\to U\) strongly in \(L^2(Q_0)\) and a.e. on \(Q_0\).

The inherited Type I bound gives a common \(L^\infty(Q_0)\) bound, and the
limit belongs to the same \(L^\infty\) ball by a.e. convergence.  For
\(q\ge2\),
\[
   \|u_k-U\|_{L^q(Q_0)}^q
   \le
   (2\sup_j\|u_j\|_{L^\infty(Q_0)}+2\|U\|_{L^\infty(Q_0)})^{q-2}
   \|u_k-U\|_{L^2(Q_0)}^2.
\]
For \(1\le q<2\), H\"older's inequality on the finite-measure set \(Q_0\)
gives
\[
   \|u_k-U\|_{L^q(Q_0)}
   \le |Q_0|^{1/q-1/2}\|u_k-U\|_{L^2(Q_0)}\to0.
\]
Choose an exhaustion \(Q^{(m)}\Subset\R^3\times(-\infty,0)\) by compact
cylinders.  The preceding argument gives a subsequence converging on
\(Q^{(1)}\); extracting successively on \(Q^{(m)}\) and taking the diagonal
subsequence gives strong local convergence on every compact cylinder.

For pressure, \cref{p1:lem:pressure} gives boundedness of
\(p_k-a_k(t)\) in \(L^{3/2}(Q^{(m)})\) for every \(m\).  Since
\(L^{3/2}\) is reflexive, each restriction has a weakly convergent
subsequence.  The same diagonal extraction gives a single limit \(Q\) in
\(L^{3/2}_{\loc}\), because the weak limits agree on overlaps by uniqueness
of distributional limits.
\end{proof}

\subsection{The ancient suitable weak limit}

\begin{lemma}[Limit equation]\label{p1:lem:limit-equation}
The limit \((U,Q)\) solves
\[
   \partial_tU+(U\cdot\nabla)U+\nabla Q=\Delta U,\qquad
   \divergence U=0
\]
in distributions on \(\R^3\times(-\infty,0)\).
\end{lemma}

\begin{proof}
For \(\varphi\in C_c^\infty(\R^3\times(-\infty,0);\R^3)\), the weak
formulation for \((u_k,p_k-a_k(t))\) is
\[
   \iint\left[
      -u_k\cdot\partial_t\varphi
      -(u_k\otimes u_k):\nabla\varphi
      +(p_k-a_k(t))\divergence\varphi
      +\nabla u_k:\nabla\varphi
   \right]\,dx\,dt=0.
\]
The linear time term passes by strong \(L^2_{\loc}\)-convergence of \(u_k\).
The pressure term passes by weak \(L^{3/2}_{\loc}\)-convergence of
\(p_k-a_k(t)\), since \(\divergence\varphi\in L^3\).  The viscous term passes
by weak \(L^2_{\loc}\)-convergence of \(\nabla u_k\), which follows from the
uniform local dissipation bound.  The nonlinear term
passes because
\[
   \|u_k\otimes u_k-U\otimes U\|_{L^1(\supp\varphi)}
   \le
   \|u_k-U\|_{L^2}
   \bigl(\|u_k\|_{L^2}+\|U\|_{L^2}\bigr)\to0.
\]
Testing with gradients of scalar functions gives \(\divergence U=0\).
\end{proof}

\begin{lemma}[Stability of suitability]\label{p1:lem:suitability}
The pair \((U,Q)\) is a suitable weak solution on
\(\R^3\times(-\infty,0)\).
\end{lemma}

\begin{proof}
Use the distributional form of the local energy inequality.  For every
non-negative \(\phi\in C_c^\infty(\R^3\times(-\infty,0))\),
\[
   2\iint |\nabla u_k|^2\phi
   \le
   \iint |u_k|^2(\partial_t\phi+\Delta\phi)
   +\iint (|u_k|^2+2P_k)u_k\cdot\nabla\phi ,
   \qquad P_k=p_k-a_k(t).
\]
This follows from \eqref{p1:eq:lei} by applying \eqref{p1:eq:lei} with temporal
cutoffs which increase to \(1\) on the time projection of \(\supp\phi\).  The
monotone convergence theorem passes the nonnegative dissipation term to the
displayed distributional form.

The terms with \(|u_k|^2\) converge strongly in \(L^1\) because
\[
   \||u_k|^2-|U|^2\|_{L^1(\supp\phi)}
   \le
   \|u_k-U\|_{L^2}
   \bigl(\|u_k\|_{L^2}+\|U\|_{L^2}\bigr)\to0.
\]
The cubic flux converges strongly since \(u_k\to U\) strongly in
\(L^3_{\loc}\), hence
\[
   |u_k|^2u_k\to |U|^2U
   \quad\hbox{in }L^1_{\loc}.
\]
For the pressure flux, \(u_k\cdot\nabla\phi\to U\cdot\nabla\phi\) strongly in
\(L^3\), while \(P_k\rightharpoonup Q\) weakly in \(L^{3/2}\).  Therefore
\[
   \iint P_k u_k\cdot\nabla\phi\to
   \iint Q\,U\cdot\nabla\phi .
\]
Finally,
\(\nabla u_k\rightharpoonup\nabla U\) in \(L^2_{\loc}\), and
\[
   w\mapsto\iint |w|^2\phi
\]
is weakly lower semicontinuous because \(\phi\ge0\) and the expression equals
\(\|\sqrt{\phi}\,w\|_{L^2}^2\).  Passing to the limit gives the local energy
inequality for \((U,Q)\) in distributional form.  The endpoint form in
\cref{p1:def:suitable} is recovered by applying the distributional inequality to
standard time cutoffs approximating characteristic functions of
\((t_1,t_2)\); weak \(L^2_{\loc}\)-continuity of local energy-class solutions
gives the endpoint traces for a.e. \(t_1<t_2\).
\end{proof}

\begin{lemma}[Inherited ancient Type I bound]\label{p1:lem:limit-bound}
For every compact interval \(I=[-S,-\sigma]\Subset(-\infty,0)\),
\[
   \|U\|_{L^\infty(\R^3\times I)}\le M\sigma^{-1/2}.
\]
\end{lemma}

\begin{proof}
\Cref{p1:lem:inherited-bound} gives
\[
   \|u_k\|_{L^\infty(\R^3\times I)}\le M\sigma^{-1/2}.
\]
Fix \(R<\infty\).  Since \(u_k\to U\) a.e. on \(B_R\times I\), the essential
supremum of \(U\) on \(B_R\times I\) cannot exceed \(M\sigma^{-1/2}\).  If it
did, \(|u_k|\) would exceed this bound by a positive amount on a positive
measure subset for large \(k\).  Letting \(R\to\infty\) through integers gives
the result on \(\R^3\times I\).
\end{proof}

\begin{proposition}[Ancient suitable Type I limit]\label{p1:prop:ancient-limit}
The pair \((U,Q)\) is an ancient suitable weak solution on
\(\R^3\times(-\infty,0)\) satisfying the interval Type I bound in
\cref{p1:lem:limit-bound}.
\end{proposition}

\begin{proof}
Combine \cref{p1:lem:limit-equation,p1:lem:suitability,p1:lem:limit-bound}.
\end{proof}

\section{Persistence of concentration and nontriviality}
\label{p1:sec:nontriviality}

\subsection{Persistence of concentration away from the terminal time}

The persistence of concentration is tracked through the velocity concentration sequence in
\cref{p1:def:admissible-type-I}, supplied by
\cref{p1:prop:local-velocity-concentration}.  The no-escape condition says that
this velocity mass does not concentrate entirely in the terminal layer \(t=0\)
after rescaling.

\subsection{Scale selection}

\begin{lemma}[Scale selection from local concentration]\label{p1:lem:scale-selection}
Let \(\lambda_k\downarrow0\) be an admissible concentration sequence in the
sense of \cref{p1:def:admissible-type-I}.  Then there are
\(\sigma\in(0,1)\), a compact cylinder
\[
   K_0=B_1\times(-1,-\sigma)\Subset\R^3\times(-\infty,0),
\]
and \(\eta_1>0\) such that
\begin{equation}\label{p1:eq:selected-lower-bound}
   \iint_{K_0}|u_k|^3\,dx\,dt\ge\eta_1
\end{equation}
for all \(k\).
\end{lemma}

\begin{proof}
By scale invariance of \(C\),
\[
   \int_{-1}^{0}\int_{B_1}|u_k|^3\,dx\,dt
   =
   C(u;z_*,\lambda_k)\ge\eta_v .
\]
The no-escape condition \eqref{p1:eq:velocity-no-escape} gives
\(\sigma\in(0,1)\) such that
\[
   \sup_k\int_{-\sigma}^{0}\int_{B_1}|u_k|^3\,dx\,dt
   \le \frac{\eta_v}{2}.
\]
Therefore
\[
   \int_{-1}^{-\sigma}\int_{B_1}|u_k|^3\,dx\,dt
   \ge \frac{\eta_v}{2}
\]
for every \(k\).  The assertion follows with
\(\eta_1=\eta_v/2\).
\end{proof}

\subsection{Nonzero ancient limits}

\begin{proposition}[Nonzero ancient limit]\label{p1:prop:nonzero}
For the scales chosen in \cref{p1:lem:scale-selection}, the limit \(U\) in
\cref{p1:prop:ancient-limit} is nonzero.  More precisely, there are a compact
cylinder \(K_0\Subset\R^3\times(-\infty,0)\) and \(\eta_1>0\) such that
\[
   \iint_{K_0}|U|^3\,dx\,dt\ge\eta_1 .
\]
\end{proposition}

\begin{proof}
Let \(K_0=B_1\times(-1,-\sigma)\) and \(\eta_1>0\) be given by
\cref{p1:lem:scale-selection}.  Since
\(K_0\Subset\R^3\times(-\infty,0)\), the compactness theorem gives, after
passing to the subsequence used to define \(U\),
\[
   u_k\to U\qquad\hbox{strongly in }L^3(K_0).
\]
Hence
\[
   \int_{K_0}|U|^3\,dx\,dt
   =
   \lim_{k\to\infty}\int_{K_0}|u_k|^3\,dx\,dt
   \ge\eta_1 .
\]
\end{proof}

Thus concentration generated by the singular point propagates through the
compactness limit.  The velocity concentration and no-escape condition are the
hypotheses that prevent the ancient limit from being identically zero.

\section{Centered variables and normalized Seregin limits}\label{p1:sec:seregin}

\subsection{The centered self-similar equation}

Set
\[
   y=\frac{x}{\sqrt{-t}},\qquad
   \tau=-\log(-t),\qquad
   V(y,\tau)=\sqrt{-t}\,U(y\sqrt{-t},t),
\]
and
\[
   \Pi(y,\tau)=(-t)Q(y\sqrt{-t},t),
\]
up to addition of a function of \(\tau\).

\begin{remark}[Pressure gauge transformation in centered variables]
\label{p1:rem:centered-gauge}
If \(Q\) is replaced by \(Q+a(t)\) for a function \(a\) of \(t\) only, then
\(\Pi\) is replaced by \(\Pi+e^{-\tau}a(-e^{-\tau})\), which is a function of
\(\tau\) only.  Thus the centered equation \eqref{p1:eq:centered} is
gauge-invariant: changing the pressure by a function of physical time amounts
to changing the centered pressure by a function of centered time.
\end{remark}

\begin{lemma}[Centered pullback identities]
\label{p1:lem:centered-pullback-identities}
The centered and physical variables are related by
\[
   t=-e^{-\tau},\qquad x=e^{-\tau/2}y,\qquad
   U(x,t)=e^{\tau/2}V(y,\tau),
\]
or equivalently
\[
   V(y,\tau)=e^{-\tau/2}U(e^{-\tau/2}y,-e^{-\tau}).
\]
For every \(\tau\in\R\), with \(t=-e^{-\tau}\),
\[
   \int_{\R^3}|U(x,t)|^3\,dx
   =
   \int_{\R^3}|V(y,\tau)|^3\,dy .
\]
\end{lemma}

\begin{proof}
The first identities are the definitions with \(-t=e^{-\tau}\).  Since
\[
   |U(x,t)|^3=e^{3\tau/2}|V(y,\tau)|^3,
   \qquad
   dx=e^{-3\tau/2}\,dy,
\]
the critical \(L^3\) identity follows by multiplication and integration.
\end{proof}

\begin{lemma}[Renormalized equation]\label{p1:lem:renormalized-equation}
The pair \((V,\Pi)\) solves \eqref{p1:eq:centered}.
\end{lemma}

\begin{proof}
Write \(s=-t=e^{-\tau}\), \(x=s^{1/2}y\), and
\[
   U(x,t)=s^{-1/2}V(y,\tau).
\]
The identities are first obtained in distributions by testing the equation
for \(U\) against the inverse change of variables.  After
\cref{p1:lem:smoothness}, the same identities hold classically for the smooth
representative.  Because \(s_t=-1\),
\(\tau_t=s^{-1}\), and \(y_t=(2s)^{-1}y\),
\[
   \partial_tU
   =
   s^{-3/2}\left(\partial_\tau V+\frac12V+\frac12y\cdot\nabla V\right),
\]
\[
   (U\cdot\nabla_x)U=s^{-3/2}(V\cdot\nabla)V,\qquad
   \Delta_xU=s^{-3/2}\Delta V,\qquad
   \nabla_xQ=s^{-3/2}\nabla\Pi.
\]
Substitution into the equation for \(U\) gives \eqref{p1:eq:centered}, and the
divergence condition is preserved by the spatial scaling.
\end{proof}

\subsection{Boundedness, smoothness, and local energy}

\begin{lemma}[Smoothness]\label{p1:lem:smoothness}
The ancient solution \(U\) is smooth on \(\R^3\times(-\infty,0)\), and
\((V,\Pi)\) is smooth on \(\R^3\times\R\) after fixing a pressure gauge.
\end{lemma}

\begin{proof}
On each compact cylinder \(Q_0\Subset\R^3\times(-\infty,0)\),
\cref{p1:lem:limit-bound} gives \(U\in L^\infty(Q_0)\).  Hence
\(U\in L^s_tL^r_x(Q_0)\) for every finite \(r,s\).  Choose \(r=s=10\), so
\[
   \frac2s+\frac3r=\frac12<1.
\]
Serrin's interior regularity criterion \cite{Serrin1962,Sohr2001} applies and
gives local H\"older regularity of \(U\) on a slightly smaller cylinder.
Since \(Q_0\) was arbitrary, this removes every possible interior singular
point in \(\R^3\times(-\infty,0)\).  The bootstrap starts from
\[
   \partial_tU-\Delta U+\nabla Q=-\divergence(U\otimes U),\qquad
   \divergence U=0,
\]
with \(U\otimes U\) locally bounded after Serrin regularity.  Interior Stokes
estimates on nested cylinders give \(W^{2,1}_q\)-regularity for every finite
\(q\).  Parabolic Sobolev embedding then improves the H\"older regularity of
\(\nabla U\), and differentiating the equation and iterating gives
\[
   U\in C^\infty(\R^3\times(-\infty,0)).
\]
The pressure is recovered locally from
\(-\Delta Q=\partial_i\partial_j(U_iU_j)\) plus a harmonic function and is
smooth after fixing a time-dependent gauge.  The centered change of variables
is a smooth diffeomorphism for \(t<0\), so \(V\) and \(\Pi\) are smooth.
\end{proof}

\begin{corollary}[Classical Type I bound]\label{p1:cor:classical-type-I-bound}
For the smooth representative from \cref{p1:lem:smoothness},
\[
   \|U(t)\|_{L^\infty(\R^3)}\le \frac{M}{\sqrt{-t}},
   \qquad t<0.
\]
\end{corollary}

\begin{proof}
Fix \(t_0<0\).  Choose \(\delta_n\downarrow0\) with
\[
   I_n=[t_0-\delta_n,t_0+\delta_n]\Subset(-\infty,0).
\]
\Cref{p1:lem:limit-bound} gives
\[
   \|U\|_{L^\infty(\R^3\times I_n)}
   \le M(-t_0-\delta_n)^{-1/2}.
\]
Since the representative is continuous, the essential spacetime bound holds
pointwise.  Letting \(n\to\infty\) gives the estimate at \(t_0\).
\end{proof}

\begin{lemma}[Renormalized boundedness and nontriviality]
\label{p1:lem:renorm-bounds}
The centered field satisfies
\[
   \|V\|_{L^\infty(\R^3\times\R)}\le M.
\]
Moreover \eqref{p1:eq:nonvanishing} holds for some compact cylinder
\(K\Subset\R^3\times\R\) and \(\eta_0>0\).
\end{lemma}

\begin{proof}
The \(L^\infty\) bound follows from \cref{p1:cor:classical-type-I-bound}:
\[
   |V(y,\tau)|=\sqrt{-t}|U(y\sqrt{-t},t)|\le M.
\]
For nontriviality, \cref{p1:prop:nonzero} gives a positive local \(L^3\) lower bound
on a compact cylinder in \((x,t)\).  Under
\((x,t)=(e^{-\tau/2}y,-e^{-\tau})\),
\[
   dx\,dt=e^{-5\tau/2}\,dy\,d\tau,\qquad
   U=e^{\tau/2}V,
\]
and hence
\[
   |U|^3\,dx\,dt=e^{-\tau}|V|^3\,dy\,d\tau .
\]
On the compact image cylinder the factor \(e^{-\tau}\) is bounded above and
below by positive constants, so the positive lower bound becomes
\eqref{p1:eq:nonvanishing}, with a possibly different \(\eta_0\).
\end{proof}

\begin{lemma}[Local energy bounds in centered variables]
\label{p1:lem:centered-local-energy}
Let \(V\) be obtained from an admissible finite-energy Type I singularity.  For every
\(R<\infty\) and
compact interval \(I\subset\R\),
\[
   \int_I\int_{B_R}|V|^2\,dy\,d\tau
   +
   \int_I\int_{B_R}|\nabla V|^2\,dy\,d\tau<\infty .
\]
After subtracting a function of \(\tau\), the associated pressure satisfies
\[
   \Pi-b(\tau)\in L^{3/2}(B_R\times I).
\]
\end{lemma}

\begin{proof}
The unrenormalized limit \((U,Q)\) is suitable and has local energy,
dissipation, and pressure bounds on compact negative-time cylinders by
\cref{p1:lem:energy-diss,p1:lem:pressure,p1:prop:ancient-limit}; the bounds for the
limit follow from the weak lower semicontinuity and weak compactness already
used in \cref{p1:lem:suitability}.  On each compact
\((y,\tau)\)-cylinder, the map
\[
   (y,\tau)\mapsto (x,t)=(e^{-\tau/2}y,-e^{-\tau})
\]
has Jacobian \(dx\,dt=e^{-5\tau/2}\,dy\,d\tau\), with
\[
   V=e^{-\tau/2}U,\qquad
   \nabla_yV=e^{-\tau}\nabla_xU,\qquad
   \Pi=e^{-\tau}Q .
\]
If \(I=[\tau_1,\tau_2]\), then \(e^{-\tau}\) is bounded above and below by
positive constants on \(I\).  Therefore
\[
   \int_I\int_{B_R}|V|^2\,dy\,d\tau
   \le C(I)\int_{-e^{-\tau_1}}^{-e^{-\tau_2}}
       \int_{B_{C(I)R}}|U|^2\,dx\,dt,
\]
and similarly
\[
   \int_I\int_{B_R}|\nabla_yV|^2\,dy\,d\tau
   \le C(I)\int_{-e^{-\tau_1}}^{-e^{-\tau_2}}
       \int_{B_{C(I)R}}|\nabla_xU|^2\,dx\,dt.
\]
For the pressure, subtract a spatial-average gauge \(a(t)\) for \(Q\) on the
corresponding \(x\)-ball and set \(b(\tau)=e^{-\tau}a(-e^{-\tau})\).  Then
\(\Pi-b=e^{-\tau}(Q-a)\), and the same bounded-weight change of variables
gives \(\Pi-b\in L^{3/2}(B_R\times I)\).
\end{proof}

\begin{remark}\label{p1:rem:no-global-tightness}
\Cref{p1:lem:centered-local-energy} is local.  It does not give uniform
\(L^3\)-tightness in \(y\), nor a global \(L^\infty_\tau L^3_y\) bound.  This
is the reason for retaining the remaining case.
\end{remark}

\subsection{Normalized Seregin ancient limits}

\begin{definition}[Admissible Type I extraction]
\label{p1:def:admissible-type-I-extraction}
An admissible Type I extraction is one of the following two blow-up
procedures:
\begin{enumerate}[label=\textup{(\roman*)}]
\item the finite-energy Type I extraction in \Cref{p1:def:admissible-type-I}
and \Cref{p1:thm:extraction};
\item the local pointwise Type I extraction in
\Cref{p0:def:paper0-admissible-local-typeI-sequence,p0:lem:local-seregin-extraction}.
\end{enumerate}
In both cases the conclusion is a smooth bounded solution of the centered equation
\eqref{p1:eq:centered} with compact local \(L^3\) nonvanishing.
\end{definition}

\begin{definition}[Raw extracted normalized Seregin ancient limit]
\label{p1:def:seregin-limit}
A smooth bounded solution \((V,\Pi)\) of \eqref{p1:eq:centered} is a
\emph{raw extracted normalized Seregin ancient limit} at \(z_*\) if it is
produced by an admissible finite-energy Type I extraction in the sense of
\Cref{p1:def:admissible-type-I-extraction}, or from an admissible local
pointwise Type I concentration sequence at \(z_*\) in the sense of
\Cref{p0:def:paper0-admissible-local-typeI-sequence,p0:lem:local-seregin-extraction},
and satisfies the compact spacetime \(L^3\) nonvanishing normalization
\eqref{p1:eq:nonvanishing}.
\end{definition}

\begin{corollary}[Raw extracted limits are nonzero]
\label{p1:cor:raw-extracted-nonzero}
Every raw extracted normalized Seregin ancient limit satisfies the compact
spacetime nonvanishing condition \eqref{p1:eq:nonvanishing}.  In particular,
it is not identically zero.
\end{corollary}

\begin{proof}
This is part of the definition of a raw extracted normalized Seregin ancient
limit.  For the finite-energy extraction the condition is produced by
\Cref{p1:lem:scale-selection,p1:prop:nonzero,p1:lem:renorm-bounds}; for the
local pointwise Type I concentration sequence it is produced by
\Cref{p0:lem:paper0-local-typeI-into-terminal-dichotomy,p0:lem:local-seregin-extraction}.
\end{proof}

\begin{lemma}[Time-translation invariant nonvanishing]
\label{p1:lem:compact-to-invariant-nonzero}
Let \((V,\Pi)\) be a raw extracted normalized Seregin ancient limit
satisfying \eqref{p1:eq:nonvanishing} on the compact cylinder
\(K\Subset\R^3\times\R\) with constant \(\eta_{\rm raw}=\eta_0>0\).  Then
there exist \(R_0>0\) and \(\eta_*>0\), depending only on \(K\) and
\(\eta_{\rm raw}\), such that
\begin{equation}\label{p1:eq:invariant-nonvanishing}
   \sup_{\tau\in\R}
   \int_{B_{R_0}}|V(y,\tau)|^3\,dy
   \ge \eta_* .
\end{equation}
\end{lemma}

\begin{proof}
Choose \(R_0>0\) and an interval \([a,b]\subset\R\) with
\(K\subset B_{R_0}\times[a,b]\).  Then
\[
   \eta_{\rm raw}
   \le
   \iint_K |V|^3\,dy\,d\tau
   \le
   \int_a^b\int_{B_{R_0}}|V(y,\tau)|^3\,dy\,d\tau
   \le (b-a)\sup_{\tau\in\R}\int_{B_{R_0}}|V(y,\tau)|^3\,dy .
\]
Take \(\eta_*=\eta_{\rm raw}/(b-a)\).
\end{proof}

\begin{definition}[Raw generated Seregin state space]
\label{p1:def:raw-generated-seregin-space}
\label{p1:def:seregin-collection}
For a fixed triple \((u,p,z_*)\), fix invariant normalization data
\((R_0,\eta_*)\) obtained from a raw extracted profile by
\Cref{p1:lem:compact-to-invariant-nonzero}.  Let
\[
   \calS_{\mathrm{raw}}(u,p,z_*;R_0,\eta_*)
\]
be the collection of raw extracted normalized Seregin ancient limits
produced by admissible finite-energy Type I or admissible local pointwise
Type I extractions at \(z_*\) which satisfy the invariant nonvanishing
condition
\[
   \sup_{\tau\in\R}
   \int_{B_{R_0}} |V(y,\tau)|^3\,dy \ge \eta_* .
\]
If different raw extractions yield different invariant normalization pairs,
the construction below is carried out separately for each pair.

The \emph{raw generated descendant hull}
\[
   \calS_{\mathrm{raw\text{-}gen}}(u,p,z_*;R_0,\eta_*)
\]
is, by definition, the smallest class containing
\(\calS_{\mathrm{raw}}(u,p,z_*;R_0,\eta_*)\) and closed under logarithmic time
translations and locally smooth limits that retain:
\begin{enumerate}[label=(\roman*)]
\item the inherited centered \(L^\infty\) bound;
\item smoothness and the centered local-pressure ancient formulation;
\item the invariant local velocity nonvanishing condition
\eqref{p1:eq:invariant-nonvanishing} with the fixed parameters
\((R_0,\eta_*)\).
\end{enumerate}
The class is closed under locally smooth limits that still satisfy the
invariant \((R_0,\eta_*)\)-nonvanishing condition; no claim is made that
this condition is lower semicontinuous under arbitrary uncentered locally
smooth limits.  An element of
\(\calS_{\mathrm{raw\text{-}gen}}(u,p,z_*;R_0,\eta_*)\) is called
a \emph{generated raw Seregin state}.  When the invariant normalization data
are clear from the extracted profile under discussion, we write
\(\calS_{\mathrm{raw}}(u,p,z_*)\) for
\(\calS_{\mathrm{raw}}(u,p,z_*;R_0,\eta_*)\) and
\(\calS_{\mathrm{raw\text{-}gen}}(u,p,z_*)\) for
\(\calS_{\mathrm{raw\text{-}gen}}(u,p,z_*;R_0,\eta_*)\).  For notational
economy, \(\calS(u,p,z_*)\) denotes this raw generated class.
\end{definition}

\begin{definition}[Mild stratum of the generated state space]
\label{p1:def:mild-stratum}
The mild stratum of the raw generated state space is
\[
   \calS_{\mathrm{mild}}(u,p,z_*;R_0,\eta_*)
   \subset
   \calS_{\mathrm{raw\text{-}gen}}(u,p,z_*;R_0,\eta_*)
\]
consisting of those generated states \(V\) satisfying the finite-interval
projected centered mild formulation, where
\[
   L=\Delta-\frac12y\cdot\nabla-\frac12,\qquad S(s)=e^{sL}.
\]
For every \(\sigma<\tau\),
\[
   V(\tau)
   =
   S(\tau-\sigma)V(\sigma)
   -
   \int_\sigma^\tau
   S(\tau-s)\mathbb P\nabla\cdot(V\otimes V)(s)\,ds
\]
in the weak-* topology of \(L^\infty(\R^3)\).
\end{definition}

\begin{remark}[Raw generated descendant hulls]
\label{p1:rem:generated-hull}
No assertion is made here that an arbitrary locally smooth limit of translates
is itself an extracted limit from the original physical sequence without a
separate diagonal generation argument.  Instead, the raw generated descendant
hull is included in the definition of the class to which the residual closure
hypothesis is applied.  This is the class on which the residual set
\(\calR(\calS)\) is formed in the final assembly.  The projected mild
formulation is imposed only on the mild stratum, not on the whole raw hull.
\end{remark}

\begin{lemma}[Raw generated closure preserves normalized ancient structure]
\label{p1:lem:raw-generated-closure-stability}
Every element of
\(\calS_{\mathrm{raw\text{-}gen}}(u,p,z_*;R_0,\eta_*)\) is a smooth bounded
centered ancient local-pressure solution, satisfies the inherited centered
\(L^\infty\) bound, and
retains the time-translation invariant local \(L^3\) nonvanishing
normalization \eqref{p1:eq:invariant-nonvanishing} with the fixed
\((R_0,\eta_*)\).
\end{lemma}

\begin{proof}
Raw extracted elements have these raw properties by
\Cref{p1:thm:extraction,p0:lem:local-seregin-extraction} together with
\Cref{p1:lem:compact-to-invariant-nonzero}, which converts the raw compact
cylinder mass lower bound into the invariant condition
\eqref{p1:eq:invariant-nonvanishing}.  Logarithmic time translations
preserve the autonomous centered equation, the centered
\(L^\infty\) bound, smoothness, and the local-pressure formulation; they also
preserve the supremum on the left of
\eqref{p1:eq:invariant-nonvanishing}, since the supremum is over all
\(\tau\in\R\).  For a locally smooth limit of such translates, smoothness
and the equation pass by local convergence; the \(L^\infty\) bound passes
by pointwise convergence.  The raw generated hull is closed only under limits
retaining the invariant nonvanishing condition with the fixed \((R_0,\eta_*)\),
so \eqref{p1:eq:invariant-nonvanishing} remains part of the structure.
\end{proof}

\begin{lemma}[Mild stratum stability]
\label{p1:lem:mild-stratum-stability}
The class \(\calS_{\mathrm{mild}}(u,p,z_*;R_0,\eta_*)\) is invariant under
logarithmic time translations.  Moreover, locally smooth limits of elements of
\(\calS_{\mathrm{mild}}\) with a common \(L^\infty\) bound remain in
\(\calS_{\mathrm{mild}}\), provided the limiting state belongs to the raw
generated hull and retains the fixed invariant nonvanishing normalization.
\end{lemma}

\begin{proof}
Time translation preserves the finite-interval projected mild identity because
the centered equation and the centered Stokes semigroup are autonomous.
Consider a locally smooth sequence \(V_j\in\calS_{\mathrm{mild}}\) with a common
\(L^\infty\) bound, converging locally smoothly to \(V\).  For fixed
\(\sigma<\tau\), the linear term \(S(\tau-\sigma)V_j(\sigma)\) converges to
\(S(\tau-\sigma)V(\sigma)\) in distributions and weak-* \(L^\infty_{\loc}\).
The Duhamel terms pass to the limit against compactly supported test fields by
the Oseen-kernel estimate in \Cref{p1:lem:centered-stokes-estimates}, the common
\(L^\infty\) bound, and local smooth convergence of \(V_j\otimes V_j\).  Indeed,
after testing against a compactly supported solenoidal field, the adjoint
Oseen kernel belongs to \(L^1_y\), with time-singularity integrable on
\((\sigma,\tau)\); truncating this \(L^1_y\) kernel outside a large ball and
using the common \(L^\infty\) bound controls the spatial tails uniformly,
while local smooth convergence handles the truncated part.  Hence
the projected mild formula holds for \(V\) on \([\sigma,\tau]\).  The
convergence against compactly supported tests upgrades to weak-* \(L^\infty\)
convergence against arbitrary \(L^1(\R^3)\) tests by approximating the
\(L^1\) test function by compactly supported functions and using the common
\(L^\infty\) bound on the sequence and on the limit.  The remaining
raw generated and nonvanishing requirements are part of the stated limiting
hypotheses.
\end{proof}

\begin{remark}
The raw compact spacetime \(L^3\) nonvanishing condition is part of the
normalization, not an additional hypothesis.  It is inherited from the
selected velocity concentration and the no-escape condition through
\cref{p1:lem:scale-selection,p1:prop:nonzero} for the finite-energy extraction,
and, for the local pointwise Type I extraction, through
\Cref{p0:lem:paper0-local-typeI-into-terminal-dichotomy,p0:lem:local-seregin-extraction}.
The raw generated descendant hull uses the
time-translation invariant consequence \eqref{p1:eq:invariant-nonvanishing};
this is the nonzero normalization used below for generated raw states.
\end{remark}

\begin{corollary}[Generated raw Seregin states are nonzero]
\label{p1:cor:normalized-nonzero}
Every generated raw Seregin state in
\(\calS_{\mathrm{raw\text{-}gen}}(u,p,z_*;R_0,\eta_*)\) satisfies the
time-translation invariant local \(L^3\) nonvanishing condition
\eqref{p1:eq:invariant-nonvanishing}; in particular it is not identically
zero.
\end{corollary}

\begin{proof}
For raw finite-energy extractions this is
\Cref{p1:lem:scale-selection,p1:prop:nonzero,p1:lem:renorm-bounds} together
with \Cref{p1:lem:compact-to-invariant-nonzero}.  For the
raw local pointwise Type I concentration sequence it is
\Cref{p0:lem:paper0-local-typeI-into-terminal-dichotomy,p0:lem:local-seregin-extraction}
together with \Cref{p1:lem:compact-to-invariant-nonzero}.
Generated closure elements retain
\eqref{p1:eq:invariant-nonvanishing} by
\Cref{p1:lem:raw-generated-closure-stability}.
\end{proof}

\begin{proof}[Proof of \cref{p1:thm:extraction}]
Choose an admissible concentration sequence
\(\lambda_k\downarrow0\) from \cref{p1:def:admissible-type-I}, and form the
rescaled sequence \eqref{p1:eq:rescaled-sequence}.  The compact ancient limit is
\cref{p1:prop:ancient-limit}; the strong and weak convergence are
\cref{p1:prop:compactness}.  The no-escape argument in
\cref{p1:lem:scale-selection}, followed by strong \(L^3\)-compactness away from
\(t=0\), gives nontriviality through \cref{p1:prop:nonzero}.  The pointwise
Type I bound is \cref{p1:cor:classical-type-I-bound}, and the centered equation,
smoothness, boundedness, and raw compact \(L^3\) nonvanishing are
\cref{p1:lem:renormalized-equation,p1:lem:smoothness,p1:lem:renorm-bounds}.
\end{proof}

\section{Liouville classes and the remaining class}\label{p1:sec:classes}

Fix a raw generated descendant hull
\[
   \calS=\calS_{\mathrm{raw\text{-}gen}}(u,p,z_*;R_0,\eta_*)
\]
in the sense of \Cref{p1:def:raw-generated-seregin-space}.  Complements below
are taken relative to this raw hull.  Write
\[
   \calS_{\mathrm{mild}}
   =
   \calS_{\mathrm{mild}}(u,p,z_*;R_0,\eta_*)
\]
for its mild stratum.  Structural conditions involving the unrenormalized
ancient field refer to
\[
   U(x,t)=(-t)^{-1/2}V(x/\sqrt{-t},-\log(-t)).
\]

\begin{definition}[Small-amplitude class]\label{p1:def:small-class}
Let \(\varepsilon_{\rm small}>0\) be the small-data Liouville constant for
bounded mild ancient solutions of \eqref{p1:eq:centered}.  We say
\(V\in\calC_{\mathrm{small}}\) if
\[
   V\in\calS_{\mathrm{mild}}
   \qquad\text{and}\qquad
   \|V\|_{L^\infty(\R^3\times\R)}\le \varepsilon_{\rm small}.
\]
\end{definition}

\begin{definition}[Stationary \(L^3\) class]\label{p1:def:stationary-class}
We say \(V\in\calC_{\mathrm{stat},L^3}\) if
\[
   V(y,\tau)=W(y)
   \quad\hbox{and}\quad
   W\in L^3(\R^3).
\]
\end{definition}

\begin{definition}[Uniformly \(L^3\)-tight class]\label{p1:def:tight-class}
We say \(V\in\calC_{L^3\mathrm{-tight}}\) if
\[
   V\in\calS_{\mathrm{mild}}
\]
and
\begin{equation}\label{p1:eq:L3-tightness}
   \lim_{R\to\infty}
   \sup_{\tau\in\R}\int_{|y|>R}|V(y,\tau)|^3\,dy=0 .
\end{equation}
\end{definition}

Every \(V\in\calS\) is bounded.  \Cref{p2:thm:tight-liouville} shows that, for
mild states, the tail condition \eqref{p1:eq:L3-tightness} gives the endpoint
sequence-\(L^3\) control needed after physical pullback; no separate
\(L^\infty_\tau L^3_y\) hypothesis is required for the tight-class exclusion.

\begin{definition}[Structure and decay class]
\label{p1:def:known-structure-decay-class}
A generated raw Seregin state \(V\in\calS_{\mathrm{raw\text{-}gen}}\) belongs
to \(\calC_{\mathrm{sd}}\) if \(V\in\calS_{\mathrm{mild}}\) and its physical
ancient representative \(U\) satisfies at least one of the following analytic
hypotheses:
\begin{enumerate}[label=(\roman*)]
\item \(U\) is axisymmetric without swirl;
\item \(U\) is axisymmetric and satisfies the pointwise scale-invariant bound
\[
   |U(x,t)|\le \frac{C}{r}
\]
for some \(C<\infty\), where \(r\) is the cylindrical radius;
\item \(U\) is axisymmetric and
\[
   rU_\theta\in L_t^\infty L_x^p(\R^3\times(-\infty,0))
\]
for some \(1\le p<\infty\);
\item \(U\) is axisymmetric, periodic in the axial variable, and
\[
   rU_\theta\in L^\infty(\R^3\times(-\infty,0));
\]
\item the physical ancient representative \(U\) satisfies the anchored
Gallay--Wayne perturbative vorticity condition of
\Cref{p3:def:gw-anchored-condition}.
\end{enumerate}
\end{definition}

\begin{remark}
Only hypotheses for which a cited Liouville theorem gives \(U\equiv0\) are
included in \(\calC_{\mathrm{sd}}\).  Broader qualitative descriptions are
not included unless they imply one of the listed hypotheses.  In the
anchored weighted-vorticity alternative, the Biot--Savart normalization
is used at the anchor times to identify the Gallay--Wayne Cauchy problem with
the anchored physical flow.  After the anchored argument gives
\(\nabla\times U\equiv0\), boundedness gives a spatially constant velocity
field, and the Type I bound removes the constant.
Membership in \(\calC_{\mathrm{sd}}\) is checked profile by profile.  No
closedness of the structural class under all generated limits is used in the
final contradiction; whenever a generated limit satisfies one of the listed
conditions and lies in the mild stratum, the corresponding Liouville theorem is
applied directly to that profile.
\end{remark}

\begin{lemma}[Structural quantities under centering]
\label{p1:lem:structure-centering}
The centered change of variables preserves axisymmetry and absence of swirl.
If \((r,z,t)\) and \((\rho,Z,\tau)\) are related by
\[
   \rho=\frac r{\sqrt{-t}},\qquad Z=\frac z{\sqrt{-t}},
   \qquad \tau=-\log(-t),
\]
then the swirl variables satisfy
\[
   \rho V_\theta(\rho,Z,\tau)=rU_\theta(r,z,t).
\]
\end{lemma}

\begin{proof}
The map \(x\mapsto y=x/\sqrt{-t}\) is radial in the cylindrical variables and
does not mix the \(e_r,e_\theta,e_z\) frame.  Since
\(V(y,\tau)=\sqrt{-t}\,U(x,t)\), the \(\theta\)-components obey
\(V_\theta=\sqrt{-t}\,U_\theta\), and the displayed identities follow.
\end{proof}

\begin{definition}[Remaining class]\label{p1:def:remaining-class}
The remaining class is
\[
   \calR(\calS)
   =
   \calS_{\mathrm{raw\text{-}gen}}(u,p,z_*;R_0,\eta_*)
   \setminus
   \left(
      \calC_{\mathrm{small}}
      \cup
      \calC_{\mathrm{stat},L^3}
      \cup
      \calC_{L^3\mathrm{-tight}}
      \cup
      \calC_{\mathrm{sd}}
   \right).
\]
\end{definition}

\begin{remark}
The first four classes are not intended to be disjoint.  The remaining class is
the complement of their union, so overlaps among the Liouville classes cause no
ambiguity.  The remaining class is a set-theoretic complement, not a new
structural hypothesis: it carries no analytic content beyond ``not in any of
the four named classes.''  The residual closure hypothesis
\Cref{p1:hyp:no-remainder} is the sole input used to remove this complement in
the final contradiction.
\end{remark}

\begin{remark}[Non-mild states are residual unless otherwise classified]
The ambient state space is the raw generated Seregin hull.  The explicit
Liouville classes include exactly the hypotheses needed to invoke their
corresponding Liouville theorems.  In particular, a generated state for which
the projected mild formulation has not been established is not placed in the
small, tight, or structure/decay Liouville classes merely because it satisfies
the corresponding pointwise or tail condition.  Such a state lies in the
remaining class unless it belongs to the stationary \(L^3\) branch.
\end{remark}

\begin{remark}[Single-profile versus collection-level tightness]
\label{p1:rem:single-vs-collection}
The exclusion of a single generated raw Seregin state in
\(\calC_{L^3\mathrm{-tight}}\) uses only the single-profile tail condition
together with mildness and boundedness, via \Cref{p2:thm:tight-liouville}.  The stronger
collection-level tightness modulus appearing in
\Cref{p2:def:tight-normalized-collection} is used only in the
auxiliary compactness and minimal-element material of Part III; no hidden common
modulus is imposed on individual generated raw Seregin states.
\end{remark}

\begin{proposition}[Exhaustion by defined classes]\label{p1:prop:exhaustion}
Let
\[
   \calS=\calS_{\mathrm{raw\text{-}gen}}(u,p,z_*;R_0,\eta_*)
\]
be a raw generated descendant hull.  Then every \(V\in\calS\) belongs either to
one of
\[
   \calC_{\mathrm{small}},\qquad
   \calC_{\mathrm{stat},L^3},\qquad
   \calC_{L^3\mathrm{-tight}},\qquad
   \calC_{\mathrm{sd}},
\]
or to
\[
   \calR(\calS)
   =
   \calS
   \setminus
   \left(
      \calC_{\mathrm{small}}
      \cup
      \calC_{\mathrm{stat},L^3}
      \cup
      \calC_{L^3\mathrm{-tight}}
      \cup
      \calC_{\mathrm{sd}}
   \right).
\]
\end{proposition}

\begin{proof}
Let
\[
   \calU=
      \calC_{\mathrm{small}}
      \cup
      \calC_{\mathrm{stat},L^3}
      \cup
      \calC_{L^3\mathrm{-tight}}
      \cup
      \calC_{\mathrm{sd}}.
\]
For \(V\in\calS\), either \(V\in\calU\), in which case it lies in at least one
Liouville class, or \(V\in\calS\setminus\calU=\calR(\calS)\).
\end{proof}

\section{Liouville hypotheses and class exclusions}\label{p1:sec:liouville}

\begin{lemma}[Centered Stokes kernel estimates]
\label{p1:lem:centered-stokes-estimates}
Let
\[
   L=\Delta-\frac12y\cdot\nabla-\frac12,
   \qquad S(s)=e^{sL}.
\]
There is a universal constant \(C_0\) such that, for all \(s>0\),
\[
   \|S(s)f\|_{L^\infty}\le e^{-s/2}\|f\|_{L^\infty}
\]
and
\[
   \|S(s)\mathbb P\nabla\cdot F\|_{L^\infty}
   \le C_0 e^{-s/2}(1-e^{-s})^{-1/2}\|F\|_{L^\infty}.
\]
\end{lemma}

\begin{proof}
The first estimate follows from the explicit Ornstein--Uhlenbeck kernel:
\[
   S(s)f(y)=e^{-s/2}
   \int_{\R^3}G_{1-e^{-s}}(e^{-s/2}y-z)f(z)\,dz,
\]
where \(G_a\) is the heat kernel with variance \(a\).  The kernel has total
mass one.  For the projected-divergence estimate, write the corresponding
kernel as the Oseen projected-gradient kernel after the same change of
variables.  More explicitly, the semigroup for
\(\Delta-\frac12y\cdot\nabla\) is
\[
   (e^{s(\Delta-\frac12y\cdot\nabla)}g)(y)
   =
   (e^{(1-e^{-s})\Delta}g)(e^{-s/2}y).
\]
Applying this to the distribution \(g=\mathbb P\nabla\cdot F\) and using the
projected-gradient Oseen kernel estimate
\[
   \|e^{\theta\Delta}\mathbb P\nabla\cdot F\|_\infty
   \le C\theta^{-1/2}\|F\|_\infty,\qquad \theta>0,
\]
gives the displayed estimate with \(\theta=1-e^{-s}\), together with the
extra factor \(e^{-s/2}\) from the zeroth-order term in \(L\).  The
projected-gradient Oseen estimate follows from the \(L^1\)-bound
\[
   \|\nabla O(\cdot,\theta)\|_{L^1}\le C\theta^{-1/2}
\]
for the gradient of the Oseen kernel \(O(\cdot,\theta)\).  This is the
bounded-data Stokes estimate used by Kato \cite{Kato1984}; it does not require
boundedness of the Leray projector itself on \(L^\infty\).
\end{proof}

\begin{remark}[Oseen interpretation in \(L^\infty\)]
\label{p1:rem:oseen-Linfty}
All occurrences of \(S(t)\mathbb P\nabla\cdot F\) in \(L^\infty\) are
understood through the projected heat-divergence, or Oseen, kernel
representation.  They do not use boundedness of the Leray projector as a
Calderon--Zygmund operator on \(L^\infty\).
\end{remark}

\begin{lemma}[Mild representation for generated mild states]
\label{p1:lem:mild-representation}
Every \(V\in\calS_{\mathrm{mild}}(u,p,z_*;R_0,\eta_*)\) satisfies, for all
\(\Theta>0\),
\[
   V(\tau)
   =
   S(\Theta)V(\tau-\Theta)
   -
   \int_0^\Theta
   S(\Theta-\sigma)\mathbb P\nabla\cdot
   (V\otimes V)(\tau-\Theta+\sigma)\,d\sigma .
\]
\end{lemma}

\begin{proof}
This is the finite-interval projected mild formulation in
\Cref{p1:def:mild-stratum}, applied with \(\sigma=\tau-\Theta\).  The centered
Stokes estimates in \Cref{p1:lem:centered-stokes-estimates} ensure that the
Duhamel term is well-defined in weak-* \(L^\infty\).
\end{proof}

\begin{lemma}[Small bounded ancient Liouville theorem]\label{p1:lem:small-liouville}
There is \(\varepsilon_{\rm small}>0\) such that every bounded mild ancient solution of
\eqref{p1:eq:centered} with
\[
   \|V\|_{L^\infty(\R^3\times\R)}\le\varepsilon_{\rm small}
\]
is identically zero.
\end{lemma}

\begin{proof}
Write the projected centered equation as
\[
   \partial_\tau V=LV-\mathbb P\nabla\cdot(V\otimes V),
   \qquad
   L=\Delta-\frac12 y\cdot\nabla-\frac12.
\]
Let \(S(s)=e^{sL}\).  By \cref{p1:lem:centered-stokes-estimates}, the nonlinear
kernel is integrable on \((0,\infty)\); set
\[
   A=C_0\int_0^\infty e^{-s/2}(1-e^{-s})^{-1/2}\,ds<\infty .
\]
For fixed \(\tau\) and \(T>0\), the mild identity on \([\tau-T,\tau]\) gives
\[
   V(\tau)
   =
   S(T)V(\tau-T)
   -
   \int_0^T
   S(T-\sigma)\mathbb P\nabla\cdot(V\otimes V)(\tau-T+\sigma)\,d\sigma .
\]
This Duhamel identity is valid for every \(T>0\) by
\Cref{p1:lem:mild-representation}.
If \(M=\|V\|_{L^\infty(\R^3\times\R)}\), then
\[
   \|V(\tau)\|_\infty\le e^{-T/2}M+AM^2.
\]
Letting \(T\to\infty\) and taking the supremum in \(\tau\) gives
\[
   M\le AM^2.
\]
If \(0<M<1/A\), this is impossible.  Choose
\(\varepsilon_{\rm small}=1/(2A)\); then
\(M\le\varepsilon_{\rm small}\) implies \(M=0\).
\end{proof}

\begin{lemma}[Stationary centered profiles]\label{p1:lem:stationary-profile}
If \(V(y,\tau)=W(y)\) is a smooth solution of \eqref{p1:eq:centered}, then after
subtracting a function of \(\tau\) from the pressure, \(W\) solves
\[
   (W\cdot\nabla)W+\nabla P
   =
   \Delta W-\frac12(W+y\cdot\nabla W),
   \qquad \divergence W=0.
\]
\end{lemma}

\begin{proof}
Since \(V\) is independent of \(\tau\), \(\partial_\tau V=0\).  Set
\[
   F(y)=
   \Delta W-\frac12(W+y\cdot\nabla W)-(W\cdot\nabla)W .
\]
The centered equation gives \(\nabla_y\Pi(\cdot,\tau)=F\) for a.e. \(\tau\).
Fix one such time \(\tau_0\) and put \(P(y)=\Pi(y,\tau_0)\).  Then
\(\nabla_y(\Pi(\cdot,\tau)-P)=0\) for a.e. \(\tau\), so
\(\Pi(y,\tau)-P(y)=b(\tau)\) is spatially constant.  Subtracting this gauge
from \(\Pi\) leaves the equation unchanged and yields the displayed
stationary self-similar equation for \(W\).
\end{proof}

\begin{lemma}[Spatially constant normalized profiles vanish]
\label{p1:lem:galilean-trivial-zero}
Let \(U\) be the physical ancient representative of a generated raw Seregin
state \(V\) in the mild stratum.  If \(U\) is spatially constant in the
physical Type I frame, then \(U\equiv0\) and hence \(V\equiv0\).
\end{lemma}

\begin{proof}
Mildness excludes parasitic spatially constant solutions: if
\(U(x,t)=c(t)\), the mild formula has zero nonlinear term and the heat
semigroup fixes constants, so \(c(t)=c(s)\) for all \(s<t<0\).  Thus
\(U(x,t)\equiv c\).  The centered profile is then
\[
   V(y,\tau)=e^{-\tau/2}c .
\]
Since generated raw Seregin states are bounded for all \(\tau\in\R\),
letting \(\tau\to-\infty\) forces \(c=0\).
\end{proof}

\begin{theorem}[Small and stationary class exclusions]
\label{p1:thm:classical-classes}
Let \(V\) be a generated raw Seregin state.  Then
\[
   V\notin \calC_{\mathrm{small}},
   \qquad
   V\notin \calC_{\mathrm{stat},L^3}.
\]
\end{theorem}

\begin{proof}
If \(V\in\calC_{\mathrm{small}}\), then
by definition \(V\in\calS_{\mathrm{mild}}\), and
\cref{p1:lem:mild-representation,p1:lem:small-liouville} give \(V\equiv0\),
contradicting the invariant nonvanishing normalization of
\Cref{p1:cor:normalized-nonzero}.

If \(V\in\calC_{\mathrm{stat},L^3}\), then
\(V(y,\tau)=W(y)\) with \(W\in L^3(\R^3)\).  By
\cref{p1:lem:stationary-profile}, \(W\) satisfies the stationary self-similar
profile equation after a pressure gauge change.
The stationary Leray theorem \Cref{p2:thm:NRS} gives \(W\equiv0\).  Again
the invariant nonvanishing normalization of \Cref{p1:cor:normalized-nonzero}
gives the contradiction.
\end{proof}

\begin{theorem}[Uniformly tight \(L^3\) closure]
\label{p1:thm:tight-liouville}
No generated raw Seregin state in a raw generated descendant hull
\(\calS_{\mathrm{raw\text{-}gen}}(u,p,z_*;R_0,\eta_*)\) belongs to
\(\calC_{L^3\mathrm{-tight}}\).
\end{theorem}

\begin{proof}
Let \(V\) be a generated raw Seregin state in
\(\calC_{L^3\mathrm{-tight}}\).  By definition
\(V\in\calS_{\mathrm{mild}}\) and satisfies the uniform tail condition.  Thus
\Cref{p2:thm:tight-liouville} gives \(V\equiv0\), contradicting
\Cref{p1:cor:normalized-nonzero}.
\end{proof}

\begin{theorem}[Structure and decay exclusions from Liouville theorems]
\label{p1:thm:known-structure-decay-liouville}
No nonzero generated raw Seregin state belongs to
\(\calC_{\mathrm{sd}}\).
\end{theorem}

\begin{proof}
Let \(V\in\calC_{\mathrm{sd}}\).  By definition of the structure/decay class,
\(V\in\calS_{\mathrm{mild}}\), and its physical representative
\[
   U(x,t)=(-t)^{-1/2}V(x/\sqrt{-t},-\log(-t))
\]
satisfies one of the listed structural hypotheses.  The inherited centered
\(L^\infty\) bound gives the physical Type I estimate
\[
   |U(x,t)|\le \|V\|_\infty(-t)^{-1/2}.
\]
By \Cref{p2:lem:physical-mild}, for every \(s<t<0\) the standard projected
mild formula holds on the interval \([s,t]\): choose \(T=t\) and apply the
finite-terminal restriction lemma.  Thus \(U\) is a smooth mild Type I ancient
solution in the sense of \Cref{p3:def:typeI-ancient}.  Therefore
\Cref{p3:thm:typeI-zero-hypotheses} gives \(U\equiv0\), and hence
\(V\equiv0\), contradicting the invariant nonvanishing normalization of
\Cref{p1:cor:normalized-nonzero}.
\end{proof}

\begin{proposition}[Exclusion of the Liouville classes]
\label{p1:prop:known-class-exclusion}
No nonzero generated raw Seregin state can belong to
\[
   \calC_{\mathrm{small}}
   \cup
   \calC_{\mathrm{stat},L^3}
   \cup
   \calC_{L^3\mathrm{-tight}}
   \cup
   \calC_{\mathrm{sd}}.
\]
\end{proposition}

\begin{proof}
If \(V\in\calC_{\mathrm{small}}\), then by definition
\(V\in\calS_{\mathrm{mild}}\), and the small mild Liouville theorem applies.
If \(V\in\calC_{\mathrm{stat},L^3}\), then the stationary self-similar profile
equation and the stationary Leray theorem apply.  These two branches are
excluded in \Cref{p1:thm:classical-classes}.  If
\(V\in\calC_{L^3\mathrm{-tight}}\), then by definition
\(V\in\calS_{\mathrm{mild}}\), so \Cref{p1:thm:tight-liouville} applies.  If
\(V\in\calC_{\mathrm{sd}}\), then by definition \(V\in\calS_{\mathrm{mild}}\)
and the physical representative satisfies the relevant Type I mild structural
hypothesis, so \Cref{p1:thm:known-structure-decay-liouville} applies.
\end{proof}

\section{Residual closure hypothesis}\label{p1:sec:criterion}

The preceding sections prove that every generated raw Seregin state in a raw
generated descendant hull associated with an admissible finite-energy Type I
extraction or an admissible local pointwise Type I concentration
sequence is either in one of the four non-residual Liouville classes or in the
complement \(\calR(\calS)\).  The four non-residual classes
(\(\calC_{\mathrm{small}}\), \(\calC_{\mathrm{stat},L^3}\),
\(\calC_{L^3\mathrm{-tight}}\), \(\calC_{\mathrm{sd}}\)) are excluded in
\Cref{p1:prop:known-class-exclusion}.  The remaining assumption for the Type I
contradiction is the residual closure statement below.  No result in the
preceding sections depends on the proof of residual closure; it is invoked
only in \Cref{p1:thm:remaining-criterion} and in the final assembly
\Cref{p1:thm:final-assembly}.

\begin{hypothesis}[Residual closure for raw generated Seregin state spaces]
\label{p1:hyp:no-remainder}
For every finite-energy suitable weak solution \((u,p)\), every terminal point
\(z_*=(x_*,T)\) satisfying the local pointwise Type I bound, and every
raw generated descendant hull produced by admissible local Type I concentration
sequences at \(z_*\) with invariant normalization data generated by a raw
extracted profile,
\[
   \calS=\calS_{\mathrm{raw\text{-}gen}}(u,p,z_*;R_0,\eta_*),
\]
one has
\[
   \calR\bigl(\calS_{\mathrm{raw\text{-}gen}}(u,p,z_*;R_0,\eta_*)\bigr)=\varnothing .
\]
The same assertion is assumed for raw generated descendant hulls produced by
admissible finite-energy Type I extractions in the finite-energy criterion
\Cref{p1:thm:remaining-criterion}.
The residual class is formed after removing only those states satisfying the
full hypotheses of the explicit Liouville branches.  Thus the residual closure
hypothesis also removes any generated raw state for which projected mildness is
not available and which is not covered by the stationary \(L^3\) branch.
\end{hypothesis}

\begin{remark}[Residual closure as a conditional assumption]
\label{p1:rem:companion-paper}
\Cref{p1:hyp:no-remainder} is the only residual-class assumption not established
from the preceding compactness and Liouville arguments in this paper.  It concerns
the raw generated descendant hull
\(\calS_{\mathrm{raw\text{-}gen}}(u,p,z_*;R_0,\eta_*)\) after the ordered removal of the
non-residual classes defined in \Cref{p1:sec:classes}.  The hypothesis is a
conditional analytic statement about the same raw generated state space
constructed here.  When this hypothesis is assumed,
\Cref{p1:thm:remaining-criterion} and
\Cref{p1:thm:final-assembly} give the Type I exclusions within the stated
finite-energy Type I extraction and local pointwise Type I concentration
sequence classes.  All uses of
\Cref{p1:hyp:no-remainder} are therefore conditional on this residual input.
\end{remark}

\begin{remark}[Interface with the residual paper]
\label{p1:rem:paper-IV-interface}
The companion residual paper is intended to prove
\[
   \calR\bigl(\calS_{\mathrm{raw\text{-}gen}}(u,p,z_*;R_0,\eta_*)\bigr)=\varnothing
\]
for the raw generated descendant hull defined in
\Cref{p1:def:raw-generated-seregin-space}.  The final assembly is formulated
only for this raw generated state space.
\end{remark}

\begin{corollary}[Rigidity in the remaining case]\label{p1:cor:remainder-rigidity}
Assume \Cref{p1:hyp:no-remainder}.
For every raw generated descendant hull
\(\calS=\calS_{\mathrm{raw\text{-}gen}}(u,p,z_*;R_0,\eta_*)\) associated with admissible
finite-energy Type I extractions or admissible local pointwise Type I
concentration sequences, every
element of \(\calR(\calS)\) is
identically zero.
\end{corollary}

\begin{proof}
The assertion follows directly from \cref{p1:hyp:no-remainder}.
\end{proof}

\begin{proposition}[Equivalence on generated raw Seregin states]
\label{p1:prop:remainder-equivalence}
For raw generated descendant hulls of generated raw Seregin states, the
residual-class hypothesis \cref{p1:hyp:no-remainder} implies
\cref{p1:cor:remainder-rigidity}.  Conversely, the corresponding rigidity
assertion removes the remaining class from every raw generated descendant hull.
\end{proposition}

\begin{proof}
If \cref{p1:hyp:no-remainder} holds, the rigidity statement follows immediately.
Conversely, assume the residual rigidity assertion for every raw generated
descendant hull: every element of \(\calR(\calS)\) is identically
zero.  If some raw generated descendant hull had \(V\in\calR(\calS)\),
the rigidity assertion would give \(V\equiv0\), directly contradicting the
invariant nonvanishing normalization of \Cref{p1:cor:normalized-nonzero}.  Hence no
nonzero generated raw Seregin state can remain in \(\calR(\calS)\).
By \Cref{p1:cor:normalized-nonzero}, every element of
\(\calS_{\mathrm{raw\text{-}gen}}(u,p,z_*;R_0,\eta_*)\), and hence every element of
\(\calR(\calS)\), is nonzero.  Therefore \(\calR(\calS)=\varnothing\).
\end{proof}

\begin{proof}[Proof of \cref{p1:thm:remaining-criterion}]
Assume \Cref{p1:hyp:no-remainder}.
Suppose that \((u,p,z_*)\) is an admissible finite-energy Type I singularity.
By \cref{p1:thm:extraction}, there exists a raw extracted normalized nonzero
Seregin ancient limit \(V\).  By
\Cref{p1:lem:compact-to-invariant-nonzero}, \(V\) satisfies the invariant
local nonvanishing condition.  Let
\[
   \calS=\calS_{\mathrm{raw\text{-}gen}}(u,p,z_*;R_0,\eta_*)
\]
be the raw generated descendant hull associated with \(z_*\), the admissible
extraction scales, and the corresponding fixed invariant normalization
component.  By construction, \(V\in\calS\).

By \cref{p1:prop:exhaustion}, either \(V\) belongs to one of the Liouville classes
or \(V\in\calR(\calS)\).  The first alternative is impossible by
\cref{p1:prop:known-class-exclusion}.  The second alternative is impossible by
\cref{p1:hyp:no-remainder}, applied to
\(\calS_{\mathrm{raw\text{-}gen}}(u,p,z_*;R_0,\eta_*)\).  Hence the assumed admissible finite-energy Type I
singularity cannot exist.
\end{proof}

The preceding theorem is conditional on the residual-class hypothesis for the
raw generated descendant hull associated with the finite-energy extraction.  The
compactness construction, pressure normalization, compact-to-invariant
nonvanishing passage, raw generated closure stability, and class exhaustion are
established above, with the local velocity concentration result recorded in
\cref{p1:prop:local-velocity-concentration}.  The structure and decay class is
restricted to the assumptions in \cref{p1:def:known-structure-decay-class} and
is routed through \Cref{p3:thm:typeI-zero-hypotheses}; any other structured
behavior lies in the remaining class and is addressed only by
\cref{p1:hyp:no-remainder}.  No whole-space critical-norm estimate on the
original solution is used in the conditional closure step.

\section{Elementary symmetries of the centered equation}
\label{p1:app:centered-symmetries}

\begin{lemma}[Logarithmic-time translation]\label{p1:lem:tau-translation}
If \(V\) solves \eqref{p1:eq:centered}, then
\[
   V_{\tau_0}(y,\tau)=V(y,\tau+\tau_0)
\]
also solves \eqref{p1:eq:centered}.  If
\[
   U^\lambda(x,t)=\lambda U(\lambda x,\lambda^2t),
\]
then the renormalized field associated with \(U^\lambda\) is
\[
   V^\lambda(y,\tau)=V(y,\tau-2\log\lambda).
\]
\end{lemma}

\begin{proof}
The first statement follows because \eqref{p1:eq:centered} is autonomous in
\(\tau\).  For the second, put \(s=-t=e^{-\tau}\).  Then
\[
   \sqrt{s}\,U^\lambda(y\sqrt{s},-s)
   =
   \lambda\sqrt{s}\,U(\lambda y\sqrt{s},-\lambda^2s)
   =
   V(y,-\log(\lambda^2s))
   =
   V(y,\tau-2\log\lambda).
\]
\end{proof}

\begin{proposition}[Dilation and translation are not symmetries]
\label{p1:prop:no-cascade}
The centered equation \eqref{p1:eq:centered} is autonomous in \(\tau\), but it is
not invariant under
\[
   V(y,\tau)\mapsto \alpha V(\alpha y,\alpha^2\tau)
\]
unless \(\alpha=1\).  It is also not invariant under constant translations
\[
   V(y,\tau)\mapsto V(y-a,\tau)
\]
unless \(a=0\).
\end{proposition}

\begin{proof}
For scaling, set \(Y=\alpha y\), \(T=\alpha^2\tau\), and
\(\widetilde V(y,\tau)=\alpha V(Y,T)\),
\(\widetilde\Pi(y,\tau)=\alpha^2\Pi(Y,T)\).  The non-drift part scales with
factor \(\alpha^3\), while
\[
   -\frac12(\widetilde V+y\cdot\nabla\widetilde V)
   =
   -\frac{\alpha}{2}(V+Y\cdot\nabla_YV)(Y,T).
\]
Using the equation for \(V\), the remaining term in the equation for \(\widetilde V\) is
\[
   \frac{\alpha-\alpha^3}{2}(V+Y\cdot\nabla_YV)(Y,T).
\]
If \(\alpha\ne1\), this forces \(V+Y\cdot\nabla V=0\).  Along each ray,
\[
   \frac{d}{dr}\bigl(rV(r\omega,T)\bigr)=0.
\]
Boundedness and smoothness at the origin force the constant to be zero, hence
\(V\equiv0\).  Thus a nonzero centered solution has no such scaling symmetry.

For translations, all translation-invariant terms transform by composition,
but the drift satisfies
\[
   y\cdot\nabla V(y-a,\tau)
   =
   (y-a)\cdot\nabla V(y-a,\tau)+a\cdot\nabla V(y-a,\tau).
\]
The extra term is absent from the original equation.  Thus spatial
translation is not a symmetry of the fixed centered equation unless \(a=0\),
apart from exceptional fixed solutions satisfying \(a\cdot\nabla V\equiv0\).
\end{proof}



\clearpage
\section*{Part III. Uniformly Tight Ancient Dynamics}

\subsection*{Overview}

Bounded ancient solutions of the three-dimensional Navier--Stokes equations
are considered in backward self-similar variables under uniform
\(L^3\)-tightness.  The compactness theory shows that bounded mild
ancient solutions are smooth, locally compact under time translation, and that
the common \(L^\infty_\tau L^3_y\) and tightness bounds used in compactness
statements are stable under locally smooth limits.  A minimal-element reduction
is then obtained for closed normalized uniformly tight collections:
any nonempty such collection contains a nonzero solution of minimal
\(L^\infty(\R^3\times\R)\) size.

For the time-translation hull of a bounded ancient solution, compact invariant
hulls and invariant probability measures are constructed by
Krylov--Bogolyubov averaging.  These measures satisfy a stationary statistical form of the
self-similar Navier--Stokes equation.  Passing to the barycenter produces an
explicit quadratic covariance term, isolating the obstruction to obtaining a
single stationary self-similar profile.
This compact-hull and invariant-measure material is auxiliary only; it is
not used in the final Type I contradiction.

Finally, the uniformly \(L^3\)-tight class is closed directly from its defining
hypotheses.  The tail condition, together with the boundedness of the ancient
solution, gives the sequence-\(L^3\) bound needed after the physical pullback.
Since this pullback preserves the critical \(L^3\)-norm on each time slice,
the endpoint ancient Liouville theorem of Albritton--Barker then forces the
pullback, and hence the centered solution, to vanish.  Therefore the uniformly
\(L^3\)-tight class of the Type I
ancient-limit reduction is excluded without any compact-hull Lyapunov
hypothesis.

\section{Introduction and main results}\label{p2:sec:introduction}

\subsection{The uniformly tight ancient class}

Consider the three-dimensional incompressible Navier--Stokes equations in
backward self-similar variables.  The centered velocity \(V\) and pressure \(P\)
solve
\begin{equation}\label{p2:eq:centered}
   \partial_\tau V+(V\cdot\nabla)V+\nabla P
   =
   \Delta V-\frac12(V+y\cdot\nabla V),
   \qquad \divergence V=0
\end{equation}
on \(\R^3\times\R\).  In \Cref{p1:thm:extraction}, admissible Type I
singularities are reduced to nonzero normalized ancient solutions of
\eqref{p2:eq:centered}.  The uniformly \(L^3\)-tight class is the class in which
\begin{equation}\label{p2:eq:L3-tightness}
   \lim_{R\to\infty}
   \sup_{\tau\in\R}\int_{|y|>R}|V(y,\tau)|^3\,dy=0 .
\end{equation}
The compactness statements also record the common bound
\begin{equation}\label{p2:eq:L3-bound}
   \sup_{\tau\in\R}\|V(\tau)\|_{L^3(\R^3)}<\infty ,
\end{equation}
but the endpoint exclusion of the tight class uses only boundedness and the
tail condition \eqref{p2:eq:L3-tightness}.

This condition on the centered limit is distinct from the original-solution condition
\[
   u\in L^\infty(0,T;L^3(\R^3)).
\]
For finite-energy suitable weak solutions, the original-solution assumption
\(u\in L^\infty_tL^3_x\) near a putative singular time is incompatible with
singularity by the theorem of Escauriaza--Seregin--\v{S}ver\'ak
\cite{ESS2003}.  Uniform \(L^3\)-tightness is imposed on centered ancient
solutions after blow-up.

\subsection{Minimal elements and compact hulls}

\begin{remark}[Auxiliary diagnostic material]
\label{p2:rem:auxiliary-dynamical-material}
The compact-hull, minimal-element, and invariant-measure material in this
part is auxiliary diagnostic material.  It is not used in the proof of
\Cref{p1:thm:final-assembly}; the final Type I contradiction uses only the
uniform \(L^3\)-tight Liouville theorem
\Cref{p2:thm:tight-liouville} and its interface corollary
\Cref{p2:cor:tight-input-for-paper-I}.  It is recorded here because the
covariance obstruction it identifies is the precise reason the residual
class cannot be removed by a stationary statistical reduction alone, and
because the same constructions are reused in the residual-branch paper.
\end{remark}

\begin{remark}[Paper I/Paper III interface and normalization alignment]
\label{p2:rem:paper-I-III-interface}
Raw extractions in Part II produce the local lower bound
\eqref{p1:eq:nonvanishing}.  By
\Cref{p1:lem:compact-to-invariant-nonzero}, every raw extracted profile
also satisfies the invariant local normalization used in
\eqref{p2:eq:local-nonzero}.  The raw generated class
\(\calS_{\mathrm{raw\text{-}gen}}(u,p,z_*;R_0,\eta_*)\) is closed under time translations and
locally smooth limits only within this invariantly normalized class, so
the closed normalized collections of \Cref{p2:def:closed-normalized}
include every raw generated descendant hull as a special case (with
\((R_0,\eta_0)\) given by the raw extraction parameters of
\Cref{p1:lem:compact-to-invariant-nonzero}).
\end{remark}

The first reduction is variational:
\[
   \hbox{nonempty tight normalized class}
   \quad\Longrightarrow\quad
   \hbox{\(L^\infty\)-minimal ancient element}.
\]
The normalization is a fixed local velocity nonvanishing condition.  It
prevents minimizing sequences from losing all mass by time translation.
Once a minimal element \(V_*\) is selected, its time-translation hull is
\[
   \calH(V_*)
   =
   \overline{\{\Theta_sV_*:\ s\in\R\}}^{C^\infty_{\loc}},
   \qquad
   (\Theta_sV)(y,\tau)=V(y,\tau+s).
\]
The hull is compact because bounded mild ancient solutions are locally smooth
with estimates depending only on their common \(L^\infty\) bound.

\subsection{Invariant measures and covariance}

Every compact invariant hull \(K\) carries invariant probability measures.
Indeed, Krylov--Bogolyubov averages
\[
   \mu_T=\frac1T\int_0^T\delta_{\Theta_sW_0}\,ds
\]
have weak subsequential limits.  Such invariant measures satisfy a stationary
statistical form of \eqref{p2:eq:centered}.  If
\[
   \overline W(y)=\int_K W(y,0)\,d\mu(W),
\]
then the averaged nonlinearity decomposes as
\[
   \int_K W_iW_j\,d\mu
   =
   \overline W_i\,\overline W_j+C_{ij}.
\]
Thus the barycenter solves the stationary self-similar equation only up to the
explicit covariance term \(C\).  This invariant-measure construction does
not prove that a recurrent compact hull contains a single stationary profile.

\subsection{\texorpdfstring{Endpoint \(L^3\) closure of the tight class}{Endpoint L3 closure of the tight class}}

The compact-hull and invariant-measure constructions identify the covariance
obstruction in a purely local topology, but they are not needed to close the
uniformly \(L^3\)-tight class.  The tail condition already gives the
sequence-\(L^3\) information required by the endpoint ancient Liouville theorem,
because the centered solution is bounded.  Under the physical change of variables
\[
   u(x,t)=(-t)^{-1/2}
   V\left(\frac{x}{\sqrt{-t}},-\log(-t)\right),
   \qquad t<0,
\]
the \(L^3\)-norm is invariant:
\[
   \|u(t)\|_{L^3_x}
   =
   \|V(-\log(-t))\|_{L^3_y}.
\]
Consequently every tail-tight centered ancient solution produces, after
restriction to any finite terminal time and translation of that terminal time
to \(0\), a bounded physical mild ancient solution whose terminal datum lies
in \(L^3(\R^3)\) and whose \(L^3\)-norm is bounded along a sequence
\(t_k\downarrow-\infty\).  The endpoint theorem of Albritton--Barker
\cite{AlbrittonBarker2019} then gives \(u\equiv0\), and the change of
variables gives \(V\equiv0\).

The endpoint closure uses only the boundedness and tail tightness already
present in the ancient solution class and introduces no additional estimate,
recurrence assumption, or Lyapunov functional.  The compactness and
invariant-measure results establish the
local smooth compactness theory, the existence of minimal elements and
compact trajectory hulls, and the stationary statistical identity with its
covariance term.  The endpoint \(L^3\) exclusion of the uniformly tight class
is then obtained by the physical pullback and the Albritton--Barker ancient
Liouville theorem.

Thus the uniformly \(L^3\)-tight class is closed under local smooth limits by
its defining hypotheses.  The compact-hull construction describes recurrent dynamics and the
covariance obstruction, but no compact-recurrence or Lyapunov assumption enters
the tight-class exclusion in \Cref{p2:thm:tight-liouville}.

\subsection{Main results}

The first main theorem is the analytic compactness and closure statement.

\begin{theorem}[Compactness and closure of uniformly tight ancient solutions]
\label{p2:thm:tight-compactness-intro}
Let \(\{V_n\}\) be a sequence of bounded mild ancient solutions of
\eqref{p2:eq:centered}.  Assume
\[
   \sup_n\|V_n\|_{L^\infty(\R^3\times\R)}<\infty .
\]
Then a subsequence converges locally smoothly to a bounded mild ancient
solution \(V\).

If in addition there are \(L<\infty\) and a modulus
\(\omega:[1,\infty)\to[0,\infty)\), \(\omega(R)\to0\), such that
\[
   \sup_n\sup_{\tau\in\R}\|V_n(\tau)\|_{L^3(\R^3)}\le L
\]
and
\[
   \sup_n\sup_{\tau\in\R}
   \int_{|y|>R}|V_n(y,\tau)|^3\,dy\le \omega(R)
   \qquad (R\ge1),
\]
then
\[
   \sup_{\tau\in\R}\|V(\tau)\|_{L^3(\R^3)}\le L,
   \qquad
   \sup_{\tau\in\R}\int_{|y|>R}|V(y,\tau)|^3\,dy\le\omega(R),
\]
and in particular \(V\) is uniformly \(L^3\)-tight.
\end{theorem}

Compatible pressure representatives also converge locally weakly modulo
time-dependent gauges under the additional \(L^\infty_\tau L^3_y\) assumption;
see \Cref{p2:lem:pressure-reconstruction}.

The second main theorem is the minimal-element reduction.

\begin{theorem}[Minimal element in a normalized uniformly tight class]
\label{p2:thm:minimal-tight-intro}
Let \(A\) be a nonempty collection of bounded mild ancient solutions satisfying
a common \(L^\infty\) bound, a common \(L^\infty_\tau L^3_y\) bound, a common
uniform \(L^3\)-tightness modulus, and a fixed local nonvanishing
normalization
\[
   \sup_{\tau\in\R}\int_{B_{R_0}}|V(y,\tau)|^3\,dy\ge \eta_0>0.
\]
Assume \(A\) is invariant under time translation and closed under locally
smooth limits which retain the normalization.  Then \(A\) contains an element
\(V_*\) minimizing
\[
   \|V\|_{L^\infty(\R^3\times\R)}
\]
among all \(V\in A\).  In particular \(V_*\not\equiv0\).
\end{theorem}

The third main theorem is the endpoint \(L^3\) tight-class Liouville
criterion.

\begin{theorem}[Uniformly tight Liouville criterion]
\label{p2:thm:tight-liouville-intro}
Let \(A\) be a collection of bounded mild ancient solutions of
\eqref{p2:eq:centered}.  Assume there is a modulus
\(\omega:[1,\infty)\to[0,\infty)\), \(\omega(R)\to0\), such that
\[
   \sup_{\tau\in\R}
   \int_{|y|>R}|V(y,\tau)|^3\,dy\le\omega(R)
   \qquad\text{for every }V\in A,\ R\ge1,
\]
and assume that \(A\) has a fixed local nonvanishing normalization:
\[
   \sup_{\tau\in\R}\int_{B_{R_0}}|V(y,\tau)|^3\,dy\ge\eta_0>0
   \qquad\text{for every }V\in A.
\]
Then
\[
   A=\varnothing .
\]
\end{theorem}

The proof of the last theorem is summarized by the implications
\[
\begin{array}{c}
\text{uniformly \(L^3\)-tight locally nonzero class }A\\[2mm]
\Downarrow\quad\text{boundedness plus tail tightness}\\[2mm]
\sup_k\|V(\tau_k)\|_{L^3(\R^3)}<\infty
\quad\text{for a sequence }\tau_k\to-\infty\\[2mm]
\Downarrow\quad\text{physical pullback and critical \(L^3\)-invariance}\\[2mm]
\text{mild ancient Navier--Stokes solution }u\\[1mm]
\text{with }\sup_k\|u(t_k)\|_{L^3(\R^3)}<\infty
\quad\text{for some }t_k\downarrow-\infty\\[2mm]
\Downarrow\quad\text{Albritton--Barker endpoint theorem}\\[2mm]
u=0\\[2mm]
\Downarrow\\[2mm]
\text{contradiction with local velocity nonvanishing in }A .
\end{array}
\]

\subsection{Relation to the Type I reduction}

\Cref{p1:thm:extraction} first reduces admissible Type I singularities
to raw extracted Seregin ancient limits.  The subsequent normalization step
fixes the time-translation invariant lower bound, and the generated descendant
hull is then organized into small, stationary \(L^3\), uniformly tight,
structure and decay, and residual classes.
The tight-class exclusion requires no additional compact-dynamical hypothesis:
the uniformly \(L^3\)-tight class contains no
nonzero generated raw ancient state in the mild stratum.

\subsection{Organization}

\Cref{p2:sec:centered-topology} fixes the centered equation, mild formulation,
time translation, and local smooth topology.  \Cref{p2:sec:smoothing-compactness}
proves smoothing, pressure reconstruction, stability, and compactness.
\Cref{p2:sec:tightness-rigidity} proves closure of uniform \(L^3\)-tightness and
recalls stationary \(L^3\) rigidity.  \Cref{p2:sec:minimal-elements} proves the
minimal-element reduction.  \Cref{p2:sec:compact-hulls} constructs compact
trajectory hulls and minimal invariant subsets.  \Cref{p2:sec:measures-covariance}
constructs invariant measures and identifies the covariance obstruction.
\Cref{p2:sec:endpoint-L3-closure} proves the physical pullback identities,
applies the endpoint \(L^3\) ancient Liouville theorem, and states the
consequence used in the Type I reduction.

\section{Centered ancient solutions and the local smooth topology}
\label{p2:sec:centered-topology}

\subsection{The centered equation}

Throughout this part \(V\) denotes a centered ancient velocity and \(P\) its
pressure.  An element of a compact hull is denoted by \(W\), and a
stationary profile by \(w(y)\).  Define
\[
   L V=\Delta V-\frac12 y\cdot\nabla V-\frac12 V .
\]
After applying the Helmholtz--Leray projection \(\mathbb P\), the centered
equation becomes
\begin{equation}\label{p2:eq:projected}
   \partial_\tau V
   =
   LV-\mathbb P\divergence(V\otimes V).
\end{equation}

\subsection{The Ornstein--Uhlenbeck--Stokes semigroup}

Let \(S(a)=e^{aL}\), \(a>0\).  It is given by
\begin{equation}\label{p2:eq:OU-semigroup}
   (S(a)f)(y)
   =
   e^{-a/2}
   \int_{\R^3}
   \frac{\exp\!\left(-\frac{|z-e^{-a/2}y|^2}{4(1-e^{-a})}\right)}
        {[4\pi(1-e^{-a})]^{3/2}}
   f(z)\,dz .
\end{equation}
The factor \(e^{-a/2}\) represents the damping caused by the term \(-V/2\).

\begin{lemma}[Centered Stokes kernel estimate]\label{p2:lem:centered-stokes-kernel}
For \(a>0\) and \(\phi\in C_c^\infty(\R^3;\R^3)\),
\begin{equation}\label{p2:eq:stokes-adjoint-estimate}
   \|(S(a)\mathbb P\divergence)^*\phi\|_{L^1(\R^3)}
   \le
   C e^{-a/2}(1-e^{-a})^{-1/2}\|\phi\|_{L^1(\R^3)}.
\end{equation}
Consequently the Duhamel term in \eqref{p2:eq:mild} is well-defined by duality
for bounded \(V\).
\end{lemma}

\begin{proof}
Set \(\lambda=e^{-a/2}\) and \(\theta=1-e^{-a}\).  From
\eqref{p2:eq:OU-semigroup},
\[
   S(a)f(y)=e^{-a/2}(e^{\theta\Delta}f)(\lambda y).
\]
Thus
\[
   S(a)^*\phi(z)
   =
   e^{-a/2}\lambda^{-3}
   (e^{\theta\Delta}\psi)(z),
   \qquad
   \psi(x)=\phi(x/\lambda),
\]
and \(\|\psi\|_{L^1}=\lambda^3\|\phi\|_{L^1}\).  The Leray projection is a
homogeneous Fourier multiplier of degree zero, so it commutes with this
dilation.  The classical heat--Stokes kernel estimate
\[
   \|(e^{\theta\Delta}\mathbb P\divergence)^*\phi\|_{L^1}
   \le C\theta^{-1/2}\|\phi\|_{L^1},
   \qquad \theta>0,
\]
follows from the Oseen kernel bound; see
\cite{FabesJonesRiviere1972,GigaMiyakawa1985,Sohr2001}.  Substituting
\(\theta=1-e^{-a}\) and using
\[
   e^{-a/2}\lambda^{-3}\|\psi\|_{L^1}=e^{-a/2}\|\phi\|_{L^1}
\]
gives \eqref{p2:eq:stokes-adjoint-estimate}.  Since
\((1-e^{-a})^{-1/2}\) is integrable at \(a=0\), the duality pairing in
\eqref{p2:eq:mild} is integrable on finite time intervals for
\(V\in L^\infty\).
\end{proof}

\subsection{Riesz transform convention}

The convention is that \(R_i\) has Fourier symbol \(i\xi_i/|\xi|\).  Thus
\(R_iR_j\) has symbol \(-\xi_i\xi_j/|\xi|^2\), and
\[
   P=R_iR_j(F_{ij})
   \quad\Longleftrightarrow\quad
   -\Delta P=\partial_i\partial_jF_{ij}.
\]
Pressures are determined only modulo functions of time unless a whole-space
Leray gauge is explicitly chosen.

\subsection{Bounded mild ancient solutions}

\begin{definition}[Bounded mild ancient solution]\label{p2:def:mild}
A vector field \(V:\R^3\times\R\to\R^3\) is a bounded mild ancient solution of
\eqref{p2:eq:centered} if
\[
   V\in L^\infty(\R^3\times\R),\qquad \divergence V=0
\]
in distributions, and for every \(\sigma<\tau\),
\begin{equation}\label{p2:eq:mild}
   V(\tau)
   =
   S(\tau-\sigma)V(\sigma)
   -
   \int_\sigma^\tau
   S(\tau-s)\mathbb P\divergence(V\otimes V)(s)\,ds
\end{equation}
in the weak-* topology of \(L^\infty(\R^3)\).
\end{definition}

The Duhamel term in \eqref{p2:eq:mild} is interpreted by duality.  For
\(\phi\in C_c^\infty(\R^3;\R^3)\),
\[
   \left\langle
      S(\tau-s)\mathbb P\divergence(V\otimes V)(s),\phi
   \right\rangle
   =
   \left\langle
      V\otimes V(s),
      (S(\tau-s)\mathbb P\divergence)^*\phi
   \right\rangle ,
\]
and \Cref{p2:lem:centered-stokes-kernel} makes the scalar integrand absolutely
integrable on \(s\in(\sigma,\tau)\).

\subsection{Time translation}

For \(s\in\R\) define
\[
   (\Theta_sV)(y,\tau)=V(y,\tau+s).
\]
The equation, the mild formulation, the \(L^\infty\) norm, and the
\(L^\infty_\tau L^3_y\) and tail bounds used below are invariant under
\(\Theta_s\).  Indeed, replacing \((\sigma,\tau)\) by
\((\sigma+s,\tau+s)\) in \eqref{p2:eq:mild} and changing variables in the time
integral gives the mild formula for \(\Theta_sV\).

\subsection{The locally smooth topology}

\begin{definition}[Local smooth topology]\label{p2:def:local-smooth-topology}
The locally smooth topology on \(C^\infty_{\loc}(\R^3\times\R;\R^3)\) is
generated by the seminorms
\[
   \|V\|_{m,N}
   =
   \max_{|\alpha|+j\le m}
   \sup_{\substack{|y|\le N\\ |\tau|\le N}}
   |\partial_y^\alpha\partial_\tau^jV(y,\tau)|,
   \qquad m,N\in\N .
\]
One compatible metric is
\[
   d(U,V)
   =
   \sum_{m,N=1}^{\infty}
   2^{-(m+N)}
   \frac{\|U-V\|_{m,N}}{1+\|U-V\|_{m,N}}.
\]
\end{definition}

Thus \(V_n\to V\) locally smoothly if and only if the convergence is in
\(C^m(K)\) for every compact cylinder
\(K\Subset\R^3\times\R\) and every \(m\ge0\).

\section{Smoothing, pressure reconstruction, and compactness}
\label{p2:sec:smoothing-compactness}

\subsection{Local smoothing}

\begin{lemma}[Smoothing of bounded mild solutions]\label{p2:lem:mild-smoothing}
Every bounded mild ancient solution is smooth on \(\R^3\times\R\) and satisfies
the projected equation \eqref{p2:eq:projected} classically.  Moreover, for each
compact cylinder \(K\Subset\R^3\times\R\) and each integer \(m\ge0\),
\[
   \|V\|_{C^m(K)}
   \le
   C\bigl(m,K,\|V\|_{L^\infty(\R^3\times\R)}\bigr).
\]
The same estimate holds uniformly for any family with a common
\(L^\infty\) bound.
\end{lemma}

\begin{proof}
Fix a compact time interval \(J\Subset\R\) containing the time projection of
\(K\), and set \(t=-e^{-\tau}\).  On the corresponding compact interval
\(t(J)\Subset(-\infty,0)\), define
\[
   U(x,t)=(-t)^{-1/2}V(x/\sqrt{-t},-\log(-t)).
\]
The map \((x,t)\mapsto (x/\sqrt{-t},-\log(-t))\) is a smooth diffeomorphism
between \(\R^3\times t(J)\) and \(\R^3\times J\), with all derivatives bounded
on compact subcylinders.

The centered mild formula \eqref{p2:eq:mild} transforms, under this parabolic
change of variables, into the usual projected Navier--Stokes mild formula on
every compact subinterval of \(t(J)\):
\[
   U(t)=e^{(t-s)\Delta}U(s)
   -
   \int_s^t e^{(t-r)\Delta}\mathbb P\divergence(U\otimes U)(r)\,dr,
   \qquad s<t.
\]
This transformation uses only the projected mild equation: the Leray projection
is homogeneous of degree zero and commutes with the parabolic dilation, while
the Ornstein--Uhlenbeck kernel in \eqref{p2:eq:OU-semigroup} becomes the heat
kernel with time parameter \(t-s\).  No pressure formulation is used at this
stage.  The \(L^\infty\) norm of \(U\) on \(\R^3\times t(J)\) is bounded by a
constant depending only on \(J\) and \(\|V\|_{L^\infty}\).

The standard local smoothing theorem for bounded projected mild
Navier--Stokes solutions \cite{Kato1984,GigaMiyakawa1985} gives first local
derivative bounds away from the initial time in any mild representation.  In
the present ancient setting one may choose the initial time strictly before the
compact subcylinder under consideration.  Iterating the differentiated
Duhamel formula and then applying the interior parabolic bootstrap for the
Stokes system \cite{Ladyzhenskaya1968,Sohr2001} gives uniform bounds for all
derivatives of \(U\) on compact subcylinders of \(\R^3\times t(J)\).  The mild
formula may then be differentiated in time, so \(U\) solves the projected
Navier--Stokes equation classically.  Returning to \((y,\tau)\) gives the
asserted estimates for \(V\) and the classical projected equation
\eqref{p2:eq:projected}.  The argument is uniform for families with a common
\(L^\infty\) bound.
\end{proof}

\subsection{Pressure reconstruction}

\begin{lemma}[Local pressure reconstruction]\label{p2:lem:local-pressure}
Let \(V\) be a smooth bounded solution of the projected centered equation
\eqref{p2:eq:projected}.  Then for every ball \(B_R\) and compact interval
\(I\Subset\R\) there exists
\[
   P_R\in C^\infty(B_R\times I),
\]
unique up to an additive function of time, such that
\[
   \partial_\tau V+(V\cdot\nabla)V+\nabla P_R
   =
   \Delta V-\frac12(V+y\cdot\nabla V),
   \qquad \divergence V=0,
\]
on \(B_R\times I\).
\end{lemma}

\begin{proof}
Set
\[
   F=
   \Delta V-\frac12(V+y\cdot\nabla V)
   -\partial_\tau V-(V\cdot\nabla)V .
\]
Since \((V\cdot\nabla)V=\divergence(V\otimes V)\) and \(V\) satisfies
\eqref{p2:eq:projected}, self-adjointness of the Leray projection gives
\[
   \int_I\int_{B_R}F\cdot\Phi\,dy\,d\tau=0
\]
for every \(\Phi\in C_c^\infty(B_R\times I;\R^3)\) satisfying
\(\divergence_y\Phi=0\).  The local de Rham theorem, applied in the spatial
variables with \(\tau\) as a parameter, gives a scalar distribution \(P_R\) on
\(B_R\times I\), unique up to a distribution depending only on time, such that
\(\nabla_yP_R=F\).  Since \(F\) is smooth, all spatial derivatives of \(P_R\)
of positive order are smooth.  Fixing, for instance, the spatial mean of
\(P_R\) over \(B_R\) to be zero for each \(\tau\) selects a representative
which is smooth on \(B_R\times I\).  This proves the local pressure
formulation.
\end{proof}

\begin{lemma}[Leray pressure and local gauges]\label{p2:lem:pressure-reconstruction}
Let \(V\) be a bounded mild ancient solution satisfying
\[
   \sup_{\tau\in\R}\|V(\tau)\|_{L^3(\R^3)}<\infty .
\]
Then the whole-space Leray pressure
\[
   P=R_iR_j(V_iV_j)
\]
belongs to \(L^\infty(\R;L^{3/2}(\R^3))\).  Moreover, if \(V_n\to V\) locally
smoothly, the \(V_n\) have a common \(L^\infty_\tau L^3_y\) bound, and
\[
   P_n=R_iR_j(V_{n,i}V_{n,j}),
\]
then on every compact cylinder \(B_R\times[\tau_1,\tau_2]\) there are bounded
measurable gauges \(a_n(\tau)\) and \(a(\tau)\) such that, after passing to a
subsequence,
\[
   P_n-a_n(\tau)
   \rightharpoonup P-a(\tau)
   \quad\hbox{weakly in }L^{3/2}(B_R\times[\tau_1,\tau_2]).
\]
Equivalently, after replacing the limiting pressure by an equivalent
representative modulo a function of time, the weak limit may be denoted by
\(P\).
\end{lemma}

\begin{proof}
The first assertion follows from the Calderon--Zygmund boundedness of Riesz
transforms on \(L^{3/2}(\R^3)\) \cite{Stein1970}:
\[
   \|P(\tau)\|_{L^{3/2}}
   \le C\|V_i(\tau)V_j(\tau)\|_{L^{3/2}}
   \le C\|V(\tau)\|_{L^3}^2 .
\]
For the local stability statement, choose
\(\chi\in C_c^\infty(B_{4R})\) with \(\chi\equiv1\) on \(B_{3R}\).  Decompose
on each time slice
\[
   P_n=R_iR_j(\chi V_{n,i}V_{n,j})
       +R_iR_j((1-\chi)V_{n,i}V_{n,j}) .
\]
The local term converges strongly in \(L^q(B_{2R}\times[\tau_1,\tau_2])\) for
every finite \(q\).  This is because \(V_n\to V\) smoothly on
\(\overline B_{4R}\times[\tau_1,\tau_2]\), so
\(\chi V_{n,i}V_{n,j}\to\chi V_iV_j\) strongly in every finite \(L^q\), and
the Riesz transforms are Calderon--Zygmund operators on \(L^q\), \(1<q<\infty\).

It remains to control the exterior part modulo constants.  Let \(K_{ij}\) be
the convolution kernel of \(R_iR_j\).  For \(x\in B_R\) and
\(|z|>4R\),
\[
   |K_{ij}(x-z)-K_{ij}(-z)|\le C_R |z|^{-4}.
\]
Define the exterior gauges
\[
   b_n(\tau)
   =
   \int_{\R^3} K_{ij}(-z)(1-\chi(z))V_{n,i}(z,\tau)V_{n,j}(z,\tau)\,dz
\]
and \(b(\tau)\) by the same formula with \(V\) in place of \(V_n\).
The cutoff removes the singularity at the origin, and the integral is
absolutely convergent at infinity because, on the dyadic annulus
\(A_m=\{2^mR<|z|<2^{m+1}R\}\),
\[
   \int_{A_m}|z|^{-3}|V_n(z,\tau)|^2\,dz
   \le
   C(2^mR)^{-3}
   \|V_n(\tau)\|_{L^3(A_m)}^2|A_m|^{1/3}
   \le C L^2(2^mR)^{-2}.
\]
The same dyadic estimate shows that \(b_n\) and \(b\) are bounded uniformly on
\([\tau_1,\tau_2]\).  The maps \(\tau\mapsto b_n(\tau)\) and
\(\tau\mapsto b(\tau)\) are measurable, because they are limits of measurable
truncated integrals.
After subtracting \(b_n(\tau)\), the exterior contribution on \(B_R\) is
\[
   E_n(x,\tau)
   =
   \int_{\R^3}
   [K_{ij}(x-z)-K_{ij}(-z)]
   (1-\chi(z))V_{n,i}(z,\tau)V_{n,j}(z,\tau)\,dz .
\]
The same dyadic estimate, now using \(|z|^{-4}\), gives
\[
   \|E_n(\cdot,\tau)\|_{L^\infty(B_R)}
   \le C(R)L^2
\]
uniformly in \(n\) and \(\tau\).  Splitting the exterior integral into
\(\{4R<|z|<R'\}\) and \(\{|z|>R'\}\), the first part converges by local smooth
convergence, while the second part is bounded uniformly by \(C L^2/R'\).  Letting
first \(n\to\infty\) and then \(R'\to\infty\) identifies the local limit of
\(E_n\) with the corresponding exterior contribution for \(V\).  Therefore,
with \(a_n=b_n\) and \(a=b\), the decompositions of \(P_n-a_n\) and \(P-a\)
have local parts converging strongly and exterior parts converging
distributionally with a uniform local \(L^{3/2}\) bound.  Hence
\[
   P_n-a_n(\tau)\rightharpoonup P-a(\tau)
   \quad\text{weakly in }L^{3/2}(B_R\times[\tau_1,\tau_2]).
\]
\end{proof}

\subsection{Stability under locally smooth limits}

\begin{lemma}[Closure of the mild formulation]\label{p2:lem:mild-limit}
Let \(V_n\) be bounded mild ancient solutions satisfying
\[
   \sup_n\|V_n\|_{L^\infty(\R^3\times\R)}<\infty .
\]
If \(V_n\to V\) locally smoothly, then \(V\) is a bounded mild ancient solution
of \eqref{p2:eq:centered}.
\end{lemma}

\begin{proof}
The common \(L^\infty\) bound passes to \(V\) by pointwise convergence on
compact subsets.  Divergence freeness passes distributionally.  Fix
\(\sigma<\tau\) and test \eqref{p2:eq:mild} for \(V_n\) against
\(\phi\in C_c^\infty(\R^3;\R^3)\).  The endpoint terms are
\[
   \langle V_n(\tau),\phi\rangle,
   \qquad
   \langle V_n(\sigma),S(\tau-\sigma)^*\phi\rangle .
\]
The first converges by local smooth convergence.  For the second, note that
\(S(\tau-\sigma)^*\phi\in L^1(\R^3)\).  Approximate this \(L^1\) function by
a compactly supported smooth function and use the common \(L^\infty\) bound to
control the \(L^1\) approximation error uniformly in \(n\).

For the Duhamel term, write
\[
   \left\langle
      S(\tau-s)\mathbb P\divergence(V_n\otimes V_n)(s),\phi
   \right\rangle
   =
   \left\langle
      V_n(s)\otimes V_n(s),
      (S(\tau-s)\mathbb P\divergence)^*\phi
   \right\rangle .
\]
For fixed \(s\), the adjoint test belongs to \(L^1\); pairing convergence
again follows by compact approximation in \(L^1\), local smooth convergence,
and the common \(L^\infty\) bound on \(V_n\otimes V_n\).  The membership in
\(L^1\) follows from \Cref{p2:lem:centered-stokes-kernel}: the operator
\((S(\tau-s)\mathbb P\divergence)^*\) acts on the compactly supported smooth
test field \(\phi\) by integration against the (smooth) adjoint OU--Stokes
kernel, which has finite \(L^1\) norm in the spatial variable for every
\(\tau-s>0\).  The approximation
error is uniform in \(n\) for this fixed \(s\).  Moreover
\Cref{p2:lem:centered-stokes-kernel} gives the integrable domination
\[
   C\sup_n\|V_n\|_{L^\infty}^2
   e^{-(\tau-s)/2}(1-e^{-(\tau-s)})^{-1/2}\|\phi\|_{L^1},
   \qquad s\in(\sigma,\tau).
\]
Dominated convergence in \(s\) passes the Duhamel integral to the limit.
Thus \(V\) satisfies \eqref{p2:eq:mild}.
\end{proof}

\subsection{Compactness of bounded families}

\begin{proposition}[Compactness of bounded ancient families]
\label{p2:prop:bounded-compactness}
Let \(\{V_n\}\) be bounded mild ancient solutions satisfying
\[
   \sup_n\|V_n\|_{L^\infty(\R^3\times\R)}<\infty .
\]
Then a subsequence converges locally smoothly to a bounded mild ancient
solution \(V\).
\end{proposition}

\begin{proof}
By \Cref{p2:lem:mild-smoothing}, the sequence is bounded in every
\(C^m\)-norm on every compact cylinder.  Arzela--Ascoli on a countable
exhaustion \(B_N\times[-N,N]\), followed by a diagonal extraction over
\((m,N)\), gives a subsequence converging locally smoothly to a smooth limit.
\Cref{p2:lem:mild-limit} shows that the limit is again a bounded mild ancient
solution.
\end{proof}

\section{\texorpdfstring{Uniform \(L^3\)-tightness and stationary \(L^3\) rigidity}{Uniform L3-tightness and stationary L3 rigidity}}
\label{p2:sec:tightness-rigidity}

\subsection{\texorpdfstring{Uniform \(L^3\)-tightness}{Uniform L3-tightness}}

\begin{definition}[Uniform \(L^3\)-tightness]\label{p2:def:L3-tight}
A bounded ancient solution \(V\) is uniformly \(L^3\)-tight if
\[
   \lim_{R\to\infty}
   \sup_{\tau\in\R}
   \int_{|y|>R}|V(y,\tau)|^3\,dy=0 .
\]
\end{definition}

For the compactness statements in this section we sometimes assume, in
addition, a common \(L^\infty_\tau L^3_y\) bound.  The endpoint Liouville
argument in \Cref{p2:sec:endpoint-L3-closure} does not take this bound as a
separate hypothesis; it obtains the required sequence estimate from boundedness
and the tail condition above.

\subsection{Closedness under locally smooth convergence}

\begin{lemma}[Closure of uniform \(L^3\)-tightness]\label{p2:lem:tight-closure}
Let \(V_n\) be bounded mild ancient solutions satisfying
\[
   \sup_n\sup_{\tau\in\R}\|V_n(\tau)\|_{L^3(\R^3)}\le L<\infty
\]
and assume that they have a common tightness modulus: for some
\(\omega:[1,\infty)\to[0,\infty)\) with \(\omega(R)\to0\),
\[
   \sup_n\sup_{\tau\in\R}
   \int_{|y|>R}|V_n(y,\tau)|^3\,dy\le \omega(R),
   \qquad R\ge1.
\]
If \(V_n\to V\) locally smoothly, then
\[
   \sup_{\tau\in\R}\|V(\tau)\|_{L^3(\R^3)}\le L
\]
and
\[
   \sup_{\tau\in\R}
   \int_{|y|>R}|V(y,\tau)|^3\,dy\le \omega(R),
   \qquad R\ge1.
\]
In particular \(V\) is uniformly \(L^3\)-tight.
\end{lemma}

\begin{proof}
Fix \(\tau\).  On \(B_R\), local smooth convergence gives
\[
   \int_{B_R}|V(y,\tau)|^3\,dy
   =
   \lim_{n\to\infty}\int_{B_R}|V_n(y,\tau)|^3\,dy
   \le L^3 .
\]
Here \(L\) bounds the \(L^3\) \emph{norm} (so \(L^3\) bounds the integral of
the cube), while \(\omega(R)\) below bounds the integral on the tail
directly, not its cube.
Letting \(R\to\infty\) and using monotone convergence gives
\(\|V(\tau)\|_{L^3}\le L\).  Since \(\tau\) was arbitrary, the same bound is
uniform in time.

For the tail bound, fix \(1\le R<R'\).  On the bounded annulus
\(\{R<|y|<R'\}\), dominated convergence gives
\[
   \int_{R<|y|<R'}|V(y,\tau)|^3\,dy
   =
   \lim_{n\to\infty}
   \int_{R<|y|<R'}|V_n(y,\tau)|^3\,dy
   \le \omega(R).
\]
Letting \(R'\to\infty\) and using monotone convergence gives
\[
   \int_{|y|>R}|V(y,\tau)|^3\,dy\le \omega(R).
\]
The estimate is independent of \(\tau\), so taking the supremum in time gives
the asserted tail bound.  Since \(\omega(R)\to0\), \(V\) is uniformly
\(L^3\)-tight.
\end{proof}

\begin{proof}[Proof of \Cref{p2:thm:tight-compactness-intro}]
The compactness assertion is \Cref{p2:prop:bounded-compactness}.  The
\(L^\infty_\tau L^3_y\) and tightness closure assertions follow from
\Cref{p2:lem:tight-closure} with the displayed constants \(L\) and \(\omega\).
\end{proof}

\subsection{Stationary time-translate limits}

The stationary self-similar profile equation used below is
\begin{equation}\label{p2:eq:stationary-profile}
   (w\cdot\nabla)w+\nabla p
   =
   \Delta w-\frac12(w+y\cdot\nabla w),
   \qquad \divergence w=0 .
\end{equation}

\begin{theorem}[Stationary limits under uniform \(L^3\)-tightness]
\label{p2:thm:stationary-limits-tight}
Let \(V\) be uniformly \(L^3\)-tight.  Let \(\tau_n\in\R\), and suppose
\[
   \Theta_{\tau_n}V\to W
\]
locally smoothly.  If \(W\) is stationary, then \(W\equiv0\).
\end{theorem}

\begin{proof}
The sequence \(\Theta_{\tau_n}V\) has the same \(L^3\) bound and tightness
modulus as \(V\).  By \Cref{p2:lem:tight-closure}, the limit \(W\) belongs to
\(L^\infty_\tau L^3_y\).  If \(W\) is stationary, write
\(W(y,\tau)=w(y)\).  Then \(w\in L^3(\R^3)\).

By \Cref{p2:lem:mild-limit}, \(W\) is a bounded mild ancient solution.  The
smoothness from \Cref{p2:lem:mild-smoothing} and the local pressure lemma
\Cref{p2:lem:local-pressure} give smooth local pressures \(p_N\) on \(B_N\) at
time \(0\).  Since
\(\partial_\tau W=0\), each \(p_N\) satisfies
\[
   (w\cdot\nabla)w+\nabla p_N
   =
   \Delta w-\frac12(w+y\cdot\nabla w)
   \quad\text{on }B_N .
\]
On overlaps the gradients of the local pressures agree, so the pressures
differ only by constants.  Adjusting those constants and passing to the
compatible union gives a scalar distribution \(p\) on \(\R^3\) for which
\eqref{p2:eq:stationary-profile} holds in distributions.  The stationary
\(L^3\) rigidity theorem \Cref{p2:thm:NRS} applies and gives \(w=0\).
\end{proof}

\subsection{Stationary self-similar rigidity}

\begin{theorem}[Stationary self-similar rigidity]\label{p2:thm:NRS}
Let \(w\in C^\infty(\R^3;\R^3)\cap L^3(\R^3)\) solve
\eqref{p2:eq:stationary-profile} in distributions for some scalar distribution
\(p\), with \(\divergence w=0\).  Then \(w\equiv0\).
\end{theorem}

\begin{proof}
We use the theorem of Ne\v{c}as, R\r{u}\v{z}i\v{c}ka, and
\v{S}ver\'ak \cite{NRS1996}, applied to the stationary Leray self-similar
system.  The smooth distributional formulation stated here is a special case
of their weak \(L^3(\R^3)\) theorem.
\end{proof}

\section{Normalized collections and minimal ancient solutions}
\label{p2:sec:minimal-elements}

\subsection{Closed normalized collections}

\begin{definition}[Closed normalized collection]\label{p2:def:closed-normalized}
A collection \(A\) of bounded mild ancient solutions is called a closed
normalized collection if:
\begin{enumerate}[label=\textup{(\roman*)}]
\item there is \(M<\infty\) such that
\[
   \|V\|_{L^\infty(\R^3\times\R)}\le M
   \qquad\text{for all }V\in A;
\]
\item there are \(R_0>0\) and \(\eta_0>0\) such that
\begin{equation}\label{p2:eq:local-nonzero}
   \sup_{\tau\in\R}
   \int_{B_{R_0}}|V(y,\tau)|^3\,dy
   \ge \eta_0
   \qquad\text{for all }V\in A;
\end{equation}
\item \(A\) is invariant under time translation:
\[
   V\in A,\ s\in\R
   \quad\Longrightarrow\quad
   \Theta_sV\in A;
\]
\item if \(V_n\in A\), \(V_n\to V\) locally smoothly, and \(V\) satisfies the
same \((R_0,\eta_0)\)-nonvanishing condition \eqref{p2:eq:local-nonzero}, then
\(V\in A\).
\end{enumerate}
\end{definition}

\subsection{Time-centering}

\begin{lemma}[Time-centering]\label{p2:lem:time-centering}
Let \(A\) be a closed normalized collection.  If \(V\in A\), then for every
\(\delta>0\) there exists \(s\in\R\) such that
\[
   \int_{B_{R_0}}|\Theta_sV(y,0)|^3\,dy
   \ge \eta_0-\delta.
\]
\end{lemma}

\begin{proof}
By \eqref{p2:eq:local-nonzero}, the supremum in time of the local \(L^3\) mass is
at least \(\eta_0\).  Choose \(s\) within \(\delta\) of this supremum.  Since
\((\Theta_sV)(y,0)=V(y,s)\), the displayed estimate follows.  Time-translation
invariance keeps \(\Theta_sV\) in \(A\).
\end{proof}

\subsection{Stability of local nonvanishing}

\begin{lemma}[Stability of local nonvanishing]\label{p2:lem:nonzero-stability}
If \(V_n\to V\) locally smoothly and
\[
   \int_{B_{R_0}}|V_n(y,0)|^3\,dy\ge \eta_0-o(1),
\]
then
\[
   \int_{B_{R_0}}|V(y,0)|^3\,dy\ge \eta_0 .
\]
\end{lemma}

\begin{proof}
Local smooth convergence implies uniform convergence on
\(\overline B_{R_0}\times\{0\}\), hence convergence in
\(L^3(B_{R_0})\).  The map \(f\mapsto\int_{B_{R_0}}|f|^3\) is continuous on
\(L^3(B_{R_0})\), and the claim follows by passing to the limit.
\end{proof}

\subsection{Lower semicontinuity of size}

\begin{lemma}[Lower semicontinuity of \(L^\infty\) size]
\label{p2:lem:lsc-size}
If \(V_n\to V\) locally smoothly and
\[
   \sup_n\|V_n\|_{L^\infty(\R^3\times\R)}<\infty,
\]
then
\[
   \|V\|_{L^\infty(\R^3\times\R)}
   \le
   \liminf_{n\to\infty}
   \|V_n\|_{L^\infty(\R^3\times\R)}.
\]
\end{lemma}

\begin{proof}
Let \(L_0\) be the liminf on the right.  If
\(|V(y_0,\tau_0)|>L_0+\eps\), choose a subsequence with
\(\|V_n\|_{L^\infty}\to L_0\).  Local uniform convergence at
\((y_0,\tau_0)\) gives
\(|V_n(y_0,\tau_0)|>L_0+\eps/2\) for all large \(n\), contradicting
\(\|V_n\|_{L^\infty}\to L_0\).  Hence \(|V|\le L_0\) pointwise.  Since the
smooth representative has pointwise supremum equal to its essential
\(L^\infty\) norm, the result follows.
\end{proof}

\subsection{Existence of a minimal element}

\begin{theorem}[Existence of a minimal ancient solution]
\label{p2:thm:minimal-existence}
Let \(A\) be nonempty and closed normalized.  Set
\[
   m=
   \inf_{V\in A}
   \|V\|_{L^\infty(\R^3\times\R)}.
\]
Then \(0<m<\infty\), and there exists \(V_*\in A\) such that
\[
   \|V_*\|_{L^\infty(\R^3\times\R)}=m.
\]
In particular \(V_*\not\equiv0\).
\end{theorem}

\begin{proof}
The common \(L^\infty\) bound gives \(m<\infty\).  The local nonvanishing gives
a positive lower bound:
\[
   \eta_0
   \le
   \sup_{\tau\in\R}\int_{B_{R_0}}|V(y,\tau)|^3\,dy
   \le
   |B_{R_0}|\|V\|_{L^\infty(\R^3\times\R)}^3,
\]
so
\[
   m\ge \left(\eta_0/|B_{R_0}|\right)^{1/3}>0.
\]

Choose \(V_n\in A\) with
\[
   m\le \|V_n\|_{L^\infty(\R^3\times\R)}\le m+\frac1n .
\]
By \Cref{p2:lem:time-centering}, after replacing \(V_n\) by a time translate,
which stays in \(A\) and preserves the \(L^\infty\) norm, we may assume
\[
   \int_{B_{R_0}}|V_n(y,0)|^3\,dy\ge \eta_0-\frac1n .
\]
The common \(L^\infty\) bound and \Cref{p2:prop:bounded-compactness} give a
locally smooth subsequential limit \(V_*\), which is a bounded mild ancient
solution.  By \Cref{p2:lem:nonzero-stability},
\[
   \int_{B_{R_0}}|V_*(y,0)|^3\,dy\ge\eta_0,
\]
so \(V_*\) retains the fixed \((R_0,\eta_0)\)-normalization.  Closedness of
\(A\) then gives \(V_*\in A\).  Finally, \Cref{p2:lem:lsc-size} gives
\[
   \|V_*\|_{L^\infty}
   \le
   \liminf_{n\to\infty}\|V_n\|_{L^\infty}
   =m.
\]
Since \(m\) is the infimum over \(A\) and \(V_*\in A\), equality holds.  The
positive lower bound on \(m\) rules out \(V_*\equiv0\).
\end{proof}

\subsection{The uniformly tight normalized class}

\begin{definition}[Uniformly tight normalized collection]
\label{p2:def:tight-normalized-collection}
Fix \(M,L,R_0,\eta_0>0\) and a modulus
\[
   \omega:[1,\infty)\to[0,\infty),
   \qquad \omega(R)\to0.
\]
A collection \(A\) of bounded mild ancient solutions is called a uniformly
tight normalized collection with data \((M,L,\omega,R_0,\eta_0)\) if:
\begin{enumerate}[label=\textup{(\arabic*)}]
\item for every \(V\in A\),
\[
   \|V\|_{L^\infty(\R^3\times\R)}\le M;
\]
\item for every \(V\in A\),
\[
   \sup_{\tau\in\R}\|V(\tau)\|_{L^3(\R^3)}\le L;
\]
\item for every \(V\in A\) and every \(R\ge1\),
\[
   \sup_{\tau\in\R}
   \int_{|y|>R}|V(y,\tau)|^3\,dy\le \omega(R);
\]
\item for every \(V\in A\),
\[
   \sup_{\tau\in\R}
   \int_{B_{R_0}}|V(y,\tau)|^3\,dy\ge\eta_0;
\]
\item \(A\) is invariant under time translation;
\item \(A\) is closed under locally smooth limits retaining item \textup{(4)}.
\end{enumerate}
\end{definition}

\begin{proposition}[Uniformly tight normalized collections are closed normalized]
\label{p2:prop:tight-closed-normalized}
Every uniformly tight normalized collection is a closed normalized collection.
\end{proposition}

\begin{proof}
The \(L^\infty\) bound, local nonvanishing, time-translation invariance, and
closedness clause in \Cref{p2:def:closed-normalized} are part of
\Cref{p2:def:tight-normalized-collection}.  It remains to verify that the
stated \(L^3\) bound and fixed tightness modulus are
preserved under locally smooth limits.  This follows from
\Cref{p2:lem:tight-closure}.
\end{proof}

\begin{proof}[Proof of \Cref{p2:thm:minimal-tight-intro}]
By \Cref{p2:prop:tight-closed-normalized}, the collection is closed normalized.
\Cref{p2:thm:minimal-existence} gives the asserted nonzero minimizer.
\end{proof}

\section{Compact trajectory hulls}
\label{p2:sec:compact-hulls}

The material in \Cref{p2:sec:compact-hulls,p2:sec:measures-covariance} is
auxiliary: it is included to clarify why a purely recurrent compactness
argument does not by itself produce a stationary profile, since the
barycenter carries a covariance defect.  None of these results is used in the
proof of \Cref{p2:thm:tight-liouville} or in the final assembly
\Cref{p1:thm:final-assembly}.

On compact hulls we use the metric induced by the seminorms
\(\|\,\cdot\,\|_{C^m(B_N\times[-N,N])}\), \(m,N\in\N\); this gives the
locally smooth topology a complete separable metric structure on bounded
sets, used by the probability-measure arguments in
\Cref{p2:sec:measures-covariance}.

\subsection{Trajectory hulls}

\begin{definition}[Trajectory hull]\label{p2:def:trajectory-hull}
For a bounded mild ancient solution \(V\), its trajectory hull is
\[
   \calH(V)
   =
   \overline{\{\Theta_sV:\ s\in\R\}}^{C^\infty_{\loc}}.
\]
\end{definition}

\subsection{Compactness of bounded hulls}

\begin{theorem}[Compact trajectory hull]\label{p2:thm:compact-hull}
If \(V\) is a bounded mild ancient solution, then \(\calH(V)\) is compact in
the locally smooth topology and invariant under every time translation
\(\Theta_s\).  Every \(W\in\calH(V)\) is a bounded mild ancient solution with
\[
   \|W\|_{L^\infty(\R^3\times\R)}
   \le
   \|V\|_{L^\infty(\R^3\times\R)}.
\]
\end{theorem}

\begin{proof}
Every translate \(\Theta_sV\) is a bounded mild ancient solution with the same
\(L^\infty\) bound as \(V\).  Hence every sequence of translates has a
locally smoothly convergent subsequence by \Cref{p2:prop:bounded-compactness}.
If \((W_n)\subset\calH(V)\), choose translates \(\Theta_{s_n}V\) with
\(d(W_n,\Theta_{s_n}V)<1/n\).  A subsequence of \(\Theta_{s_n}V\) converges
locally smoothly, and the corresponding subsequence of \(W_n\) has the same
limit.  Thus \(\calH(V)\) is sequentially compact.  Since the locally smooth
topology is metrizable, it is compact.

The \(L^\infty\) estimate and the mild equation pass to elements of the hull
by local smooth convergence and \Cref{p2:lem:mild-limit}.  For invariance, note
that \(\Theta_t\) maps the orbit of \(V\) onto itself.  If \(W_n\to W\)
locally smoothly, then \(\Theta_tW_n\to\Theta_tW\) locally smoothly because
time translation only shifts the compact time interval on which the
seminorms are evaluated.  Hence \(\Theta_t\calH(V)=\calH(V)\).

\end{proof}

\begin{corollary}[Uniform tightness on hulls]\label{p2:cor:tightness-on-hulls}
If \(V\) is uniformly \(L^3\)-tight with bound \(L\) and modulus \(\omega\),
then every \(W\in\calH(V)\) satisfies
\[
   \sup_{\tau\in\R}\|W(\tau)\|_{L^3(\R^3)}\le L
\]
and
\[
   \sup_{\tau\in\R}\int_{|y|>R}|W(y,\tau)|^3\,dy\le\omega(R),
   \qquad R\ge1.
\]
The same estimates hold uniformly on every compact invariant subset
\(K\subset\calH(V)\).
\end{corollary}

\begin{proof}
Every translate of \(V\) has the same \(L^3\) bound and tightness modulus.
Applying \Cref{p2:lem:tight-closure} to any sequence of translates converging
locally smoothly to \(W\in\calH(V)\) gives the displayed estimates for \(W\).
Taking the supremum over \(W\in K\subset\calH(V)\) gives the final assertion.
\end{proof}

\subsection{Compact invariant sets}

\begin{definition}[Compact invariant set]\label{p2:def:compact-invariant}
A nonempty compact subset \(K\) of the locally smooth space of bounded mild
ancient solutions is invariant if
\[
   \Theta_tK=K
   \qquad\text{for every }t\in\R .
\]
It is minimal if it contains no proper nonempty compact invariant subset.
\end{definition}

\begin{lemma}[Continuity of the translation action]\label{p2:lem:action-continuity}
Let \(K\) be a compact invariant set.  For each \(t\in\R\),
\(\Theta_t:K\to K\) is a homeomorphism, and the action
\[
   \R\times K\ni(t,W)\mapsto \Theta_tW\in K
\]
is jointly continuous.
\end{lemma}

\begin{proof}
For fixed \(t\), convergence \(W_n\to W\) locally smoothly implies
\[
   \|\Theta_tW_n-\Theta_tW\|_{C^m(B_R\times[-T,T])}
   =
   \|W_n-W\|_{C^m(B_R\times[-T+t,T+t])}\to0.
\]
Thus \(\Theta_t\) is continuous, and its inverse is \(\Theta_{-t}\).  If
\(t_n\to t\) and \(W_n\to W\), then on any fixed compact cylinder the
difference
\[
   \Theta_{t_n}W_n-\Theta_tW
\]
is bounded by a term involving \(W_n-W\) on a slightly larger compact cylinder
and a term involving the uniform continuity of finitely many derivatives of
\(W\) on that larger cylinder.  Both terms tend to zero, proving joint
continuity.
\end{proof}

\subsection{Minimal compact invariant subsets}

\begin{proposition}[Minimal compact invariant subsets]\label{p2:prop:minimal-subset}
Let \(K\) be a nonempty compact invariant set.  Then \(K\) contains a
nonempty compact minimal invariant subset.
\end{proposition}

\begin{proof}
Let \(\mathcal F\) be the family of nonempty compact invariant subsets of
\(K\), ordered by reverse inclusion.  For a chain \(\{K_\alpha\}\), the
intersection \(\bigcap_\alpha K_\alpha\) is nonempty by compactness and the
finite intersection property, since the chain is totally ordered by ordinary
inclusion.  The intersection is compact and invariant: invariance follows
from \(\Theta_tK_\alpha=K_\alpha\) for all \(\alpha\) and the fact that
\(\Theta_t\) is a homeomorphism.  Thus every chain has an upper bound in the
reverse-inclusion order.  Zorn's lemma gives a maximal element for that order,
equivalently a minimal compact invariant subset under ordinary inclusion.
\end{proof}

\subsection{The zero trajectory}

\begin{lemma}[The zero trajectory]\label{p2:lem:zero-trajectory}
The singleton \(\{0\}\) is compact and invariant.  If a minimal compact
invariant set contains the zero trajectory, then it equals \(\{0\}\).
\end{lemma}

\begin{proof}
The zero solution satisfies \(\Theta_t0=0\) for all \(t\).  Therefore
\(\{0\}\) is a nonempty compact invariant subset of any invariant set
containing zero.  Minimality forces equality.
\end{proof}

\section{Invariant measures and the covariance obstruction}
\label{p2:sec:measures-covariance}

\subsection{Krylov--Bogolyubov averages}

The following is the standard Krylov--Bogolyubov averaging argument
\cite{KryloffBogoliouboff1937,Billingsley1999}.

\begin{theorem}[Krylov--Bogolyubov averages]\label{p2:thm:KB}
Let \(K\) be a compact invariant set and \(W_0\in K\).  For \(T>0\), define
\[
   \mu_T
   =
   \frac1T\int_0^T\delta_{\Theta_sW_0}\,ds .
\]
Then every sequence \(T_n\to\infty\) has a subsequence such that
\[
   \mu_{T_n}\rightharpoonup\mu
\]
weakly as Borel probability measures on \(K\), and every such limit \(\mu\)
is invariant:
\[
   (\Theta_t)_\#\mu=\mu
   \qquad\text{for all }t\in\R .
\]
If \(K\) is minimal, then every invariant Borel probability measure on \(K\)
has support equal to \(K\).
\end{theorem}

\begin{proof}
The orbit map \(s\mapsto\Theta_sW_0\) is continuous by
\Cref{p2:lem:action-continuity}; thus \(\mu_T\) is the push-forward of normalized
Lebesgue measure on \([0,T]\).  Since \(K\) is compact and metrizable, the
space of Borel probability measures on \(K\) is weakly sequentially compact
\cite{Billingsley1999}.  Hence \(\mu_{T_n}\) has a weakly convergent
subsequence.

For \(F\in C(K)\) and \(t>0\),
\[
   \int_K F(\Theta_tW)\,d\mu_T(W)-\int_KF(W)\,d\mu_T(W)
   =
   \frac1T\int_T^{T+t}F(\Theta_sW_0)\,ds
   -
   \frac1T\int_0^tF(\Theta_sW_0)\,ds .
\]
The absolute value is at most \(2t\|F\|_{C(K)}/T\).  Letting
\(T=T_{n_k}\to\infty\) gives invariance for \(t>0\); negative times follow by
applying the positive-time result to \(F\circ\Theta_t\).

If \(K\) is minimal and \(\mu\) is invariant, then \(\supp\mu\) is nonempty,
closed, and invariant.  Indeed, if \(U\) is a neighborhood of
\(\Theta_tW\), then \(\Theta_{-t}U\) is a neighborhood of \(W\), and invariance
shows \(\mu(U)=\mu(\Theta_{-t}U)>0\) whenever \(W\in\supp\mu\).  Hence
\(\Theta_t(\supp\mu)\subset\supp\mu\).  Applying the same argument with
\(-t\) gives the reverse inclusion, so \(\supp\mu\) is invariant.  It is
nonempty because \(\mu(K)=1\), and it is compact because it is closed in
compact \(K\).  Minimality forces \(\supp\mu=K\).
\end{proof}

\subsection{Support on minimal hulls}

\begin{corollary}[Nonzero support]\label{p2:cor:nonzero-support}
Let \(K\) be a minimal compact invariant set with \(K\ne\{0\}\).  Then zero is
not in \(K\), and every invariant probability measure on \(K\) has support
equal to \(K\).
\end{corollary}

\begin{proof}
If zero were in \(K\), \Cref{p2:lem:zero-trajectory} would give \(K=\{0\}\),
contrary to the assumption.  The support conclusion is the minimal-support part
of \Cref{p2:thm:KB}.
\end{proof}

\subsection{Stationary statistical identity}

\begin{theorem}[Stationary statistical form]\label{p2:thm:statistical-identity}
Let \(K\) be a compact invariant set and \(\mu\) an invariant Borel
probability measure on \(K\).  Then for every solenoidal
\(\phi\in C_c^\infty(\R^3;\R^3)\),
\begin{equation}\label{p2:eq:statistical-identity}
\begin{aligned}
0
=
\int_K
\bigg[
   &\int_{\R^3}
   W(y,0)\cdot
   \left(\Delta\phi+\phi+\frac12 y\cdot\nabla\phi\right)\,dy  \\
   &+
   \int_{\R^3}
   W_i(y,0)W_j(y,0)\partial_j\phi_i(y)\,dy
\bigg]\,d\mu(W).
\end{aligned}
\end{equation}
\end{theorem}

\begin{proof}
Define
\[
   F(W)=\int_{\R^3}W(y,0)\cdot\phi(y)\,dy .
\]
Local smooth convergence makes \(F\) continuous on \(K\): if \(W_n\to W\),
then \(W_n(\cdot,0)\to W(\cdot,0)\) uniformly on the compact support of
\(\phi\).  Invariance gives
\[
   \int_K\frac{F(\Theta_hW)-F(W)}{h}\,d\mu(W)=0
   \qquad (0<|h|\le1).
\]
Choose \(R\) with \(\supp\phi\subset B_R\).  On
\(\overline B_R\times[-1,1]\), compactness of \(K\) in the locally smooth
topology gives
\[
   M_\phi
   =
   \sup_{W\in K}
   \|\partial_\tau W\|_{L^\infty(\overline B_R\times[-1,1])}
   <\infty .
\]
For \(0<|h|\le1\),
\[
   \frac{F(\Theta_hW)-F(W)}{h}
   =
   \int_{\R^3}
   \left(\frac1h\int_0^h\partial_\tau W(y,s)\,ds\right)\cdot\phi(y)\,dy,
\]
and the absolute value is bounded by
\(|B_R|\|\phi\|_{L^\infty}M_\phi\), uniformly in \(W\) and \(h\).  For each
fixed \(W\), smoothness gives uniform convergence of the inner average to
\(\partial_\tau W(y,0)\) on \(\supp\phi\).  Dominated convergence on \(K\)
yields
\[
   \int_K\int_{\R^3}\partial_\tau W(y,0)\cdot\phi(y)\,dy\,d\mu(W)=0 .
\]
For each \(W\), the smooth local equation \eqref{p2:eq:centered} at time zero,
tested against solenoidal \(\phi\), gives
\[
\begin{aligned}
\int_{\R^3}\partial_\tau W\cdot\phi\,dy
=
&\int_{\R^3} W\cdot\Delta\phi\,dy
 +\int_{\R^3} W\cdot\phi\,dy
 +\frac12\int_{\R^3}W\cdot(y\cdot\nabla\phi)\,dy\\
&+\int_{\R^3}W_iW_j\partial_j\phi_i\,dy .
\end{aligned}
\]
The pressure term vanishes because \(\divergence\phi=0\); the drift identity
uses
\[
   \int y_j\partial_jW_i\,\phi_i\,dy
   =
   -3\int W\cdot\phi\,dy-\int W\cdot(y\cdot\nabla\phi)\,dy .
\]
The right-hand side is a continuous bounded function of \(W\in K\).  Averaging
over \(K\) gives \eqref{p2:eq:statistical-identity}.
\end{proof}

\subsection{Barycenter and covariance}

\begin{theorem}[Barycenter and covariance]\label{p2:thm:barycenter-covariance}
Let \(K\) and \(\mu\) be as in \Cref{p2:thm:statistical-identity}.  Define
\[
   \overline W(y)=\int_K W(y,0)\,d\mu(W)
\]
and
\[
   C_{ij}(y)
   =
   \int_K W_i(y,0)W_j(y,0)\,d\mu(W)
   -
   \overline W_i(y)\overline W_j(y).
\]
Then \(\overline W\) and \(C\) are \(C^\infty_{\loc}\), \(\overline W\) is
solenoidal, and for every solenoidal
\(\phi\in C_c^\infty(\R^3;\R^3)\),
\begin{equation}\label{p2:eq:barycenter-covariance}
   \int_{\R^3}
   \overline W\cdot
   \left(\Delta\phi+\phi+\frac12 y\cdot\nabla\phi\right)\,dy
   +
   \int_{\R^3}
   \overline W_i\overline W_j\partial_j\phi_i\,dy
   +
   \int_{\R^3}C_{ij}\partial_j\phi_i\,dy
   =0 .
\end{equation}
In particular, \(\overline W\) is a weak stationary self-similar solution
against solenoidal test fields if and only if
\[
   \int_{\R^3}C_{ij}\partial_j\phi_i\,dy=0
\]
for every solenoidal test field \(\phi\).
\end{theorem}

\begin{proof}
Compactness of \(K\) in the locally smooth topology gives, for every compact
ball and every derivative order, a uniform bound on
\(\partial_y^\alpha W(y,0)\) for \(W\in K\).  Hence one may differentiate
\(\overline W\) under the integral sign.  More explicitly, for every
multiindex \(\alpha\) and every \(y\) in a fixed compact ball,
\[
   \partial_y^\alpha\overline W(y)
   =
   \int_K\partial_y^\alpha W(y,0)\,d\mu(W),
\]
because the difference quotients are dominated by the next derivative
seminorm, uniformly on \(K\).  The same argument applies to the second moment
\[
   M_{ij}(y)=\int_KW_i(y,0)W_j(y,0)\,d\mu(W),
\]
using Leibniz' rule and the uniform bounds on derivatives of \(W\).  This
gives \(C^\infty_{\loc}\) regularity of \(\overline W\), \(M\), and
\(C=M-\overline W\otimes\overline W\).  Since every \(W\) is divergence free,
Fubini against scalar test functions gives \(\divergence\overline W=0\):
for \(\eta\in C_c^\infty(\R^3)\),
\[
   \int_{\R^3}\overline W\cdot\nabla\eta\,dy
   =
   \int_K\int_{\R^3}W(y,0)\cdot\nabla\eta(y)\,dy\,d\mu(W)=0.
\]

For each compactly supported test field \(\phi\), the maps
\[
   W\mapsto\int_{\R^3}W(y,0)\cdot\phi(y)\,dy,
   \qquad
   W\mapsto\int_{\R^3}W_i(y,0)W_j(y,0)\partial_j\phi_i(y)\,dy
\]
are continuous and bounded on \(K\), again by compactness in
\(C^\infty_{\loc}\).  Thus all integrals below are finite, and Fubini's theorem
applies.

Applying \Cref{p2:thm:statistical-identity} and passing the linear term through
the \(K\)-integral gives the first integral in
\eqref{p2:eq:barycenter-covariance}.  Passing the nonlinear term through the
integral gives
\[
   \int_K W_iW_j\,d\mu
   =
   \overline W_i\overline W_j+C_{ij}.
\]
Substitution yields \eqref{p2:eq:barycenter-covariance}.  Comparing this identity
with the pressure-free stationary weak formulation for \(\overline W\) gives
the stated equivalence with vanishing covariance contribution.
\end{proof}

\section{\texorpdfstring{Endpoint \(L^3\) closure of the uniformly tight class}{Endpoint L3 closure of the uniformly tight class}}
\label{p2:sec:endpoint-L3-closure}

\subsection{Physical pullback}

The pullback below may be unbounded as \(t\uparrow0\); it is bounded only on
every truncated ancient interval \((-\infty,T]\) with \(T<0\) by the factor
\((-T)^{-1/2}\) below.  For this reason the Albritton--Barker endpoint Liouville
theorem is applied after translating such a truncated interval to terminal
time \(0\), not directly to the unbounded pullback on \((-\infty,0)\).

Let \(V\) be a smooth centered ancient velocity and let \(P\) be a compatible
local pressure.  For \(t<0\) define
\begin{equation}\label{p2:eq:physical-pullback}
   u(x,t)=(-t)^{-1/2}
   V\left(\frac{x}{\sqrt{-t}},-\log(-t)\right),
\end{equation}
and
\begin{equation}\label{p2:eq:physical-pressure-pullback}
   p(x,t)=(-t)^{-1}
   P\left(\frac{x}{\sqrt{-t}},-\log(-t)\right).
\end{equation}
Equivalently,
\[
   V(y,\tau)=e^{-\tau/2}u(e^{-\tau/2}y,-e^{-\tau}),
   \qquad
   P(y,\tau)=e^{-\tau}p(e^{-\tau/2}y,-e^{-\tau}).
\]

\begin{lemma}[Pullback solves physical Navier--Stokes]
\label{p2:lem:physical-pullback}
Let \(V\) be a bounded mild ancient solution of \eqref{p2:eq:centered}.  Then
the pullback \(u\) in \eqref{p2:eq:physical-pullback} is a smooth divergence-free
solution of
\begin{equation}\label{p2:eq:physical-NS}
   \partial_tu+(u\cdot\nabla)u+\nabla p=\Delta u,
   \qquad \divergence u=0,
\end{equation}
on \(\R^3\times(-\infty,0)\), with pressure \(p\) locally given by
\eqref{p2:eq:physical-pressure-pullback}.  Moreover, for every \(T<0\),
\[
   \|u\|_{L^\infty(\R^3\times(-\infty,T])}
   \le
   (-T)^{-1/2}\|V\|_{L^\infty(\R^3\times\R)}.
\]
\end{lemma}

\begin{proof}
By \Cref{p2:lem:mild-smoothing}, \(V\) is smooth.  By
\Cref{p2:lem:local-pressure}, on every compact parabolic cylinder there is a
smooth pressure representative \(P\), unique up to a function of \(\tau\), for
which \eqref{p2:eq:centered} holds classically.  Adding a function of \(\tau\) to
\(P\) only adds a function of \(t\) to \(p\), and hence does not affect
\(\nabla p\).

Set
\[
   \lambda(t)=(-t)^{-1/2},\qquad
   y=\lambda(t)x,\qquad
   \tau=-\log(-t).
\]
Then \(\dot \tau=\lambda^2\) and
\[
   \partial_t y=\frac12\lambda^2 y,
   \qquad
   \dot\lambda=\frac12\lambda^3 .
\]
Using \(u=\lambda V(y,\tau)\), one obtains
\[
   \partial_tu
   =
   \lambda^3
   \left(
      \partial_\tau V+\frac12y\cdot\nabla V+\frac12V
   \right),
\]
while
\[
   (u\cdot\nabla_x)u=\lambda^3(V\cdot\nabla_y)V,
   \qquad
   \Delta_xu=\lambda^3\Delta_yV,
   \qquad
   \nabla_xp=\lambda^3\nabla_yP .
\]
Substituting these identities into \eqref{p2:eq:centered} gives
\eqref{p2:eq:physical-NS}.  The divergence identity follows similarly:
\[
   \divergence_xu=\lambda^2\divergence_yV=0.
\]
Finally, if \(t\le T<0\), then \((-t)^{-1/2}\le(-T)^{-1/2}\), and the stated
\(L^\infty\) bound follows directly from \eqref{p2:eq:physical-pullback}.
\end{proof}

\subsection{\texorpdfstring{Critical \(L^3\) invariance}{Critical L3 invariance}}

\begin{lemma}[Invariance of the critical norm]
\label{p2:lem:L3-pullback-invariance}
Let \(u\) be the physical pullback of \(V\).  Then, for every \(t<0\),
\begin{equation}\label{p2:eq:L3-pullback-invariance}
   \|u(t)\|_{L^3(\R^3)}
   =
   \|V(-\log(-t))\|_{L^3(\R^3)} .
\end{equation}
\end{lemma}

\begin{proof}
With \(y=x/\sqrt{-t}\), \eqref{p2:eq:physical-pullback} gives
\[
   |u(x,t)|^3=(-t)^{-3/2}|V(y,-\log(-t))|^3,
   \qquad
   dx=(-t)^{3/2}dy .
\]
Multiplying these two identities and integrating over \(\R^3\) gives
\eqref{p2:eq:L3-pullback-invariance}.
\end{proof}

\begin{lemma}[Tail tightness gives an endpoint sequence]
\label{p2:lem:tail-tight-sequence-L3}
Let \(V\) be a bounded function on \(\R^3\times\R\) satisfying
\[
   \lim_{R\to\infty}
   \sup_{\tau\in\R}\int_{|y|>R}|V(y,\tau)|^3\,dy=0 .
\]
Then, for every sequence \(\tau_k\in\R\),
\[
   \sup_k\|V(\tau_k)\|_{L^3(\R^3)}<\infty .
\]
\end{lemma}

\begin{proof}
Let \(M=\|V\|_{L^\infty(\R^3\times\R)}\).  Choose \(R\ge1\) such that
\[
   \sup_{\tau\in\R}\int_{|y|>R}|V(y,\tau)|^3\,dy\le1 .
\]
Then, for every \(k\),
\[
   \int_{\R^3}|V(y,\tau_k)|^3\,dy
   \le M^3 |B_R|
      + \sup_{\tau\in\R}\int_{|y|>R}|V(y,\tau)|^3\,dy
   \le M^3|B_R|+1 .
\]
This gives the claimed endpoint sequence bound.
\end{proof}

\subsection{Mildness after finite terminal-time restriction}

\begin{lemma}[Physical pullbacks are mild ancient solutions on finite intervals]
\label{p2:lem:physical-mild}
Let \(V\) be a bounded mild ancient solution of \eqref{p2:eq:centered}, and let
\(u\) be its physical pullback.  For every \(T<0\), the restriction of \(u\)
to \(\R^3\times(-\infty,T]\) is a bounded mild ancient solution of the
physical Navier--Stokes equations after translating the terminal time \(T\) to
\(0\).
\end{lemma}

\begin{proof}
Let \(s<t\le T<0\), and set
\[
   \sigma=-\log(-s),\qquad \tau=-\log(-t).
\]
Then \(\sigma<\tau\).  Apply the centered mild formula for \(V\) on
\([\sigma,\tau]\):
\[
   V(\tau)
   =
   S(\tau-\sigma)V(\sigma)
   -
   \int_\sigma^\tau
   S(\tau-\alpha)\mathbb P\nabla\cdot(V\otimes V)(\alpha)\,d\alpha .
\]
Put
\[
   a=\tau-\sigma,\qquad \lambda=e^{-a/2}=\sqrt{\frac{-t}{-s}},
   \qquad \theta=1-e^{-a}=\frac{t-s}{-s}.
\]
For the linear term, use the explicit Ornstein--Uhlenbeck semigroup formula
\eqref{p2:eq:OU-semigroup}.  Writing \(x=\sqrt{-t}\,y\) and changing
variables \(z=\zeta/\sqrt{-s}\), we obtain
\[
\begin{aligned}
(-t)^{-1/2}(S(a)V(\sigma))(y)
&=(-t)^{-1/2}e^{-a/2}
  \int_{\R^3}
  \frac{\exp\!\left(-\frac{|e^{-a/2}y-z|^2}{4(1-e^{-a})}\right)}
       {[4\pi(1-e^{-a})]^{3/2}}
  V(z,\sigma)\,dz \\
&=\int_{\R^3}
  \frac{\exp\!\left(-\frac{|x-\zeta|^2}{4(t-s)}\right)}
       {[4\pi(t-s)]^{3/2}}
  u(\zeta,s)\,d\zeta \\
&=(e^{(t-s)\Delta}u(s))(x).
\end{aligned}
\]
For the Duhamel term, write \(r=-e^{-\alpha}\) for
\(\alpha\in[\sigma,\tau]\).  Then \(dr=(-r)\,d\alpha\),
\[
   V(\xi,\alpha)\otimes V(\xi,\alpha)
   =(-r)\,(u\otimes u)(\sqrt{-r}\,\xi,r),
\]
and the homogeneity of \(\mathbb P\nabla\cdot\) together with the same
scaling calculation gives
\[
\begin{aligned}
&(-t)^{-1/2}
  \bigl(S(\tau-\alpha)\mathbb P\nabla\cdot(V\otimes V)(\alpha)\bigr)(y) \\
&\qquad =
  (-r)\,
  \bigl(e^{(t-r)\Delta}\mathbb P\nabla\cdot(u\otimes u)(r)\bigr)(x).
\end{aligned}
\]
Indeed, if \(F_\alpha(\xi)=V(\xi,\alpha)\otimes V(\xi,\alpha)\) and
\(H_r(x)=u(x,r)\otimes u(x,r)\), then
\[
   F_\alpha(\xi)=(-r)H_r(\sqrt{-r}\,\xi),
\]
so
\[
   \mathbb P\nabla_\xi\cdot F_\alpha
   =(-r)^{3/2}
   \bigl(\mathbb P\nabla_x\cdot H_r\bigr)(\sqrt{-r}\,\xi),
\]
and using \(S(b)f(\eta)=e^{-b/2}(e^{(1-e^{-b})\Delta}f)(e^{-b/2}\eta)\)
with \(b=\tau-\alpha\) yields the displayed identity because
\(e^{-b/2}\sqrt{-r}=\sqrt{-t}\) and \((-r)(1-e^{-b})=t-r\).
Therefore, after the change of variables \(r=-e^{-\alpha}\),
\[
\begin{aligned}
&(-t)^{-1/2}
  \int_\sigma^\tau
  S(\tau-\alpha)\mathbb P\nabla\cdot(V\otimes V)(\alpha)\,d\alpha \\
&\qquad =
  \int_s^t
  e^{(t-r)\Delta}\mathbb P\nabla\cdot(u\otimes u)(r)\,dr .
\end{aligned}
\]
Under the parabolic change of variables
\[
   u(x,t)=(-t)^{-1/2}V(x/\sqrt{-t},-\log(-t)),
\]
the centered mild formula therefore becomes
\[
   u(t)
   =
   e^{(t-s)\Delta}u(s)
   -
   \int_s^t
   e^{(t-r)\Delta}\mathbb P\divergence(u\otimes u)(r)\,dr
\]
in the weak-* topology of \(L^\infty(\R^3)\).  Since \(t\le T<0\), the
physical pullback is bounded on \(\R^3\times(-\infty,T]\), and the Oseen
kernel estimate makes the Duhamel term finite on every finite interval.  Hence
the restriction of \(u\) to \((-\infty,T]\), after translating \(T\) to \(0\),
is a bounded mild ancient solution.
\end{proof}

\subsection{\texorpdfstring{Endpoint \(L^3\) ancient Liouville theorem}{Endpoint L3 ancient Liouville theorem}}

\begin{theorem}[Albritton--Barker endpoint ancient Liouville theorem]
\label{p2:thm:AB-endpoint}
Let \(u\) be a bounded mild ancient solution of the three-dimensional
Navier--Stokes equations on \(\R^3\times(-\infty,0]\).  Suppose that there
is a sequence \(s_k\downarrow-\infty\) such that
\[
   \sup_k\|u(s_k)\|_{L^3(\R^3)}<\infty,
\]
and suppose \(u(0)\in L^3(\R^3)\).  Then \(u\equiv0\).  The statement is
written after translating the terminal time to \(0\).
\end{theorem}

\begin{proof}
We use the \(L^3\)-sequence formulation of Albritton--Barker's ancient
Liouville theorem \cite[Theorem 1.2]{AlbrittonBarker2019}, after
translating the terminal time to \(0\).
\end{proof}

\begin{remark}[Quantitative Lorentz/Besov form]
\label{p2:rem:AB-quantitative}
Albritton--Barker also give a quantitative endpoint statement
\cite[Theorem 4.1]{AlbrittonBarker2019} with additional quantitative
endpoint hypotheses.  Those hypotheses are not used in this paper.  The only
imported statement needed below is the \(L^3\)-sequence Liouville theorem in
\Cref{p2:thm:AB-endpoint}.
\end{remark}

\begin{remark}[Notation match with Albritton--Barker]
\label{p2:rem:AB-notation}
In the tight-class application below, the terminal datum is \(u_T(0)=u(T)\).
Tail tightness and the critical \(L^3\) identity give \(u(T)\in L^3(\R^3)\),
so the \(L^3\) terminal hypothesis of \Cref{p2:thm:AB-endpoint} is
automatic.  The backward sequence bound is obtained in \(L^3\), which is
the form used below.
\end{remark}

\begin{theorem}[Tail-tight centered ancient solutions vanish]
\label{p2:thm:tight-liouville}
Let \(V\) be a bounded mild ancient solution of \eqref{p2:eq:centered}.  Assume
\[
   \lim_{R\to\infty}
   \sup_{\tau\in\R}\int_{|y|>R}|V(y,\tau)|^3\,dy=0 .
\]
Then \(V\equiv0\).
\end{theorem}

\begin{proof}
Let \(u\) be the physical pullback of \(V\).  Fix \(T<0\).  By
\Cref{p2:lem:physical-mild}, after translating the terminal time \(T\) to \(0\),
the restriction of \(u\) to \((-\infty,T]\) is a bounded mild ancient
solution.  Moreover \(u(T)\in L^3(\R^3)\): indeed, if
\(\tau_T=-\log(-T)\), then \(\|u(T)\|_{L^3}=\|V(\tau_T)\|_{L^3}\), and the
latter is finite by boundedness on \(B_R\) and tail tightness outside \(B_R\).
Choose \(t_k=-e^k\).  For all sufficiently large \(k\), \(t_k<T\),
and \(\tau_k=-\log(-t_k)=-k\).  By
\Cref{p2:lem:tail-tight-sequence-L3},
\[
   \sup_k\|V(-k)\|_{L^3(\R^3)}<\infty .
\]
Hence \eqref{p2:eq:L3-pullback-invariance} gives
\[
   \|u(t_k)\|_{L^3(\R^3)}
   =
   \|V(-k)\|_{L^3(\R^3)}
\]
uniformly along the same sequence, after discarding the finitely many indices
for which \(t_k\ge T\).
Set \(u_T(x,\sigma)=u(x,\sigma+T)\) for \(\sigma\le0\) and
\(s_k=t_k-T\).  Then \(s_k\downarrow-\infty\) and
\[
   \|u_T(s_k)\|_{L^3(\R^3)}
   =
   \|u(t_k)\|_{L^3(\R^3)}
\]
is uniformly bounded.  Applying \Cref{p2:thm:AB-endpoint} to \(u_T\) gives
\(u_T\equiv0\) on \(\R^3\times(-\infty,0]\), equivalently \(u=0\) on
\(\R^3\times(-\infty,T]\).  Since \(T<0\) was arbitrary, \(u=0\) on
\(\R^3\times(-\infty,0)\).  The inverse change of variables then gives
\[
   V(y,\tau)=e^{-\tau/2}u(e^{-\tau/2}y,-e^{-\tau})=0
\]
for every \((y,\tau)\).
\end{proof}

\section{Tight-Class Liouville Criterion and Type I Reduction}
\label{p2:sec:tight-liouville-interface}

\subsection{Normalized tight collections}

\begin{theorem}[No nonempty normalized uniformly tight collection]
\label{p2:thm:no-tight-normalized-collection}
No uniformly tight normalized collection in the sense of
\Cref{p2:def:tight-normalized-collection} is nonempty.
\end{theorem}

\begin{proof}
Suppose \(A\) is such a collection and choose \(V\in A\).  By the tightness
modulus in item \textup{(3)} of \Cref{p2:def:tight-normalized-collection}, \(V\)
satisfies the tail condition in \Cref{p2:thm:tight-liouville}.  Hence
\Cref{p2:thm:tight-liouville} gives \(V\equiv0\).  This contradicts the local
velocity normalization in the same definition:
\[
   \sup_{\tau\in\R}
   \int_{B_{R_0}}|V(y,\tau)|^3\,dy\ge\eta_0>0 .
\]
Hence \(A\) is empty.
\end{proof}

\begin{proof}[Proof of \Cref{p2:thm:tight-liouville-intro}]
If \(A\) were nonempty, choose \(V\in A\).  The common tail modulus gives the
tail condition in \Cref{p2:thm:tight-liouville}, so \(V\equiv0\).  This contradicts
the fixed local nonvanishing normalization in \Cref{p2:thm:tight-liouville-intro}.
\end{proof}

\subsection{Use in the Type I reduction}

\begin{definition}[Tight-class Liouville consequence]
\label{p2:def:tight-class-input}
The tight-class Liouville consequence is the assertion that no generated raw
Seregin state satisfying the invariant nonvanishing normalization of
\Cref{p1:cor:normalized-nonzero} belongs to \(\calC_{L^3\mathrm{-tight}}\).
\end{definition}

\begin{corollary}[Tight-class consequence for the Type I reduction]
\label{p2:cor:tight-input-for-paper-I}
\Cref{p2:def:tight-class-input} holds.
\end{corollary}

\begin{proof}
Suppose, contrary to \Cref{p2:def:tight-class-input}, that a generated raw
Seregin state \(V\) belongs to \(\calC_{L^3\mathrm{-tight}}\).  By definition
of \(\calC_{L^3\mathrm{-tight}}\), \(V\in\calS_{\mathrm{mild}}\) and satisfies
the tail condition
\[
   \lim_{R\to\infty}
   \sup_{\tau\in\R}\int_{|y|>R}|V(y,\tau)|^3\,dy=0 .
\]
\Cref{p2:thm:tight-liouville} applies and gives \(V\equiv0\), contradicting
\Cref{p1:cor:normalized-nonzero}.  Hence no generated raw Seregin state can
belong to \(\calC_{L^3\mathrm{-tight}}\).
\end{proof}

The tight-class consequence used in the Type I reduction is therefore proved
without the residual-class hypothesis.  The proof relies only on boundedness, the
tail-tightness condition of the uniformly \(L^3\)-tight class, the
local-in-time physical pullback, and the endpoint ancient Liouville theorem of
Albritton--Barker.



\clearpage
\section*{Part IV. Structured Ancient Solution Exclusions}

\subsection*{Overview}

The final part concerns bounded ancient solutions of the three-dimensional
Navier--Stokes equations arising from Type I blow-up arguments.  Explicit PDE
hypotheses are stated under which such solutions are spatially constant; in the
Type I class, these constants vanish.  The perturbative weighted-vorticity
anchored Gallay--Wayne perturbative exclusion used by the conditional Type I
reduction is also proved.

For the unforced equations, the cited axisymmetric Liouville theorems cover the
following cases: absence of swirl, pointwise scale-invariant control, finite
\(L_t^\infty L_x^p\) control of the swirl variable \(\Gamma=ru_\theta\), and
axial periodicity with bounded \(\Gamma\).  When one of these hypotheses holds
for a nonzero Type I blow-up limit, the cited theorems reduce the solution to a
spatially constant state in the physical Type I frame, and that constant is
then forced to vanish, contradicting the nontriviality of the limit.

The perturbative part uses Gallay--Wayne's forward small-data theorem only
through anchored forward rescalings launched from arbitrarily early physical
times.  The anchored small-data estimate forces the vorticity at any fixed
physical time to vanish, and the section concludes with a summary of the
Liouville hypotheses and cited PDE results.

\section{Introduction}

\subsection{Ancient solutions from Type I blow-up}

The Type I blow-up argument in \Cref{p1:thm:extraction} produces, after
parabolic rescaling around a Type I singular point, a nonzero ancient solution
\((U,P)\) of
\begin{equation}\label{p3:eq:NS}
   \partial_tU+(U\cdot\nabla)U+\nabla P=\Delta U,
   \qquad \nabla\cdot U=0,
   \qquad t<0,
\end{equation}
satisfying the Type I ancient bound
\begin{equation}\label{p3:eq:typeI}
   \sup_{t<0}\sqrt{-t}\,
   \|U(t)\|_{L^\infty(\mathbb R^3)}<\infty .
\end{equation}
Thus every result proving \(U\equiv0\) for an ancient solution satisfying
\eqref{p3:eq:typeI} has the same consequence in the Type I problem:
\[
   \text{PDE hypothesis}
   \quad\Longrightarrow\quad
   U\equiv0
   \quad\Longrightarrow\quad
   \text{contradiction with nonzero blow-up limit}.
\]

The hypotheses treated in this part are precisely those to which a cited
Liouville theorem or the perturbative argument in
\Cref{p3:sec:weighted-vorticity} applies:
\[
   \begin{gathered}
   \text{axisymmetric Liouville hypotheses},\qquad
   \text{anchored Gallay--Wayne perturbative vorticity}.
   \end{gathered}
\]
No conclusion is asserted for decaying or symmetry-restricted ancient solutions
outside the stated hypotheses.  Each conclusion is tied to an explicit
PDE hypothesis and a specified analytic result.

For the Type I application in \Cref{p1:thm:extraction}, the exclusions used are
precisely the cases in which the stated PDE hypothesis and the cited or proved
Liouville theorem imply \(U\equiv0\).  Ancient solutions outside these
hypotheses are treated separately.

\subsection{Axisymmetric Liouville hypotheses}

Several unforced axisymmetric cases fall under cited Liouville theorems.
Koch--Nadirashvili--Seregin--\v{S}ver\'ak treat the no-swirl case
and the pointwise scale-invariant condition \(|u|\le C/r\).  The
swirl-bearing finite-\(L^p\) condition
\[
   \Gamma=ru_\theta\in L_t^\infty L_x^p,\qquad 1\le p<\infty,
\]
is treated by Lei--Zhang--Zhao.  The \(z\)-periodic case with bounded
\(\Gamma\) is treated by Lei--Ren--Zhang.

For Type I ancient limits, a constant conclusion after the relevant Galilean
reduction must be converted into zero in the blow-up frame.  The proof
time-shifts away from the terminal time \(t=0\), applies the bounded ancient
Liouville theorem, patches the resulting constants on backward intervals, and
uses
\eqref{p3:eq:typeI}: if \(U(x,t)\equiv c\), then
\[
   \sup_{t<0}\sqrt{-t}\,\|U(t)\|_\infty
   =
   \sup_{t<0}\sqrt{-t}\,|c|
\]
is finite only when \(c=0\).

\subsection{Anchored Gallay--Wayne hypotheses}

The perturbative branch is stated directly for the physical Type I ancient
representative \(U\).  From anchor times \(a_n\to-\infty\) and scales
\(\lambda_n\), one launches the Gallay--Wayne forward rescaled vorticity
equation from the anchored physical vorticity data
\[
   w_n^0(\xi)=\lambda_n^2(\nabla\times U)(\lambda_n\xi,a_n).
\]
If these anchored initial data remain in a fixed perturbative
\(L^2_\sigma(m)\)-ball, with the Biot--Savart normalization at each anchor,
then Gallay--Wayne's forward decay estimate applies on every anchored forward
interval.  Evaluating those anchored forward solutions at any fixed physical
time \(T<0\) and sending \(n\to\infty\) forces \(\nabla\times U(\cdot,T)=0\),
hence \(U\equiv0\).  This avoids treating the backward Type I centered
vorticity as a single Gallay--Wayne forward orbit.

\subsection{Main results}

The first result states the unforced axisymmetric Liouville consequences
relevant to the Type I argument.

\begin{theorem}[Axisymmetric Liouville consequences for Type I ancient solutions]
\label{p3:thm:axisymmetric-liouville}
Let \(U\) be a smooth mild Type I ancient solution of \eqref{p3:eq:NS}
in the sense of \Cref{p3:def:typeI-ancient}.  If \(U\) satisfies any of the following axisymmetric
assumptions, then \(U\equiv0\):
\begin{enumerate}[label=\textup{(\roman*)}]
\item axisymmetric without swirl;
\item axisymmetric and satisfying a pointwise scale-invariant bound
\[
   |U(x,t)|\le \frac{C}{r};
\]
\item axisymmetric with
\[
   \Gamma=rU_\theta\in L_t^\infty L_x^p
   \bigl(\mathbb R^3\times(-\infty,0)\bigr),
   \qquad 1\le p<\infty;
\]
\item axisymmetric, periodic in the physical axial variable \(z\), and with
\[
   \Gamma=rU_\theta\in L^\infty
   \bigl(\mathbb R^3\times(-\infty,0)\bigr).
\]
\end{enumerate}
Consequently none of these assumptions can hold for a nonzero Type I blow-up
limit.
\end{theorem}

\begin{theorem}[Anchored Gallay--Wayne perturbative exclusion]
\label{p3:thm:weighted-vorticity-vanishing-intro}
Let \(U\) be a smooth mild Type I ancient solution on
\(\mathbb R^3\times(-\infty,0)\).  If \(U\) satisfies the anchored
Gallay--Wayne perturbative vorticity condition of
\Cref{p3:def:gw-anchored-condition}, then
\[
   U\equiv0.
\]
\end{theorem}

The anchored condition is stated directly in terms of the physical ancient
representative \(U\).  Gallay--Wayne's theorem is used only on forward Cauchy
problems launched from the anchor times \(a_n\).

\begin{theorem}[Consequences for Type I blow-up limits]
\label{p3:thm:typeI-consequences}
Let \(U\) be a nonzero Type I blow-up limit.  Then \(U\) cannot satisfy any of
the axisymmetric assumptions in \Cref{p3:thm:axisymmetric-liouville}; nor can
\(U\) satisfy the anchored Gallay--Wayne perturbative vorticity condition of
\Cref{p3:def:gw-anchored-condition}.
\end{theorem}

Only the explicit hypotheses in the theorem are used.

\subsection{Relation to the preceding parts}

\Cref{p1:thm:extraction} reduces Type I singularities to nonzero ancient
limits satisfying \eqref{p3:eq:typeI}.  \Cref{p2:thm:tight-liouville} treats
uniformly \(L^3\)-tight ancient solutions.  The axisymmetric and perturbative
weighted-vorticity exclusions used by the Type I argument are proved in
\Cref{p3:sec:axisymmetric-consequences,p3:sec:weighted-vorticity}.  Cited
Liouville theorems rule out the axisymmetric cases, and the anchored
Gallay--Wayne argument handles the perturbative weighted-vorticity case.
\Cref{p3:thm:typeI-zero-hypotheses} does not depend on
\Cref{p2:thm:tight-liouville}.

\subsection{Summary of hypotheses}

\begin{center}
\begin{tabular}{p{0.36\textwidth}p{0.18\textwidth}p{0.34\textwidth}}
\hline
Hypothesis & Conclusion & Analytic ingredient \\
\hline
axisymmetric no swirl & \(U\equiv0\) & KNSS plus Type I constant removal \\
pointwise \(|u|\le C/r\) & \(U\equiv0\) & KNSS plus \(r\to\infty\) constant removal \\
\(\Gamma\in L_t^\infty L_x^p\), \(1\le p<\infty\) & \(U\equiv0\) & KNSS/LZZ split plus Type I constant removal \\
\(z\)-periodic bounded \(\Gamma\) & \(U\equiv0\) & LRZ plus Type I constant removal \\
anchored Gallay--Wayne perturbative vorticity & \(U\equiv0\) & anchored Gallay--Wayne decay plus physical uniqueness comparison \\
\hline
\end{tabular}
\end{center}

\subsection{Organization}

\Cref{p3:sec:typeI} introduces the Type I ancient setting, self-similar variables, the
invariance facts, and the time-shift lemma.
\Cref{p3:sec:axisymmetric-liouville-theorems} fixes axisymmetric notation and the
cited Liouville theorems.  \Cref{p3:sec:axisymmetric-consequences} proves the
unforced axisymmetric Liouville consequences.
\Cref{p3:sec:weighted-vorticity} proves the anchored Gallay--Wayne perturbative exclusion, and
\Cref{p3:sec:summary} summarizes the consequences.

\section{Ancient Type I limits and self-similar variables}
\label{p3:sec:typeI}

\subsection{Type I ancient solutions}

\begin{definition}[Type I ancient solution]
\label{p3:def:typeI-ancient}
A smooth mild solution \(U\) of \eqref{p3:eq:NS} on
\(\mathbb R^3\times(-\infty,0)\) is called Type I ancient if
\eqref{p3:eq:typeI} holds.  Mild is understood in the standard velocity
formulation \cite{Kato1984,Sohr2001}: for every \(-\infty<s<t<0\),
\begin{equation}\label{p3:eq:mild-formula}
   U(t)=e^{(t-s)\Delta}U(s)
   -\int_s^t e^{(t-\sigma)\Delta}\mathbb P\nabla\cdot
      (U(\sigma)\otimes U(\sigma))\,d\sigma
\end{equation}
in distributions, where \(\mathbb P\) is the Leray projector.
\end{definition}

\subsection{Nonzero Type I blow-up limits}

\begin{definition}[Nonzero Type I blow-up limit]
\label{p3:def:nonzero-typeI-limit}
A Type I ancient solution \(U\) is called a nonzero Type I blow-up limit if it
arises as the limit of parabolic rescalings around a Type I singular point and
the limiting solution is not identically zero.  The arguments below require
only the conclusion \(U\not\equiv0\).
\end{definition}

\begin{lemma}[Nontriviality of blow-up limits]
\label{p3:lem:nonzero-typeI-limit}
   A nonzero Type I blow-up limit is not identically zero.
\end{lemma}

\begin{proof}
The definition includes the condition \(U\not\equiv0\).
\end{proof}

\subsection{Backward self-similar variables}

For \(t<0\), define
\begin{equation}\label{p3:eq:self-similar}
   V(y,\tau)=\sqrt{-t}\,U(x,t),\qquad
   Q(y,\tau)=(-t)P(x,t),\qquad
   y=\frac{x}{\sqrt{-t}},\qquad
   \tau=-\log(-t).
\end{equation}
Equivalently,
\[
   x=e^{-\tau/2}y,\qquad t=-e^{-\tau}.
\]

\begin{lemma}[Self-similar equation]\label{p3:lem:self-similar-equation}
Let \((U,P)\) be a smooth solution of \eqref{p3:eq:NS}, and define \((V,Q)\) by
\eqref{p3:eq:self-similar}.  Then \(V\) is defined on
\(\mathbb R^3\times\mathbb R\) and satisfies
\begin{equation}\label{p3:eq:centered}
   \partial_\tau V+(V\cdot\nabla)V+\nabla Q
   =
   \Delta V-\frac12(V+y\cdot\nabla V),
   \qquad \nabla\cdot V=0.
\end{equation}
Moreover \eqref{p3:eq:typeI} is equivalent to
\[
   \sup_{\tau\in\mathbb R}
   \|V(\tau)\|_{L^\infty(\mathbb R^3)}<\infty .
\]
\end{lemma}

\begin{proof}
Let \(a=\sqrt{-t}=e^{-\tau/2}\).  Then \(V=aU\), \(y=x/a\), and
\[
   \partial_t\tau=\frac1{-t}=\frac1{a^2},
   \qquad
   \partial_t y=\frac{y}{2a^2},
   \qquad
   \partial_t(a^{-1})=\frac1{2a^3}.
\]
Since \(U=a^{-1}V\), differentiating at fixed \(x\) gives
\[
   \partial_tU
   =
   a^{-3}
   \left(
      \partial_\tau V+\frac12 y\cdot\nabla V+\frac12V
   \right).
\]
Spatial derivatives scale as
\[
   \nabla_xU=a^{-2}\nabla_yV,\qquad
   \Delta_xU=a^{-3}\Delta_yV,
\]
and the pressure scaling gives
\[
   \nabla_xP=a^{-3}\nabla_yQ.
\]
Thus
\[
   (U\cdot\nabla_x)U=a^{-3}(V\cdot\nabla_y)V.
\]
Multiplying \eqref{p3:eq:NS} by \(a^3\) yields \eqref{p3:eq:centered}; the
divergence equation follows from
\(\nabla_x\cdot U=a^{-2}\nabla_y\cdot V\).  The spatial map
\(y\mapsto x=e^{-\tau/2}y\) is a bijection for each \(\tau\), and
\[
   \|V(\tau)\|_\infty
   =
   \sqrt{-t}\,\|U(t)\|_\infty .
\]
Taking suprema over \(t<0\), equivalently \(\tau\in\mathbb R\), gives the
last assertion.
\end{proof}

\subsection{Preservation of axisymmetry and swirl variables}

For axisymmetric flows we write
\[
   x=(r\cos\theta,r\sin\theta,z),\qquad
   U=U_re_r+U_\theta e_\theta+U_ze_z,
\]
and in centered variables
\[
   y=(\rho\cos\theta,\rho\sin\theta,Z),\qquad
   V=V_\rho e_\rho+V_\theta e_\theta+V_Ze_Z.
\]
The scalar swirl variables are
\[
   \Gamma_U(r,z,t)=rU_\theta(r,z,t),\qquad
   \Gamma_V(\rho,Z,\tau)=\rho V_\theta(\rho,Z,\tau).
\]

\begin{lemma}[Preservation of axisymmetric structure]
\label{p3:lem:structure-preserved}
Axisymmetry and absence of swirl are preserved under
\eqref{p3:eq:self-similar}.  If \(U\) is axisymmetric, then
\[
   \Gamma_V(\rho,Z,\tau)=\Gamma_U(r,z,t),
   \qquad
   (r,z)=\sqrt{-t}\,(\rho,Z).
\]
Fixed physical \(z\)-periodicity is preserved by physical time shifts, but
under the centered change of variables it becomes a \(\tau\)-dependent period
in the \(Z\)-coordinate.
\end{lemma}

\begin{proof}
The self-similar map is a scalar spatial dilation and a scalar multiplication
of the velocity, so it commutes with rotations around the symmetry axis.
The cylindrical basis is unchanged by radial dilation, and
\[
   V_\theta(\rho,Z,\tau)
   =
   \sqrt{-t}\,U_\theta(r,z,t),
   \qquad
   r=\sqrt{-t}\rho,\quad z=\sqrt{-t}Z .
\]
Thus \(U_\theta\equiv0\) if and only if \(V_\theta\equiv0\), and
\[
   \Gamma_V=\rho V_\theta
   =
   \rho\sqrt{-t}\,U_\theta
   =
   rU_\theta
   =
   \Gamma_U .
\]
If \(U\) has physical axial period \(L\), then \(V\) has period
\(L/\sqrt{-t}=Le^{\tau/2}\) in \(Z\), not a fixed period.
\end{proof}

\subsection{Time shifts away from the singular time}

\begin{lemma}[Time shifts away from the singular time]
\label{p3:lem:time-shift}
Let \(U\) be a smooth Type I ancient solution and fix \(T<0\).  Then
\[
   W_T(x,s)=U(x,T+s),\qquad \Pi_T(x,s)=P(x,T+s),\qquad s<0,
\]
is a smooth bounded mild ancient solution of \eqref{p3:eq:NS} on
\(\mathbb R^3\times(-\infty,0)\).  Moreover the axisymmetric, no-swirl,
\(L^p\)-swirl, bounded-\(\Gamma\), and physical \(z\)-periodic properties are
preserved.
\end{lemma}

\begin{proof}
Time translation preserves the differential equation because
\(\partial_sW_T(x,s)=\partial_tU(x,T+s)\) and all spatial derivatives are
unchanged.  For the mild formulation, apply \eqref{p3:eq:mild-formula} to \(U\)
between the physical times \(T+s<T+t<0\):
\[
   U(T+t)=e^{(t-s)\Delta}U(T+s)
   -\int_{T+s}^{T+t}e^{(T+t-\eta)\Delta}\mathbb P\nabla\cdot
      (U(\eta)\otimes U(\eta))\,d\eta .
\]
With \(\eta=T+\sigma\), this becomes the mild formula for \(W_T\) on the
interval \((s,t)\).  For \(s<0\), \(T+s<T<0\), so the Type I bound gives
\[
   \|W_T(s)\|_\infty
   =
   \|U(T+s)\|_\infty
   \le
   \frac1{\sqrt{-T}}
   \sup_{\eta<0}\sqrt{-\eta}\,\|U(\eta)\|_\infty .
\]
The invariance assertions follow because only the time variable has been
translated.  In particular
\[
   \Gamma_{W_T}(r,z,s)=\Gamma_U(r,z,T+s),
\]
so \(L_s^\infty L_x^p\) and \(L^\infty\) bounds for \(\Gamma\) restrict to the
shifted solution, and a fixed physical axial period remains the same.
\end{proof}

\subsection{PDE hypotheses}

The arguments are stated directly in terms of hypotheses on ancient
solutions.  The axisymmetric hypotheses are:
\begin{enumerate}[label=\textup{(\roman*)}]
\item axisymmetry and \(U_\theta\equiv0\);
\item axisymmetry and the pointwise bound \(|U(x,t)|\le C/r\) for \(r>0\);
\item axisymmetry and
\[
   \Gamma=rU_\theta\in L_t^\infty L_x^p
   \bigl(\mathbb R^3\times(-\infty,0)\bigr),
   \qquad 1\le p<\infty;
\]
\item axisymmetry, periodicity in the physical axial variable \(z\), and
\(\Gamma\in L^\infty\).
\end{enumerate}
The additional hypothesis used later is the anchored Gallay--Wayne
perturbative vorticity condition of \Cref{p3:def:gw-anchored-condition}.

\section{Axisymmetric notation and cited Liouville theorems}
\label{p3:sec:axisymmetric-liouville-theorems}

\subsection{Cylindrical variables and the swirl scalar}

For an axisymmetric solution \(u\), write
\[
   u=u_re_r+u_\theta e_\theta+u_ze_z,\qquad
   \Gamma=ru_\theta,\qquad
   b=u_re_r+u_ze_z .
\]
When \(\Gamma\) is placed in \(L^p(\mathbb R^3)\), it is regarded as the
corresponding function of \(x=(r\cos\theta,r\sin\theta,z)\), independent of
\(\theta\):
\[
   \|\Gamma(t)\|_{L^p(\mathbb R^3)}^p
   =
   2\pi\int_{\mathbb R}\int_0^\infty
      |\Gamma(r,z,t)|^p r\,dr\,dz,\qquad 1\le p<\infty .
\]

\subsection{Constant solutions in the Type I frame}

\begin{definition}[Reduction to a spatial constant in the Type I frame]
The cited Liouville theorems are used below only through the consequence that,
after passing to the finite-terminal shift in the physical Type I frame, the
solution is spatially constant.  The Type I bound then forces that constant to
be zero.
\end{definition}

\begin{lemma}[Spatially constant mild solutions]\label{p3:lem:spatial-constant-mild}
Let \(W\) be a smooth mild solution of \eqref{p3:eq:NS} on
\(\mathbb R^3\times(-\infty,0)\).  If \(W(x,t)=b(t)\) for all \(x\) and \(t\),
then \(b\) is constant in time.
\end{lemma}

\begin{proof}
Insert \(W(x,t)=b(t)\) into \eqref{p3:eq:mild-formula}.  The heat semigroup
leaves constant vector fields fixed, and
\(\nabla\cdot(b(\sigma)\otimes b(\sigma))=0\).  Thus
\(b(t)=b(s)\) for every \(s<t<0\).
\end{proof}

\begin{lemma}[Patching constants on backward intervals]
\label{p3:lem:patch-constants}
Let \(U:\mathbb R^3\times(-\infty,0)\to\mathbb R^3\) be a function.  Assume
that for every \(T<0\) there exists \(c_T\in\mathbb R^3\) such that
\[
   U(x,t)=c_T\qquad(x\in\mathbb R^3,\ t<T).
\]
Then there is a single \(c\in\mathbb R^3\) such that
\[
   U(x,t)=c\qquad(x\in\mathbb R^3,\ t<0).
\]
\end{lemma}

\begin{proof}
If \(T_1<T_2<0\), the two constants agree on the nonempty interval
\((-\infty,T_1)\).  Hence all \(c_T\) are equal.  For any \(t<0\), choose
\(T\) with \(t<T<0\).
\end{proof}

\begin{lemma}[Constant Type I ancient solutions vanish]
\label{p3:lem:galilean-typeI-zero}
Let \(U\) be a Type I ancient solution.  If \(U(x,t)\equiv c\), then
\(c=0\).
\end{lemma}

\begin{proof}
The Type I bound gives
\[
   \sqrt{-t}\,|c|
   \le
   \sup_{s<0}\sqrt{-s}\,\|U(s)\|_\infty<\infty
   \qquad(t<0).
\]
Letting \(t\to-\infty\) gives \(c=0\).
\end{proof}

\subsection{Cited Liouville theorems}

The following Liouville theorems are used in the stated forms.

\begin{remark}[Match with the cited papers]\label{p3:rem:axi-citations}
The axisymmetric theorems below are imported from the cited references and are
recorded with the hypotheses used here.  In the Type I application they are
applied to the finite-terminal time shifts
\[
   W_T(x,s)=U(x,T+s),\qquad s<0,\quad T<0,
\]
which are bounded ancient mild solutions by
\Cref{p3:lem:time-shift}.  The resulting constants are then patched in \(T\)
and removed by the Type I bound.  The notation matches the cited papers
modulo the centered-versus-physical translation recorded in
\Cref{p3:lem:structure-preserved}.
\end{remark}

\begin{theorem}[KNSS no-swirl Liouville theorem]
\label{p3:thm:knss-noswirl}
Every bounded ancient mild axisymmetric no-swirl solution of the unforced
Navier--Stokes equations is reduced, in the form used here, to a spatially
constant axial drift in the physical Type I frame.
\end{theorem}

\begin{proof}
\smallskip
\noindent\emph{Imported theorem.}
We use the axisymmetric no-swirl Liouville theorem of
Koch, Nadirashvili, Seregin, and \v{S}ver\'ak
\cite[Theorem 5.2]{KNSS2009}.

\smallskip
\noindent\emph{Required hypotheses.}
The theorem is stated for bounded ancient mild solutions of the unforced
Navier--Stokes equations in the whole space which are axisymmetric and have no
swirl.  No drift-only or circulation-only hypothesis is being substituted for
the no-swirl assumption in this branch.

\smallskip
\noindent\emph{Verification for our profile.}
Let \(U\) be the Type I ancient physical representative.  For each \(T<0\),
the shifted solution
\[
   W_T(x,s)=U(x,T+s),\qquad s<0,
\]
is a smooth bounded mild ancient solution by \Cref{p3:lem:time-shift}.  The
no-swirl branch of \(\calC_{\mathrm{sd}}\) requires axisymmetry together with
\(u_\theta\equiv0\), and this structural hypothesis is preserved under the
time shift.  Therefore the cited theorem applies to \(W_T\).

\smallskip
\noindent\emph{Conclusion in the Type I frame.}
For each \(T<0\), the cited theorem gives that \(W_T\) is a spatially
constant axial drift in the frame used here.  The resulting constants patch in
\(T\) by
\Cref{p3:lem:patch-constants}, and \Cref{p3:lem:galilean-typeI-zero} then
forces the patched constant to vanish.
\end{proof}

\begin{theorem}[KNSS pointwise scale-invariant condition]
\label{p3:thm:knss-pointwise}
Every bounded ancient mild axisymmetric solution satisfying
\[
   |u(x,t)|\le \frac{C}{r}
   \qquad(r>0,\ t<0)
\]
is constant up to a Galilean transformation.  Only the spatially constant
conclusion in the physical Type I frame is used below; if the source theorem is
stated modulo Galilean transformations, this branch is invoked only after
reducing to that constant representative.  Zero is obtained only after the
constant-removal step.
\end{theorem}

\begin{proof}
\smallskip
\noindent\emph{Imported theorem.}
We use the pointwise scale-invariant axisymmetric Liouville theorem of
Koch--Nadirashvili--Seregin--\v{S}ver\'ak
\cite[Theorem 5.3]{KNSS2009}.

\smallskip
\noindent\emph{Required hypotheses.}
The source theorem applies to bounded ancient mild axisymmetric whole-space
solutions satisfying the full velocity condition
\[
   r|u(x,t)|\le C .
\]
Equivalently, away from the axis this is \(|u(x,t)|\le C/r\).  This branch
therefore uses the full velocity hypothesis.  Drift-only assumptions or
circulation-only assumptions are not included in this class.

\smallskip
\noindent\emph{Verification for our profile.}
Let \(U\) be the Type I ancient physical representative.  For each \(T<0\),
the shifted solution
\[
   W_T(x,s)=U(x,T+s),\qquad s<0,
\]
is a smooth bounded mild ancient solution by \Cref{p3:lem:time-shift}.  The
pointwise branch of \(\calC_{\mathrm{sd}}\) requires the full velocity bound
\(r|u|\le C\), not merely a bound on the drift \(u^r e_r+u^z e_z\) or on the
circulation \(\Gamma=ru_\theta\), and this structural hypothesis is preserved
under the time shift.  Therefore the cited theorem applies to \(W_T\).

\smallskip
\noindent\emph{Conclusion in the Type I frame.}
For each \(T<0\), the cited theorem gives a spatially constant solution in the
frame used here.  The resulting constants patch in \(T\) by
\Cref{p3:lem:patch-constants}.  In the Type I ancient frame, the patched
constant is eliminated by \Cref{p3:lem:galilean-typeI-zero}; in the pointwise
\(1/r\) branch it is also excluded directly by letting \(r\to\infty\) in the
pointwise bound, as recorded in \Cref{p3:rem:pointwise-r-infty}.
\end{proof}

\begin{theorem}[Lei--Zhang--Zhao]
\label{p3:thm:lzz}
Let \(u\) be a bounded ancient mild axisymmetric solution.  If
\[
   \Gamma=ru_\theta\in L_t^\infty L_x^p
   \bigl(\mathbb R^3\times(-\infty,0)\bigr),
   \qquad 1\le p<\infty,
\]
then the no-swirl case \(u_\theta\equiv0\) is handled by
\Cref{p3:thm:knss-noswirl}, while in the nontrivial-swirl case the cited
theorem yields that \(u\) is a constant vector field.
\end{theorem}

\begin{proof}
\smallskip
\noindent\emph{Imported theorem.}
We use the ancient axisymmetric circulation Liouville theorem of
Lei--Zhang--Zhao \cite[Theorem 1.3]{LZZ}.

\smallskip
\noindent\emph{Required hypotheses.}
The theorem assumes a bounded ancient mild axisymmetric solution in
\(\R^3\times(-\infty,0)\), nontrivial swirl, and finite critical circulation integrability
\[
   \Gamma=ru_\theta
   \in L_t^\infty L_x^p(\R^3\times(-\infty,0)),
   \qquad 1\le p<\infty .
\]
The no-swirl case \(u_\theta\equiv0\) is handled separately by
\Cref{p3:thm:knss-noswirl}; this theorem is invoked for the finite
\(L^p\)-circulation branch.

\smallskip
\noindent\emph{Verification for our profile.}
Let \(U\) be the Type I ancient physical representative.  For each \(T<0\),
the shifted solution
\[
   W_T(x,s)=U(x,T+s),\qquad s<0,
\]
is a smooth bounded mild ancient axisymmetric solution by
\Cref{p3:lem:time-shift}.  The finite-\(L^p\) circulation hypothesis is
preserved under the time shift.

If \((W_T)_\theta\equiv0\), the no-swirl branch
\Cref{p3:thm:knss-noswirl} applies.  If
\((W_T)_\theta\not\equiv0\), the Lei--Zhang--Zhao theorem applies and gives
that \(W_T\) is a constant vector field.

\smallskip
\noindent\emph{Conclusion in the Type I frame.}
In either case \(W_T\) is spatially constant.  The constants patch in \(T\) by
\Cref{p3:lem:patch-constants}, and
\Cref{p3:lem:galilean-typeI-zero} eliminates the patched constant in the
Type I ancient frame.
\end{proof}

\begin{theorem}[Lei--Ren--Zhang]
\label{p3:thm:lrz}
Let \(u\) be a bounded ancient mild axisymmetric solution, periodic in \(z\),
with
\[
   \Gamma=ru_\theta\in L^\infty
   \bigl(\mathbb R^3\times(-\infty,0)\bigr).
\]
Then \(u\) is a constant axial vector field.
\end{theorem}

\begin{proof}
\smallskip
\noindent\emph{Imported theorem.}
We use the periodic axisymmetric ancient Liouville theorem of
Lei--Ren--Zhang \cite[Theorem 1.1]{LRZ}.

\smallskip
\noindent\emph{Required hypotheses.}
The theorem is stated on the quotient
\(\R^2\times\mathbb T^1\) for bounded mild ancient axisymmetric solutions with
bounded swirl circulation \(\Gamma=ru_\theta\).  The periodic branch in this
paper is deliberately no larger than that source theorem: the physical
representative is assumed \(z\)-periodic and satisfies the same bounded
\(\Gamma\) condition.

\smallskip
\noindent\emph{Verification for our profile.}
Let \(U\) be the Type I ancient physical representative.  For each \(T<0\),
the shifted solution
\[
   W_T(x,s)=U(x,T+s),\qquad s<0,
\]
is a smooth bounded mild ancient solution by \Cref{p3:lem:time-shift}.  In
the periodic branch the periodicity is imposed on the physical axial variable
\(z\), not on the centered field, and the branch also requires bounded
circulation \(\Gamma=ru_\theta\).  These structural hypotheses are preserved
under the time shift.  Therefore the cited theorem applies to \(W_T\) after
descending to \(\R^2\times\mathbb T^1\) and normalizing the period by a
Navier--Stokes scaling, which does not change boundedness, mildness,
axisymmetry, or boundedness of \(\Gamma\).

\smallskip
\noindent\emph{Conclusion in the Type I frame.}
For each \(T<0\), the cited theorem gives a spatially constant solution in the
frame used here.  The resulting constants patch in \(T\) by
\Cref{p3:lem:patch-constants}, and \Cref{p3:lem:galilean-typeI-zero} then
eliminates the patched constant in the Type I ancient frame.
\end{proof}

\section{Unforced axisymmetric Liouville consequences}
\label{p3:sec:axisymmetric-consequences}

\subsection{No swirl}

\begin{proposition}[No-swirl Type I vanishing]\label{p3:prop:noswirl-typeI}
Let \(U\) be a smooth Type I ancient solution.  Assume \(U\) is axisymmetric
and \(U_\theta\equiv0\).  Then \(U\equiv0\).
\end{proposition}

\begin{proof}
Fix \(T<0\) and set \(W_T(x,s)=U(x,T+s)\).  By \Cref{p3:lem:time-shift},
\(W_T\) is a bounded ancient mild axisymmetric no-swirl solution.  Applying
\Cref{p3:thm:knss-noswirl} gives that \(W_T\) is a spatially constant axial
drift, say \(W_T(x,s)=b_T(s)e_z\).  By
\Cref{p3:lem:spatial-constant-mild}, \(b_T(s)\) is constant in \(s\).  Thus
\(U(x,t)=c_T\) for \(t<T\).  Since \(T<0\) is arbitrary,
\Cref{p3:lem:patch-constants} gives \(U(x,t)=c\) for all \(t<0\).  Finally,
\Cref{p3:lem:galilean-typeI-zero} gives \(c=0\).
\end{proof}

\subsection{Pointwise scale-invariant control}

\begin{proposition}[Pointwise \(1/r\) Type I vanishing]\label{p3:prop:pointwise-typeI}
Let \(U\) be a smooth Type I ancient solution.  Assume \(U\) is axisymmetric
and satisfies
\[
   |U(x,t)|\le \frac{C}{r}
   \qquad(r>0,\ t<0).
\]
Then \(U\equiv0\).
\end{proposition}

\begin{proof}
Fix \(T<0\).  By \Cref{p3:lem:time-shift}, the shifted solution
\[
   W_T(x,s)=U(x,T+s),\qquad s<0,
\]
is a bounded ancient mild axisymmetric solution and satisfies the same
pointwise bound
\[
   |W_T(x,s)|\le \frac{C}{r}.
\]
By \Cref{p3:thm:knss-pointwise}, \(W_T\) is spatially constant in the frame
used here.  Write \(W_T\equiv c_T\).  The pointwise bound gives
\[
   |c_T|\le \frac{C}{r}\qquad\text{for every }r>0,
\]
hence \(c_T=0\) by letting \(r\to\infty\).  Equivalently, one may patch the
constants in \(T\) and apply \Cref{p3:lem:galilean-typeI-zero}.  Thus
\(W_T\equiv0\), so \(U\equiv0\) on \(( -\infty,T)\).  Since \(T<0\) was
arbitrary, \(U\equiv0\).
\end{proof}

\begin{remark}[Pointwise bound rules out a constant directly]
\label{p3:rem:pointwise-r-infty}
In the pointwise scale-invariant case, even before invoking the Type I
constant-removal lemma, the pointwise bound \(|c|\le C/r\) for all \(r>0\)
forces \(c=0\) by letting \(r\to\infty\) when the Liouville theorem returns
a constant.  The Type I argument used in the proof above is consistent with
this stronger conclusion.
\end{remark}

\subsection{\texorpdfstring{Finite \(L^p\) control of \(\Gamma\)}
{Finite Lp control of Gamma}}

\begin{proposition}[Finite-\(L^p\) swirl Type I vanishing]
\label{p3:prop:lp-swirl-typeI}
Let \(U\) be a smooth Type I ancient solution.  Assume \(U\) is axisymmetric
and, for some \(1\le p<\infty\),
\[
   \Gamma=rU_\theta\in L_t^\infty L_x^p
   \bigl(\mathbb R^3\times(-\infty,0)\bigr).
\]
Then \(U\equiv0\).
\end{proposition}

\begin{proof}
Fix \(T<0\).  By \Cref{p3:lem:time-shift},
\[
   W_T(x,s)=U(x,T+s),\qquad s<0,
\]
is a bounded ancient mild axisymmetric solution.  The finite-\(L^p\)
circulation hypothesis is preserved by time shift:
\[
   \Gamma_{W_T}\in L_s^\infty L_x^p
   \bigl(\mathbb R^3\times(-\infty,0)\bigr).
\]
If \((W_T)_\theta\equiv0\), then the no-swirl Liouville branch gives that
\(W_T\) is spatially constant.  If \((W_T)_\theta\not\equiv0\), then the
Lei--Zhang--Zhao nontrivial-swirl theorem applies and again gives that
\(W_T\) is spatially constant.  In either case, \(W_T\equiv c_T\).
Patch the constants in \(T\) using \Cref{p3:lem:patch-constants}, and then
apply \Cref{p3:lem:galilean-typeI-zero}.
\end{proof}

\subsection{Periodic bounded swirl}

The periodic hypothesis below is always imposed on the physical ancient
representative \(U\), not on the centered profile \(V\).  Under the
self-similar transform a fixed physical \(z\)-period becomes a
\(\tau\)-dependent period in \(Z\), so it is not a fixed-period condition on
\(V\) (\Cref{p3:lem:structure-preserved}).

\begin{proposition}[Periodic bounded-swirl Type I vanishing]
\label{p3:prop:periodic-swirl-typeI}
Let \(U\) be a smooth Type I ancient solution.  Assume \(U\) is axisymmetric,
periodic in the physical axial variable \(z\), and satisfies
\[
   \Gamma=rU_\theta\in L^\infty
   \bigl(\mathbb R^3\times(-\infty,0)\bigr).
\]
Then \(U\equiv0\).
\end{proposition}

\begin{proof}
Fix \(T<0\).  By \Cref{p3:lem:time-shift}, the shifted solution \(W_T\) is
bounded, ancient, mild, axisymmetric, physically \(z\)-periodic, and has
bounded \(\Gamma\).  \Cref{p3:thm:lrz} gives that \(W_T\) is a constant axial vector
field.  The constants patch as \(T\) varies, and the Type I bound eliminates
the resulting constant by \Cref{p3:lem:galilean-typeI-zero}.
\end{proof}

\subsection{Proof of the axisymmetric Liouville theorem}

\begin{proof}[Proof of \Cref{p3:thm:axisymmetric-liouville}]
The four cases are precisely the hypotheses of
\Cref{p3:prop:noswirl-typeI,p3:prop:pointwise-typeI,p3:prop:lp-swirl-typeI,p3:prop:periodic-swirl-typeI}.
Each proposition gives \(U\equiv0\).  This contradicts
\Cref{p3:lem:nonzero-typeI-limit} when \(U\) is a nonzero Type I blow-up limit.
\end{proof}

\section{Anchored Gallay--Wayne perturbative exclusion}
\label{p3:sec:weighted-vorticity}

\subsection{The forward rescaled vorticity equation}

The Gallay--Wayne variable used here is the forward self-similar vorticity
time \(\theta\); its generator \(\Lambda\) below has the opposite drift
convention from the centered velocity generator used in Parts II--III.  The
argument is formulated directly in physical variables: from arbitrarily early
physical times \(a<0\) we launch genuine forward Gallay--Wayne Cauchy
problems and compare them with the corresponding anchored forward rescalings of
the same ancient solution.  No claim is made that the backward Type I centered
vorticity is itself a single Gallay--Wayne forward orbit.

Let \(w=\nabla\times v\).  The Gallay--Wayne forward rescaled vorticity equation is
\begin{equation}\label{p3:eq:weighted-vorticity}
   \partial_\theta w
   =
   \Lambda w-(v\cdot\nabla)w+(w\cdot\nabla)v,
   \qquad \nabla\cdot w=0,
\end{equation}
where
\begin{equation}\label{p3:eq:lambda}
   \Lambda w=\Delta w+\frac12\xi\cdot\nabla w+w,
\end{equation}
and \(v\) is recovered from \(w\) by the three-dimensional Biot--Savart law.
For \(m\ge0\), set
\[
   L^2_\sigma(m)
   =
   \left\{
      f\in L^2(\mathbb R^3;\mathbb R^3):
      \nabla\cdot f=0,\quad \|f\|_m<\infty
   \right\},
\]
where
\[
   \|f\|_m^2
   =
   \int_{\mathbb R^3}(1+|\xi|)^{2m}|f(\xi)|^2\,d\xi .
\]

\subsection{Gallay--Wayne small-data estimate}

\begin{theorem}[Gallay--Wayne small-data estimate]\label{p3:thm:gw-small-data}\label{p3:thm:gallay}
Let \(0<\mu\le1\) and \(m>2\mu+\frac12\).  There exist constants
\(r_m>0\) and \(K_m\ge1\) such that if
\[
   w_0\in L^2_\sigma(m),
   \qquad
   \|w_0\|_m\le r_m,
\]
then
\eqref{p3:eq:weighted-vorticity} has a unique global mild solution
\(w(\theta)=\Phi_\theta w_0\) in \(C([0,\infty);L^2_\sigma(m))\) and
\[
   \|w(\theta)\|_m
   \le
   K_m e^{-\mu\theta}\|w_0\|_m,\qquad \theta\ge0.
\]
In particular, if \(m>\frac72\), the choice \(\mu=1\) is admissible.
\end{theorem}

\begin{proof}
\smallskip
\noindent\emph{Imported theorem.}
We use Gallay--Wayne \cite[Theorem 2.3]{GallayWayne2002} for the rescaled
three-dimensional vorticity equation.

\smallskip
\noindent\emph{Required hypotheses.}
Their theorem is stated in the weighted divergence-free space
\(L^2_\sigma(m)\).  It assumes \(0<\mu\le1\), \(m>2\mu+1/2\), and small initial
data \(\|w_0\|_m\le r_0(m,\mu)\).  In this formulation no additional moment,
orthogonality, spectral-projection, or codimension condition is used beyond
membership in the stated small ball of \(L^2_\sigma(m)\).

\smallskip
\noindent\emph{Variable convention.}
The equation \eqref{p3:eq:weighted-vorticity} is the same forward
self-similar vorticity equation after the notational change
\(\xi\leftrightarrow y\); the sign of the confining drift is the forward
Gallay--Wayne convention used in this weighted-vorticity subsection.
\end{proof}

\subsection{Anchored Gallay--Wayne condition}

\begin{definition}[Anchored Gallay--Wayne perturbative vorticity condition]
\label{p3:def:gw-anchored-condition}\label{p3:def:weighted-vorticity-condition}
Let \(U\) be a smooth mild Type I ancient solution on
\(\mathbb R^3\times(-\infty,0)\), and set
\[
   \Omega=\nabla\times U .
\]
Fix \(m>\frac72\).  Let \(r_m>0\) and \(K_m\ge1\) be the constants in
\Cref{p3:thm:gw-small-data} with \(\mu=1\).  Fix \(0<\rho_m<r_m\).

We say that \(U\) satisfies the anchored Gallay--Wayne perturbative
vorticity condition if there exist anchor times
\[
   a_n\to-\infty
\]
and scales
\[
   \lambda_n>0
\]
such that
\[
   \lambda_n^2-a_n\to+\infty,
\]
and, for every \(n\), the anchored initial vorticity
\[
   w_n^0(\xi)
   :=
   \lambda_n^2\Omega(\lambda_n\xi,a_n)
\]
satisfies
\[
   w_n^0\in L^2_\sigma(m),
   \qquad
   \|w_n^0\|_m\le \rho_m.
\]
Moreover, the corresponding anchored initial velocity is the Biot--Savart
representative:
\[
   \lambda_n U(\lambda_n\xi,a_n)
   =
   \operatorname{BS}[w_n^0](\xi),
\]
where
\[
   \operatorname{BS}[w](\xi)
   =
   -\frac1{4\pi}
   \int_{\mathbb R^3}
   \frac{(\xi-\eta)\times w(\eta)}{|\xi-\eta|^3}\,d\eta .
\]
\end{definition}

\subsection{Anchored forward Gallay--Wayne rescalings}

\begin{lemma}[Local weighted-vorticity persistence for anchored smooth data]
\label{p3:lem:local-weighted-vorticity-persistence}
Let \(m>\frac72\), let \(a\in\mathbb R\), and let \(\lambda>0\).  Let
\[
   v^0\in C_b^\infty(\mathbb R^3)\cap C_0(\mathbb R^3),
   \qquad
   \nabla_\xi\cdot v^0=0,
\]
and suppose
\[
   w^0=\nabla_\xi\times v^0\in L^2_\sigma(m).
\]
Assume moreover that \(v^0\) is the decaying Biot--Savart representative of
\(w^0\):
\[
   v^0=\operatorname{BS}[w^0].
\]
Define the physical initial velocity at time \(a\) by
\[
   u_0(x)=\lambda^{-1}v^0(x/\lambda).
\]
Let \(u\) be the bounded mild Navier--Stokes solution with initial data
\(u_0\) on a sufficiently short interval \([a,a+\delta]\).  Let
\[
   \omega=\nabla_x\times u,
   \qquad
   t(\theta)=a+\lambda^2(e^\theta-1),
\]
and define the anchored forward self-similar vorticity and velocity by
\[
   w(\xi,\theta)
   :=
   \lambda^2 e^\theta
   \omega(\lambda e^{\theta/2}\xi,t(\theta)).
\]
\[
   v(\xi,\theta)
   :=
   \lambda e^{\theta/2}
   u(\lambda e^{\theta/2}\xi,t(\theta)).
\]
Then, for some \(\theta_0>0\),
\[
   w\in C([0,\theta_0];L^2_\sigma(m))
\]
\[
   w(\cdot,0)=w^0,
\]
and \(w\) solves the Gallay--Wayne rescaled vorticity equation
\[
   \partial_\theta w
   =
   \Lambda w-(v\cdot\nabla_\xi)w+(w\cdot\nabla_\xi)v,
   \qquad v=\operatorname{BS}[w],
\]
with initial data \(w^0\).
\end{lemma}

\begin{proof}
At time \(a\),
\[
   u_0(x)=\lambda^{-1}v^0(x/\lambda),
\]
so
\[
   \omega_0(x)
   =
   \nabla_x\times u_0(x)
   =
   \lambda^{-2}(\nabla_\xi\times v^0)(x/\lambda)
   =
   \lambda^{-2}w^0(x/\lambda).
\]
Therefore
\[
   w(\xi,0)
   =
   \lambda^2\omega_0(\lambda\xi)
   =
   w^0(\xi),
\]
and
\[
   v(\xi,0)
   =
   \lambda u_0(\lambda\xi)
   =
   v^0(\xi).
\]

The Biot--Savart normalization at \(\theta=0\) removes the possible bounded
harmonic velocity mode.  Since \(v^0=\operatorname{BS}[w^0]\in C_0\), the
scaled physical initial velocity
\[
   u_0(x)=\lambda^{-1}v^0(x/\lambda)
\]
belongs to \(C_0(\mathbb R^3)\).  The standard local mild theory in
\(C_0\cap C_b^\infty\) gives a short-time solution \(u(t)\in C_0\), and
therefore the rescaled velocity \(v(\theta)\) also belongs to \(C_0\) for
\(0\le\theta\le\theta_0\).

For each such \(\theta\), the difference
\[
   h(\theta):=v(\theta)-\operatorname{BS}[w(\theta)]
\]
is divergence-free and curl-free.  Hence \(h(\theta)\) is spatially constant.
But both \(v(\theta)\) and the Biot--Savart representative
\(\operatorname{BS}[w(\theta)]\) are in the decaying gauge, so
\(h(\theta)\in C_0(\mathbb R^3)\).  A spatially constant \(C_0\) function is
zero.  Thus
\[
   v=\operatorname{BS}[w]
\]
on \([0,\theta_0]\).

For \(C_0\cap C_b^\infty\) divergence-free data, the bounded mild solution is
smooth on a short time interval and has bounded spatial derivatives there.  Its
vorticity solves the physical vorticity equation classically.  The anchored forward
self-similar change of variables above transforms that physical vorticity
equation into
\[
   \partial_\theta w
   =
   \Lambda w-(v\cdot\nabla_\xi)w+(w\cdot\nabla_\xi)v.
\]
A standard weighted energy estimate with weight \((1+|\xi|)^{2m}\), using the
boundedness of \(v\) and \(\nabla_\xi v\) on a sufficiently short interval, gives
\[
   \sup_{0\le\theta\le\theta_0}\|w(\theta)\|_m<\infty
\]
for some \(\theta_0>0\), and continuity in \(L^2_\sigma(m)\).  Hence the
transformed vorticity belongs to the Gallay--Wayne uniqueness class on
\([0,\theta_0]\).
\end{proof}

\begin{lemma}[Bounded physical pullback of anchored Gallay--Wayne data]
\label{p3:lem:gw-bounded-physical-pullback}
Let \(m>\frac72\), let \(w^0\in L^2_\sigma(m)\) satisfy the
Gallay--Wayne small-data hypothesis.
Assume, in addition, that
\[
   v^0=\operatorname{BS}[w^0]\in
   C_b^\infty(\mathbb R^3)\cap C_0(\mathbb R^3),
   \qquad
   \nabla_\xi\cdot v^0=0,
   \qquad
   \nabla_\xi\times v^0=w^0
\]
in distributions.  Let
\[
   w^{GW}(\theta)=\Phi_\theta w^0,
   \qquad
   v^{GW}(\theta)=\operatorname{BS}[w^{GW}(\theta)]
\]
be the Gallay--Wayne mild solution.  Fix \(a<T<0\) and \(\lambda>0\), and define
\[
   A(t)=\lambda^2+t-a,
   \qquad
   \theta(t)=\log\frac{A(t)}{\lambda^2},
\]
\[
   \widetilde U(x,t)
   =
   A(t)^{-1/2}
   v^{GW}\!\left(\frac{x}{\sqrt{A(t)}},\theta(t)\right),
   \qquad a\le t\le T.
\]
Then \(\widetilde U\) is a bounded mild Navier--Stokes solution on every finite
interval \([a,b]\subset[a,T]\).  In particular,
\[
   \sup_{a\le t\le b}\|\widetilde U(t)\|_{L^\infty}<\infty .
\]
\end{lemma}

\begin{proof}
The inverse Gallay--Wayne change of variables sends the rescaled vorticity
equation to the physical vorticity equation and sends
\(v^{GW}=\operatorname{BS}[w^{GW}]\) to a physical velocity field
\(\widetilde U\).  It remains to justify that this velocity is bounded and
mild on each compact physical interval.

At \(\theta=0\), the assumptions give
\[
   v^0=\operatorname{BS}[w^0]\in C_b^\infty\cap C_0,
   \qquad
   \nabla_\xi\times v^0=w^0\in L^2_\sigma(m).
\]
Define the physical initial velocity
\[
   u_0(x)=\lambda^{-1}v^0(x/\lambda).
\]
The standard \(C_0\cap C_b^\infty\)-data mild theory gives a bounded mild physical
solution with initial data \(u_0\) on a short interval starting at \(t=a\).  By
\Cref{p3:lem:local-weighted-vorticity-persistence}, its anchored transformed
vorticity lies in
\[
   C([0,\theta_0];L^2_\sigma(m))
\]
for some \(\theta_0>0\) and solves the Gallay--Wayne equation with initial data
\(w^0\).  By uniqueness in the Gallay--Wayne class, it agrees with
\(w^{GW}\) on this initial interval.
Therefore
\[
   \sup_{0\le\theta\le\theta_0}\|v^{GW}(\theta)\|_{L^\infty}<\infty .
\]

Next fix any finite rescaled time interval \([0,\Theta]\).  Since
\(m>\frac32\), Cauchy--Schwarz against the weight gives
\[
   \|w^{GW}(\theta)\|_{L^1}
   \le
   C_m\|w^{GW}(\theta)\|_m,
   \qquad 0\le\theta\le\Theta,
\]
so \Cref{p3:thm:gw-small-data} yields
\[
   \sup_{0\le\theta\le\Theta}\|w^{GW}(\theta)\|_{L^1}<\infty .
\]
For every \(0<\theta_0<\Theta\), standard parabolic smoothing for the mild
solution of \eqref{p3:eq:weighted-vorticity} gives
\[
   \sup_{\theta_0\le\theta\le\Theta}\|w^{GW}(\theta)\|_{L^p}<\infty
\]
for any chosen \(p>3\).  Splitting the Biot--Savart kernel into the regions
\(|\xi-\eta|\le1\) and \(|\xi-\eta|>1\) gives the standard estimate
\[
   \|\operatorname{BS}[w]\|_{L^\infty}
   \le
   C_p\bigl(\|w\|_{L^1}+\|w\|_{L^p}\bigr),
   \qquad p>3,
\]
because the far part is controlled by \(\|w\|_{L^1}\) and the near part is
controlled by H\"older since
\(|x|^{-2}\mathbf 1_{\{|x|\le1\}}\in L^{p/(p-1)}(\mathbb R^3)\) for \(p>3\).
Hence
\[
   \sup_{\theta_0\le\theta\le\Theta}\|v^{GW}(\theta)\|_{L^\infty}<\infty .
\]
Together with the initial short-time bounded mild solution, this gives
\[
   \sup_{0\le\theta\le\Theta}\|v^{GW}(\theta)\|_{L^\infty}<\infty .
\]

Now let \([a,b]\subset[a,T]\), and set
\[
   \Theta=\log\frac{\lambda^2+b-a}{\lambda^2}.
\]
Since \(A(t)\ge\lambda^2>0\) on \([a,b]\), the inverse transform gives
\[
   \|\widetilde U(t)\|_{L^\infty}
   =
   A(t)^{-1/2}\|v^{GW}(\theta(t))\|_{L^\infty}
   \le
   \lambda^{-1}\sup_{0\le\theta\le\Theta}\|v^{GW}(\theta)\|_{L^\infty}
\]
for \(a\le t\le b\).  Therefore
\[
   \sup_{a\le t\le b}\|\widetilde U(t)\|_{L^\infty}<\infty .
\]

Finally, the Duhamel formula for the Gallay--Wayne mild solution transforms
under the inverse forward self-similar map into the standard physical
Navier--Stokes mild formula on \([a,b]\).  Therefore \(\widetilde U\) is a
bounded mild Navier--Stokes solution on every compact subinterval of
\([a,T]\).
\end{proof}

\begin{lemma}[Physical uniqueness comparison for anchored Gallay--Wayne flows]
\label{p3:lem:anchored-gw-physical-comparison}
Let \(m>\frac72\).  Let \(U\) be a smooth mild Navier--Stokes solution on
\(\mathbb R^3\times[a,T]\), with \(a<T<0\).  Let \(w^0\in L^2_\sigma(m)\)
satisfy the Gallay--Wayne small-data hypothesis, and let
\[
   w^{GW}(\theta)=\Phi_\theta w^0,
   \qquad
   v^{GW}(\theta)=\operatorname{BS}[w^{GW}(\theta)]
\]
be the Gallay--Wayne mild solution.

Fix \(\lambda>0\), set
\[
   A(t)=\lambda^2+t-a,
   \qquad
   \theta(t)=\log\frac{A(t)}{\lambda^2},
\]
and define the physical pullback of the Gallay--Wayne solution by
\[
   \widetilde U(x,t)
   =
   A(t)^{-1/2}
   v^{GW}\!\left(\frac{x}{\sqrt{A(t)}},\theta(t)\right),
   \qquad a\le t\le T.
\]
Assume the anchored initial velocity and vorticity agree with \(U\):
\[
   \lambda U(\lambda\xi,a)=\operatorname{BS}[w^0](\xi),
\]
and
\[
   w^0(\xi)=\lambda^2(\nabla_x\times U)(\lambda\xi,a).
\]
Then
\[
   \widetilde U(x,t)=U(x,t)
   \qquad
   \text{for }a\le t\le T.
\]
\end{lemma}

\begin{proof}
The initial agreement gives
\[
   v^0(\xi):=\operatorname{BS}[w^0](\xi)=\lambda U(\lambda\xi,a).
\]
Because \(\operatorname{BS}[w^0]\) is the decaying Biot--Savart
representative, \(v^0\in C_0(\mathbb R^3)\).  Since \(U\) is smooth and
bounded on the finite terminal interval under consideration, the anchored
representative also lies in \(C_b^\infty\).  Hence
\[
   v^0\in C_b^\infty(\mathbb R^3)\cap C_0(\mathbb R^3).
\]
Moreover,
\[
   \nabla_\xi\cdot v^0(\xi)
   =
   \lambda^2(\nabla_x\cdot U)(\lambda\xi,a)
   =
   0,
\]
since \(U\) is divergence-free.
The anchored vorticity agreement gives
\[
   \nabla_\xi\times v^0=w^0.
\]
Hence the hypotheses of
\Cref{p3:lem:gw-bounded-physical-pullback} are satisfied.  Therefore the
inverse Gallay--Wayne pullback \(\widetilde U\) is a bounded mild
Navier--Stokes solution on every compact subinterval of \([a,T]\).  At
\(t=a\), \(A(a)=\lambda^2\), so
\[
   \widetilde U(x,a)
   =
   \lambda^{-1}
   v^{GW}\!\left(\frac{x}{\lambda},0\right)
   =
   \lambda^{-1}
   \operatorname{BS}[w^0]\!\left(\frac{x}{\lambda}\right).
\]
By the anchored Biot--Savart normalization, this equals \(U(x,a)\).

Let \(Z=U-\widetilde U\).  On any compact subinterval \([a,b]\subset[a,T]\),
the two mild solutions satisfy
\[
   Z(t)
   =
   -\int_a^t
   e^{(t-s)\Delta}\mathbb P\nabla\cdot
   \bigl(Z\otimes U+\widetilde U\otimes Z\bigr)(s)\,ds .
\]
Using the Oseen estimate,
\[
   \|e^{(t-s)\Delta}\mathbb P\nabla\cdot F\|_\infty
   \le C(t-s)^{-1/2}\|F\|_\infty,
\]
we get
\[
   \sup_{a\le t\le b}\|Z(t)\|_\infty
   \le
   C(M_U+M_{\widetilde U})(b-a)^{1/2}
   \sup_{a\le t\le b}\|Z(t)\|_\infty,
\]
where
\[
   M_U=\sup_{a\le t\le b}\|U(t)\|_\infty,
   \qquad
   M_{\widetilde U}=\sup_{a\le t\le b}\|\widetilde U(t)\|_\infty .
\]
Choosing \(b-a\) small enough forces \(Z\equiv0\) on \([a,b]\).  Iterating over
finitely many subintervals gives \(Z\equiv0\) on \([a,T]\).
\end{proof}

\begin{lemma}[Anchored transforms satisfy the Gallay--Wayne equation]
\label{p3:lem:anchored-gw-transform}
Let \(U\) be a smooth mild Type I ancient solution satisfying
\Cref{p3:def:gw-anchored-condition}.  For each anchor-scale pair
\((a_n,\lambda_n)\), define
\[
   t_n(\theta)=a_n+\lambda_n^2(e^\theta-1),
   \qquad
   x=\lambda_n e^{\theta/2}\xi,
\]
and
\[
   w_n(\xi,\theta)
   =
   \lambda_n^2e^\theta
   \Omega(\lambda_n e^{\theta/2}\xi,t_n(\theta)),
\]
\[
   v_n(\xi,\theta)
   =
   \lambda_n e^{\theta/2}
   U(\lambda_n e^{\theta/2}\xi,t_n(\theta)).
\]
Then \(w_n\) is the restriction, on the interval
\[
   0\le \theta<\log\frac{\lambda_n^2-a_n}{\lambda_n^2},
\]
of the Gallay--Wayne mild solution with initial data \(w_n^0\).  In
particular, for every \(\theta\ge0\) with \(t_n(\theta)<0\),
\[
   \|w_n(\theta)\|_m
   \le
   K_m e^{-\theta}\|w_n^0\|_m
   \le
   K_m\rho_m .
\]
\end{lemma}

\begin{proof}
The change of variables is exactly the forward self-similar change of
variables used by Gallay--Wayne, translated to the initial time \(a_n\)
and length scale \(\lambda_n\).  Therefore the vorticity equation for
\(\Omega=\nabla\times U\) becomes
\[
   \partial_\theta w_n
   =
   \Lambda w_n-(v_n\cdot\nabla)w_n+(w_n\cdot\nabla)v_n .
\]
At \(\theta=0\),
\[
   w_n(\xi,0)=\lambda_n^2\Omega(\lambda_n\xi,a_n)=w_n^0(\xi),
\]
and the Biot--Savart normalization in
\Cref{p3:def:gw-anchored-condition} gives
\[
   v_n(\xi,0)=\lambda_n U(\lambda_n\xi,a_n)=\operatorname{BS}[w_n^0](\xi).
\]

Let \(w_n^{GW}(\theta)=\Phi_\theta w_n^0\) be the Gallay--Wayne mild solution
with initial data \(w_n^0\), and let
\[
   v_n^{GW}(\theta)=\operatorname{BS}[w_n^{GW}(\theta)].
\]
For
\[
   A_n(t)=\lambda_n^2+t-a_n,
   \qquad
   \theta_n(t)=\log\frac{A_n(t)}{\lambda_n^2},
\]
define the corresponding physical velocity
\[
   \widetilde U_n(x,t)
   =
   A_n(t)^{-1/2}
   v_n^{GW}\!\left(
      \frac{x}{\sqrt{A_n(t)}},
      \theta_n(t)
   \right),
   \qquad a_n\le t<0.
\]
For \(\theta=0\), the identity
\[
   w_n(\xi,0)=w_n^{GW}(\xi,0)=w_n^0(\xi)
\]
holds by definition.  Now fix
\[
   0<\theta<\log\frac{\lambda_n^2-a_n}{\lambda_n^2},
\]
and set \(T=t_n(\theta)\).  Then \(a_n<T<0\), so by
\Cref{p3:lem:anchored-gw-physical-comparison},
\[
   \widetilde U_n(x,t)=U(x,t)
\]
for \(a_n\le t\le T\).  Evaluating the equality at time \(T\) and taking curls gives
\[
   w_n(\xi,\theta)=w_n^{GW}(\xi,\theta).
\]
Since \(\theta\) was arbitrary in the stated interval, \(w_n\) is exactly the
restriction of the Gallay--Wayne mild solution launched from \(w_n^0\).
Therefore \Cref{p3:thm:gw-small-data} with \(\mu=1\) yields
\[
   \|w_n(\theta)\|_m
   =
   \|w_n^{GW}(\theta)\|_m
   \le
   K_m e^{-\theta}\|w_n^0\|_m
   \le
   K_m\rho_m .
\]
\end{proof}

\subsection{Anchored perturbative exclusion}

\begin{proposition}[Anchored Gallay--Wayne perturbative branch vanishes]
\label{p3:prop:gw-anchored-typeI-zero}
Let \(U\) be a smooth mild Type I ancient solution on
\(\mathbb R^3\times(-\infty,0)\).  If \(U\) satisfies the anchored
Gallay--Wayne perturbative vorticity condition of
\Cref{p3:def:gw-anchored-condition}, then
\[
   U\equiv0.
\]
\end{proposition}

\begin{proof}
Let
\[
   \Omega=\nabla\times U .
\]
Fix \(T<0\).  For all sufficiently large \(n\), \(a_n<T\).  Define
\[
   A_n(T)=\lambda_n^2+T-a_n,
   \qquad
   \theta_n(T)=\log\frac{A_n(T)}{\lambda_n^2}.
\]
Then
\[
   A_n(T)\to+\infty
\]
by the anchoring condition \(\lambda_n^2-a_n\to+\infty\).

At physical time \(T\), the anchored transform gives
\[
   w_n(\xi,\theta_n(T))
   =
   A_n(T)\,
   \Omega\bigl(\sqrt{A_n(T)}\,\xi,T\bigr).
\]
By \Cref{p3:lem:anchored-gw-transform},
\[
   \|w_n(\theta_n(T))\|_m
   \le
   K_m\rho_m .
\]

Let \(\varphi\in C_c^\infty(\mathbb R^3)^3\).  Changing variables
\(x=\sqrt{A_n(T)}\,\xi\), we obtain
\[
\begin{aligned}
   \left|\int_{\mathbb R^3}\Omega(x,T)\cdot\varphi(x)\,dx\right|
   &=
   A_n(T)^{1/2}
   \left|
      \int_{\mathbb R^3}
      w_n(\xi,\theta_n(T))
      \cdot
      \varphi(\sqrt{A_n(T)}\,\xi)
      \,d\xi
   \right|  \\
   &\le
   A_n(T)^{1/2}
   \|w_n(\theta_n(T))\|_m
   \left\|
      (1+|\xi|)^{-m}
      \varphi(\sqrt{A_n(T)}\,\xi)
   \right\|_{L^2_\xi}.
\end{aligned}
\]
Since \((1+|\xi|)^{-m}\le1\),
\[
   \left\|
      (1+|\xi|)^{-m}
      \varphi(\sqrt{A_n(T)}\,\xi)
   \right\|_{L^2_\xi}
   \le
   A_n(T)^{-3/4}\|\varphi\|_{L^2}.
\]
Therefore
\[
   \left|\int_{\mathbb R^3}\Omega(x,T)\cdot\varphi(x)\,dx\right|
   \le
   K_m\rho_m
   A_n(T)^{-1/4}
   \|\varphi\|_{L^2}.
\]
Letting \(n\to\infty\) gives
\[
   \int_{\mathbb R^3}\Omega(x,T)\cdot\varphi(x)\,dx=0.
\]
Thus
\[
   \Omega(\cdot,T)=0
\]
in distributions.  Since \(T<0\) was arbitrary,
\[
   \Omega\equiv0
   \qquad
   \text{on }\mathbb R^3\times(-\infty,0).
\]

For each fixed \(T<0\), \(U(\cdot,T)\) is smooth, bounded, divergence-free,
and curl-free.  Hence each component of \(U(\cdot,T)\) is a bounded harmonic
function on \(\mathbb R^3\), so \(U(\cdot,T)\) is spatially constant.  Since
this holds for every \(T<0\), write
\[
   U(x,t)=c(t).
\]
By \Cref{p3:lem:spatial-constant-mild}, \(c(t)\) is constant in time, and
then \Cref{p3:lem:galilean-typeI-zero} gives \(c=0\).  Therefore
\(U\equiv0\).
\end{proof}

\begin{remark}
The anchored argument above is deliberately formulated in physical variables.
It avoids treating the backward Type I centered vorticity as a single
Gallay--Wayne forward orbit.  Gallay--Wayne's estimates are used only on
ordinary forward Cauchy problems launched from the anchor times \(a_n\).
\end{remark}

\begin{proof}[Proof of \Cref{p3:thm:weighted-vorticity-vanishing-intro}]
This is \Cref{p3:prop:gw-anchored-typeI-zero}.
\end{proof}

\section{Summary of PDE consequences}
\label{p3:sec:summary}

\subsection{Zero conclusions}

\begin{theorem}[Zero conclusions for structured Type I ancient solutions]
\label{p3:thm:typeI-zero-hypotheses}
Let \(U\) be a smooth mild Type I ancient solution in the sense of
\Cref{p3:def:typeI-ancient}.  If \(U\) satisfies one of the following
hypotheses:
\begin{enumerate}[label=\textup{(\roman*)}]
\item axisymmetry and \(U_\theta\equiv0\);
\item axisymmetry and the pointwise bound \(|U(x,t)|\le C/r\);
\item axisymmetry and finite-\(L^p\) control of the swirl scalar,
\[
   \Gamma\in L_t^\infty L_x^p,\qquad 1\le p<\infty;
\]
\item axisymmetry, \(z\)-periodicity, and bounded \(\Gamma\);
\item \(U\) satisfies the anchored Gallay--Wayne perturbative vorticity
condition of \Cref{p3:def:gw-anchored-condition}.
\end{enumerate}
then \(U\equiv0\).
\end{theorem}

\begin{proof}
Items (i)--(iv) are \Cref{p3:thm:axisymmetric-liouville}.  Item (v) is
\Cref{p3:prop:gw-anchored-typeI-zero}.
\end{proof}

\begin{corollary}[Structured hypotheses cannot hold for nonzero Type I blow-up limits]
\label{p3:cor:structured-hypotheses-nonzero-blowup}
A nonzero Type I blow-up limit cannot satisfy any of the hypotheses in
\Cref{p3:thm:typeI-zero-hypotheses}.
\end{corollary}

\begin{proof}
If a nonzero Type I blow-up limit satisfied one of the hypotheses, then
\Cref{p3:thm:typeI-zero-hypotheses} would give \(U\equiv0\), contradicting
\Cref{p3:lem:nonzero-typeI-limit}.
\end{proof}

\subsection{Consequences for Type I blow-up arguments}

\begin{proof}[Proof of \Cref{p3:thm:typeI-consequences}]
The assertion is \Cref{p3:cor:structured-hypotheses-nonzero-blowup}.
\end{proof}

\begin{theorem}[Combined zero conclusion]
\label{p3:thm:stated-hypotheses-zero}
Let \(U\) be a Type I ancient solution obtained as a nonzero Type I blow-up
limit.  If \(U\) satisfies one of the hypotheses in
\Cref{p3:thm:typeI-zero-hypotheses}, then \(U\equiv0\), a contradiction.
\end{theorem}

\begin{proof}
The assertion is precisely the conclusion of
\Cref{p3:thm:typeI-zero-hypotheses}.
\end{proof}

\subsection{Summary table}

\begin{center}
\begin{tabular}{p{0.36\textwidth}p{0.18\textwidth}p{0.34\textwidth}}
\hline
Hypothesis & Conclusion & Analytic ingredient \\
\hline
no swirl & \(U\equiv0\) & KNSS plus Type I constant removal \\
pointwise \(1/r\) & \(U\equiv0\) & KNSS plus \(r\to\infty\) constant removal \\
\(\Gamma\in L_t^\infty L_x^p\) & \(U\equiv0\) & KNSS/LZZ split plus Type I constant removal \\
periodic bounded \(\Gamma\) & \(U\equiv0\) & LRZ plus Type I constant removal \\
anchored Gallay--Wayne perturbative vorticity & \(U\equiv0\) & anchored Gallay--Wayne decay plus physical uniqueness comparison \\
\hline
\end{tabular}
\end{center}

The no-swirl, pointwise scale-invariant, finite-\(L^p\)-swirl, and periodic
bounded-swirl cases first reduce, by the cited axisymmetric Liouville inputs
and the theorem-level KNSS/LZZ split in the finite-\(L^p\) branch, to spatial
constants or axial drifts in the physical Type I frame; the manuscript then
removes those constants internally.  The perturbative weighted-vorticity case
follows from the anchored Gallay--Wayne forward-rescaling argument together
with the physical uniqueness comparison.
Thus every zero conclusion used in the Type I application rests on one of the
explicit PDE hypotheses in the table.

\clearpage
\section*{Part V. Final Assembly: Conditional Local Type I Exclusion}
\label{p1:sec:final-assembly}

\subsection*{Overview}

This part collects the preceding results and proves the local pointwise Type I
exclusion conditional on the residual closure hypothesis.  The argument is
linear: local concentration, local endpoint weak-Serrin no-escape, admissible
local Type I concentration sequence, raw Seregin extraction, passage to the
raw generated state space with invariant nonvanishing, exhaustion into
Liouville classes plus residual class, exclusion of the four non-residual
classes, and residual closure.

\section{Final assembly}

\begin{theorem}[Local Type I regularity under residual closure]
\label{p1:thm:final-assembly}
Let \((u,p)\) be a finite-energy suitable weak solution on
\(\R^3\times(0,T)\), and let \(z_*=(x_*,T)\) be a terminal point satisfying
the local pointwise Type I bound, for some \(\rho>0\),
\[
  \esssup_{T-\rho^2<t<T}
  \sqrt{T-t}\,\|u(t)\|_{L^\infty(B_\rho(x_*))}<\infty .
\]
Assume the residual closure hypothesis \Cref{p1:hyp:no-remainder} holds for
every raw generated descendant hull
\[
   \calS_{\mathrm{raw\text{-}gen}}(u,p,z_*;R_0,\eta_*)
\]
associated with invariant normalization data obtained from admissible local
Type I concentration sequences at \(z_*\).
Then \(z_*\) is regular.
\end{theorem}

\begin{proof}
Assume for contradiction that \(z_*=(x_*,T)\) is singular.

\emph{Step 1: concentration.}
By \Cref{p0:cor:velocity-concentration}, there are \(\eta_v>0\) and
\(r_n\downarrow0\) such that
\[
   C(z_*,r_n)\ge\eta_v .
\]

\emph{Step 2: terminal no-escape.}
The local pointwise Type I bound gives
\[
   u\in L_t^{2,\infty}L_x^\infty(Q_\rho(z_*)).
\]
By \Cref{p0:lem:terminal-A-from-weak-serrin}, for each \(0<\theta<1\) the
terminal local \(A\)-quantity is uniformly bounded along the chosen
small-scale sequence.  Hence \Cref{p0:thm:local-weak-serrin-typeI} gives that
the rescaled sequence satisfies
\[
   \lim_{\sigma\downarrow0}\sup_n
   \int_{-\sigma}^{0}\int_{B_1}|u_n|^3=0 .
\]

\emph{Step 3: nontrivial ancient extraction.}
\Cref{p0:cor:paper0-local-typeI-admissible} promotes the sequence to
an admissible local Type I concentration sequence.  The concentration and no-escape
properties give, by
\Cref{p0:lem:paper0-local-typeI-into-terminal-dichotomy}, a fixed compact
negative-time cylinder carrying positive \(L^3\) mass.  Then
\Cref{p0:lem:local-seregin-extraction} yields a raw extracted normalized
Seregin ancient limit \(V\) satisfying
\[
   \iint_K |V|^3>0
\]
for some compact \(K\Subset\R^3\times\R\).  By
\Cref{p1:lem:compact-to-invariant-nonzero}, \(V\) also satisfies the
time-translation invariant local \(L^3\) nonvanishing condition
\eqref{p1:eq:invariant-nonvanishing} for some data \((R_0,\eta_*)\).

\emph{Step 4: class decomposition.}
Let
\[
   \calS=\calS_{\mathrm{raw\text{-}gen}}(u,p,z_*;R_0,\eta_*)
\]
be the raw generated descendant hull formed with the fixed invariant
normalization data obtained from the extracted profile \(V\).  By construction, the
extracted profile \(V\) belongs to \(\calS\).  By \Cref{p1:prop:exhaustion},
\[
   V\in
   \calC_{\mathrm{small}}
   \cup
   \calC_{\mathrm{stat},L^3}
   \cup
   \calC_{L^3\mathrm{-tight}}
   \cup
   \calC_{\mathrm{sd}}
   \cup
   \calR(\calS).
\]

\emph{Step 5: contradiction.}
The first four alternatives are excluded by
\Cref{p1:prop:known-class-exclusion}; each class definition includes exactly
the hypotheses needed by its corresponding Liouville theorem.  The residual
alternative is excluded by \Cref{p1:hyp:no-remainder}, applied to the raw
generated descendant hull
\(\calS_{\mathrm{raw\text{-}gen}}(u,p,z_*;R_0,\eta_*)\).  No alternative is possible, contradicting
the invariant nonvanishing normalization of \Cref{p1:cor:normalized-nonzero}.
This contradiction proves that \(z_*\) is regular.
\end{proof}

\begin{remark}[Conditional conclusion]
Assuming \Cref{p1:hyp:no-remainder}, \Cref{p1:thm:final-assembly} gives the
local pointwise Type I exclusion within the stated raw generated local pointwise
Type I extraction class.
\end{remark}

\clearpage
\begin{thebibliography}{99}

\bibitem{AlbrittonBarker2019}
D.~Albritton and T.~Barker,
\emph{On local Type I singularities of the Navier--Stokes equations and
Liouville theorems},
J. Math. Fluid Mech. \textbf{21} (2019), Article 43.

\bibitem{CKN1982}
L.~Caffarelli, R.~Kohn, and L.~Nirenberg,
\emph{Partial regularity of suitable weak solutions of the Navier--Stokes
 equations},
Comm. Pure Appl. Math. \textbf{35} (1982), no. 6, 771--831,
doi:\ \nolinkurl{10.1002/cpa.3160350604}.

\bibitem{ESS2003}
L.~Escauriaza, G.~A.~Seregin, and V.~{\v S}verak,
\emph{\(L_{3,\infty}\)-solutions of the Navier--Stokes equations and
 backward uniqueness},
Russian Math. Surveys \textbf{58} (2003), no. 2, 211--250,
doi:\ \nolinkurl{10.1070/RM2003v058n02ABEH000609}.

\bibitem{Giga1986}
Y.~Giga,
\emph{Solutions for semilinear parabolic equations in \(L^p\) and regularity
 of weak solutions of the Navier--Stokes system},
J. Differential Equations \textbf{62} (1986), no. 2, 186--212,
doi:\ \nolinkurl{10.1016/0022-0396(86)90096-3}.

\bibitem{GustafsonKangTsai2007}
S.~Gustafson, K.~Kang, and T.-P.~Tsai,
\emph{Interior regularity criteria for suitable weak solutions of the
 Navier--Stokes equations},
Comm. Math. Phys. \textbf{273} (2007), no. 1, 161--176,
doi:\ \nolinkurl{10.1007/s00220-007-0214-6}.

\bibitem{Leray1934}
J.~Leray,
\emph{Sur le mouvement d'un liquide visqueux emplissant l'espace},
Acta Math. \textbf{63} (1934), no. 1, 193--248,
doi:\ \nolinkurl{10.1007/BF02547354}.

\bibitem{Lin1998}
F.-H.~Lin,
\emph{A new proof of the Caffarelli--Kohn--Nirenberg theorem},
Comm. Pure Appl. Math. \textbf{51} (1998), no. 3, 241--257,
doi:\ \nolinkurl{10.1002/(SICI)1097-0312(199803)51:3<241::AID-CPA2>3.0.CO;2-A}.

\bibitem{Scheffer1976}
V.~Scheffer,
\emph{Partial regularity of solutions to the Navier--Stokes equations},
Pacific J. Math. \textbf{66} (1976), no. 2, 535--552,
doi:\ \nolinkurl{10.2140/pjm.1976.66.535}.

\bibitem{Seregin2012}
G.~Seregin,
\emph{A certain necessary condition of potential blow up for Navier--Stokes
 equations},
Comm. Math. Phys. \textbf{312} (2012), no. 3, 833--845,
doi:\ \nolinkurl{10.1007/s00220-011-1391-x}.

\bibitem{Serrin1962}
J.~Serrin,
\emph{On the interior regularity of weak solutions of the Navier--Stokes
 equations},
Arch. Rational Mech. Anal. \textbf{9} (1962), 187--195,
doi:\ \nolinkurl{10.1007/BF00253344}.

\bibitem{Stein1970}
E.~M. Stein,
\emph{Singular Integrals and Differentiability Properties of Functions},
Princeton Mathematical Series, No.~30, Princeton University Press, 1970.

\bibitem{Aubin1963}
J.-P. Aubin,
\emph{Un theoreme de compacite},
C. R. Acad. Sci. Paris \textbf{256} (1963), 5042--5044.

\bibitem{Kato1984}
T. Kato,
\emph{Strong \(L^p\)-solutions of the Navier--Stokes equation in
\(\mathbb R^m\), with applications to weak solutions},
Math. Z. \textbf{187} (1984), no. 4, 471--480,
doi:\url{https://doi.org/10.1007/BF01174182}.

\bibitem{KNSS2009}
G. Koch, N. Nadirashvili, G. Seregin, and V. \v{S}verak,
\emph{Liouville theorems for the Navier--Stokes equations and applications},
Acta Math. \textbf{203} (2009), no. 1, 83--105.

\bibitem{LionsMagenes}
J.-L. Lions and E. Magenes,
\emph{Non-homogeneous boundary value problems and applications. Vol. I},
Springer-Verlag, 1972.

\bibitem{NRS1996}
J. Necas, M. Ruzicka, and V. \v{S}verak,
\emph{On Leray's self-similar solutions of the Navier--Stokes equations},
Acta Math. \textbf{176} (1996), no. 2, 283--294.

\bibitem{Simon1987}
J. Simon,
\emph{Compact sets in the space \(L^p(0,T;B)\)},
Ann. Mat. Pura Appl. (4) \textbf{146} (1987), 65--96,
doi:\url{https://doi.org/10.1007/BF01762360}.

\bibitem{Sohr2001}
H. Sohr,
\emph{The Navier--Stokes equations. An elementary functional analytic approach},
Birkhauser, 2001.

\bibitem{Billingsley1999}
P. Billingsley,
\emph{Convergence of Probability Measures}, second ed.,
Wiley, New York, 1999.

\bibitem{FabesJonesRiviere1972}
E. B. Fabes, B. F. Jones, and N. M. Riviere,
\emph{The initial value problem for the Navier--Stokes equations with data in
\(L^p\)},
Arch. Rational Mech. Anal. 45 (1972), 222--240.

\bibitem{GigaMiyakawa1985}
Y. Giga and T. Miyakawa,
\emph{Solutions in \(L_r\) of the Navier--Stokes initial value problem},
Arch. Rational Mech. Anal. 89 (1985), 267--281.

\bibitem{KryloffBogoliouboff1937}
N. Kryloff and N. Bogoliouboff,
\emph{La theorie generale de la mesure dans son application a l'etude des
systemes dynamiques de la mecanique non lineaire},
Ann. of Math. (2) 38 (1937), 65--113.

\bibitem{Ladyzhenskaya1968}
O. A. Ladyzhenskaya, V. A. Solonnikov, and N. N. Ural'tseva,
\emph{Linear and Quasilinear Equations of Parabolic Type},
American Mathematical Society, Providence, 1968.

\bibitem{GallayWayne2002}
T. Gallay and C. E. Wayne,
\emph{Long-time asymptotics of the Navier--Stokes and vorticity equations on
\(\mathbb R^3\)},
Philos. Trans. Roy. Soc. London Ser. A \textbf{360} (2002), no. 1799,
2155--2188, doi:10.1098/rsta.2002.1068.

\bibitem{LRZ}
Z. Lei, X. Ren, and Q. S. Zhang,
\emph{A Liouville theorem for axially symmetric Navier--Stokes equations on
\(\mathbb R^2\times\mathbb T^1\)},
Math. Ann. \textbf{383} (2022), no. 1--2, 415--431,
doi:10.1007/s00208-020-02128-9.

\bibitem{LZZ}
Z. Lei, Q. Zhang, and N. Zhao,
\emph{Improved Liouville theorems for axially symmetric Navier--Stokes
equations},
Sci. Sin. Math. \textbf{47} (2017), no. 10, 1183--1198,
doi:10.1360/N012016-00149.

\end{thebibliography}

\end{document}
